% Leçon 6 — Grammaire : ……才/ 就...... Et 还 ; 又 ; 再 — Chinois Tle A — Cours signé OnBuch+
% Programme officiel MINESEC. Compiler avec : tectonic ZH06-grammaire-et.tex
\newcommand{\DOCMATIERE}{Chinois}
\newcommand{\DOCNIVEAU}{Tle A}
\newcommand{\DOCMODULE}{Module 1 : Vie familiale et intégration sociale}
\newcommand{\DOCLECON}{ZH06}
\newcommand{\DOCTITRE}{Leçon 6 — Grammaire : ……才/ 就...... Et 还 ; 又 ; 再}
% OnBuch+ — préambule « cours signés » ENRICHI (compilé avec tectonic / XeLaTeX)
% Système visuel « Pop » d'OnBuch+ : fond crème (jamais blanc), bordures encre
% épaisses, ombres « dures » décalées, titres Archivo Black en capitales.
% Version MATHÉMATIQUES : graphes (pgfplots/tikz), boîtes pédagogiques
% (définition, propriété, méthode, attention, le savais-tu, exercice résolu,
% exercice, TP, bilan), page de garde et signature OnBuch+ ancrée en bas.
%
% Macros attendues dans main.tex AVANT \input{preamble} :
%   \DOCMATIERE  \DOCNIVEAU  \DOCMODULE  \DOCLECON (numéro)  \DOCTITRE
\documentclass[11pt]{article}

\usepackage{fontspec}
\usepackage[french]{babel} % français = langue principale ; le chinois via \zh{..} / textebox (xeCJK)
\usepackage{xeCJK}
\usepackage{pifont} % \ding{51} \ding{55} utilisés par les modèles
\usepackage[a4paper,margin=2cm,top=3.4cm,bottom=2.8cm,headheight=1.4cm,footskip=1.3cm]{geometry}
\usepackage{amsmath,amssymb,amsfonts}
\usepackage{mathtools} % \xmapsto, \coloneqq... souvent utilisés par les modèles
\usepackage{cancel} % simplifications barrées (bilans de forces)
% \abs{z} (module, valeur absolue) et \norm{u} (norme), utilisés par les modèles
\providecommand{\abs}{}\let\abs\relax\DeclarePairedDelimiter{\abs}{\lvert}{\rvert}
\providecommand{\norm}{}\let\norm\relax\DeclarePairedDelimiter{\norm}{\lVert}{\rVert}
\usepackage[table]{xcolor} % [table] : fournit aussi \rowcolors (alternance de lignes)
\usepackage{pagecolor}
\usepackage{tikz}
\usetikzlibrary{arrows.meta,positioning,calc,shapes.geometric,shapes.misc,shapes.arrows,decorations.pathmorphing,decorations.pathreplacing,decorations.markings,patterns,shadows,mindmap,trees,fit,backgrounds,matrix,intersections,angles,quotes}
\usepackage{pgfplots}
\usepackage[version=4]{mhchem} % équations nucléaires \ce{^{235}_{92}U}
\usepackage[european,siunitx]{circuitikz} % schémas électriques
\pgfplotsset{compat=1.18}
\usepgfplotslibrary{fillbetween,groupplots} % aires entre courbes (calcul intégral)
\usepackage{enumitem}
\usepackage{fancyhdr}
% [most] charge toutes les bibliothèques tcolorbox (listings, minted, raster,
% documentation...) et gonfle inutilement la mémoire TeX sur un document long
% (~300 boîtes) : on ne charge que ce qui est réellement utilisé.
\usepackage[skins,breakable]{tcolorbox}
\usepackage{graphicx}
\usepackage{colortbl}
\usepackage{booktabs}
\usepackage{tabularx}
\usepackage{diagbox} % case d'en-tête diagonale des tableaux à double entrée
% Colonnes extensibles centrée / gauche / droite (C, L, R), souvent utilisées par les modèles
\newcolumntype{C}{>{\centering\arraybackslash}X}
\newcolumntype{L}{>{\raggedright\arraybackslash}X}
\newcolumntype{R}{>{\raggedleft\arraybackslash}X}
\usepackage{array}
\usepackage{multicol}
\usepackage{multirow}
\usepackage{siunitx}
% parse-numbers=false : accepte aussi \SI{2,5 \times 10^{-3}}{...}
\sisetup{locale=FR,output-decimal-marker={,},group-minimum-digits=4,parse-numbers=false}
\DeclareSIUnit\ohm{\ensuremath{\symup{\Omega}}} % U+2126 absent de la police
\DeclareSIUnit\division{div} % réglages d'oscilloscope (ms/div, V/div)
\DeclareSIUnit\tr{tr} % tours (vitesse de rotation, ex. \SI{1500}{\tr\per\minute})
\DeclareSIUnit\molar{M} % concentration molaire (ex. \SI{3}{\molar}, \SI{1}{\micro\molar})
\DeclareSIUnit\an{an} % année (le modèle confond parfois avec \year, un compteur TeX) — ex. \SI{5}{\kilo\meter\per\an}
\DeclareSIUnit\FCFA{FCFA} % franc CFA (ex. \SI{500}{\FCFA\per\kilo\gram})
\DeclareSIUnit\degreeBrix{\textdegree Brix} % teneur en sucre d'un sirop/jus (réfractométrie)
\DeclareSIUnit\month{mois} % le modèle utilise parfois \month, un compteur TeX, comme unité de durée
\usepackage{qrcode}
% \degree : commande fréquemment utilisée par les modèles pour un angle ou une
% température en dehors de \SI{}{} (non fournie par les paquets chargés).
\providecommand{\degree}{^{\circ}}
% \qed (fin de démonstration), utilisé par les modèles sans amsthm
\providecommand{\qed}{\hfill$\square$}
% \barre : double barre des valeurs interdites dans un tableau de variations (tkz-tab)
\providecommand{\barre}{\ensuremath{\|}}
% environnement proof (amsthm non chargé) utilisé par les modèles
\ifdefined\proof\else\newenvironment{proof}[1][Démonstration]{\par\noindent\textit{#1.}\ }{\qed\par}\fi
\usepackage{lastpage}
\usepackage{hyperref}
\hypersetup{hidelinks}

% ============================================================
% Palette « Pop » OnBuch+ — mêmes hex que la classe P (plus_ui.dart)
% ============================================================
\definecolor{poporange}{HTML}{F59321}
\definecolor{popdark}{HTML}{E07A0C}
\definecolor{popcream}{HTML}{FFF6E8}
\definecolor{popink}{HTML}{1C1714}
\definecolor{poppurple}{HTML}{7A5AE0}
\definecolor{popgreen}{HTML}{2E9E6A}
\definecolor{poppink}{HTML}{E0457B}
\definecolor{popblue}{HTML}{2F7DE1}
\definecolor{popgold}{HTML}{C98A12}
\definecolor{popmuted}{HTML}{8A7F76}
\definecolor{popline}{HTML}{EADFD2}
\definecolor{popgray}{HTML}{8A7F76}
\colorlet{popgrey}{popgray}
% Alias tolérants (le modèle nomme parfois les couleurs différemment)
\colorlet{popwhite}{white}
\colorlet{popblack}{popink}
\colorlet{popred}{poppink}
\colorlet{popyellow}{popgold}
% Teintes claires (fonds de tableaux/encadrés) — le modèle nomme parfois
% les couleurs avec un suffixe L (ex. popblueL) sans les avoir définies.
\colorlet{poporangeL}{poporange!20}
\colorlet{popdarkL}{popdark!20}
\colorlet{poppurpleL}{poppurple!20}
\colorlet{popgreenL}{popgreen!20}
\colorlet{poppinkL}{poppink!20}
\colorlet{popblueL}{popblue!20}
\colorlet{popgoldL}{popgold!20}
\colorlet{popredL}{popred!20}
\colorlet{popgrayL}{popgray!20}
% Teintes claires pour fonds de tableaux / zones
\colorlet{popblueL}{popblue!10!white}
\colorlet{poporangeL}{poporange!14!white}
\colorlet{popgreenL}{popgreen!12!white}
\colorlet{poppurpleL}{poppurple!12!white}
\colorlet{poppinkL}{poppink!12!white}
\colorlet{popgoldL}{popgold!14!white}

\pagecolor{popcream}

% ============================================================
% Typographie
% ============================================================
\defaultfontfeatures{Ligatures=TeX}
\setmainfont{PlusJakartaSans-Medium.ttf}[Path=../../../fonts/,BoldFont=PlusJakartaSans-Bold.ttf]
% Symboles, API et emojis absents de la police principale : repli sur DejaVu Sans (caractères actifs)
\newfontface\symfont{DejaVuSans.ttf}[Path=../../../fonts/]
\makeatletter
\newcommand\symactive[1]{\catcode#1=13 \begingroup\lccode`\~=#1 \lowercase{\endgroup\def~}{{\symfont\char#1}}}
\newcommand\symblank[1]{\catcode#1=13 \begingroup\lccode`\~=#1 \lowercase{\endgroup\def~}{}}
\newcommand\symhyph[1]{\catcode#1=13 \begingroup\lccode`\~=#1 \lowercase{\endgroup\def~}{\mbox{-}}}
\newcount\symc
\newcommand\symrange[2]{\symc=#1 \@whilenum\symc<#2 \do{\expandafter\symactive\expandafter{\the\symc}\advance\symc by 1 }}
\newcommand\symblankrange[2]{\symc=#1 \@whilenum\symc<#2 \do{\expandafter\symblank\expandafter{\the\symc}\advance\symc by 1 }}
\makeatother
\symrange{"0250}{"02FF}\symactive{"2713}\symactive{"2714}\symactive{"2717}\symactive{"2718}\symactive{"2610}\symactive{"2611}\symactive{"2612}
\symrange{"2600}{"27BF}
\catcode"2705=13 \begingroup\lccode`\~="2705 \lowercase{\endgroup\def~}{{\symfont\char"2713}}\symblank{"26BD}\symblank{"26A1}
\symrange{"2460}{"257F}\symactive{"223C}\symactive{"2248}\symactive{"2260}
\symhyph{"2011}\symhyph{"2E1A}\symblankrange{"1F300}{"1FAFF}\symblank{"FE0F}
% Caractères chinois : WenQuanYi Zen Hei (sous-ensemble GB2312 + pinyin), gras/italique simulés
\setCJKmainfont{WenQuanYiZenHei-subset.ttf}[Path=../../../fonts/,AutoFakeBold=3,AutoFakeSlant=0.2]
\xeCJKsetup{CJKspace=false}
% Pinyin : voyelles à caron (ǎ ǐ ǒ ǔ ǖ ǘ ǚ ǜ) absentes de la police principale -> caractères actifs
\newfontface\pyfont{WenQuanYiZenHei-subset.ttf}[Path=../../../fonts/]
\newcommand\pyactive[1]{\catcode#1=13 \begingroup\lccode`\~=#1 \lowercase{\endgroup\def~}{{\pyfont\char#1}}}
\pyactive{"01CD}\pyactive{"01CE}\pyactive{"01CF}\pyactive{"01D0}\pyactive{"01D1}\pyactive{"01D2}\pyactive{"01D3}\pyactive{"01D4}\pyactive{"01D5}\pyactive{"01D6}\pyactive{"01D7}\pyactive{"01D8}\pyactive{"01D9}\pyactive{"01DA}\pyactive{"01DB}\pyactive{"01DC}
% Maths en sans-serif (Fira Math) avec chiffres et lettres droites de Plus
% Jakarta, pour que les formules se fondent dans le texte.
\usepackage[math-style=ISO,bold-style=ISO]{unicode-math}
\setmathfont{FiraMath-Regular.otf}[Path=../../../fonts/]
\setmathfont{PlusJakartaSans-Medium.ttf}[Path=../../../fonts/,range={up/num,up/{Latin,latin},\mathup}]
\usepackage{icomma} % virgule décimale sans espace en mode math
% Caractères mathématiques Unicode absents des polices de texte (Plus Jakarta,
% Archivo Black) : rendus par leur équivalent mathématique, en texte comme en
% formule — sinon ils s'affichent en carré vide (ex. « Arithmétique dans ℤ »).
\usepackage{newunicodechar}
\newunicodechar{ℝ}{\ensuremath{\mathbb{R}}}
\newunicodechar{ℤ}{\ensuremath{\mathbb{Z}}}
\newunicodechar{ℂ}{\ensuremath{\mathbb{C}}}
\newunicodechar{ℕ}{\ensuremath{\mathbb{N}}}
\newunicodechar{ℚ}{\ensuremath{\mathbb{Q}}}
\newunicodechar{𝔻}{\ensuremath{\mathbb{D}}}
\newunicodechar{α}{\ensuremath{\alpha}}
\newunicodechar{β}{\ensuremath{\beta}}
\newunicodechar{γ}{\ensuremath{\gamma}}
\newunicodechar{δ}{\ensuremath{\delta}}
\newunicodechar{ε}{\ensuremath{\varepsilon}}
\newunicodechar{θ}{\ensuremath{\theta}}
\newunicodechar{λ}{\ensuremath{\lambda}}
\newunicodechar{μ}{\ensuremath{\mu}}
\newunicodechar{σ}{\ensuremath{\sigma}}
\newunicodechar{τ}{\ensuremath{\tau}}
\newunicodechar{φ}{\ensuremath{\varphi}}
\newunicodechar{ψ}{\ensuremath{\psi}}
\newunicodechar{ω}{\ensuremath{\omega}}
\newunicodechar{Σ}{\ensuremath{\Sigma}}
% majuscules produites par \MakeUppercase dans les titres (α-aminés -> Α)
\newunicodechar{Α}{\ensuremath{\alpha}}
\newunicodechar{Β}{\ensuremath{\beta}}
\newunicodechar{Γ}{\ensuremath{\Gamma}}
\newunicodechar{Δ}{\ensuremath{\Delta}}
\newunicodechar{Ε}{\ensuremath{\varepsilon}}
\newunicodechar{Θ}{\ensuremath{\Theta}}
\newunicodechar{Λ}{\ensuremath{\Lambda}}
\newunicodechar{Μ}{\ensuremath{\mu}}
\newunicodechar{Τ}{\ensuremath{\tau}}
\newunicodechar{Φ}{\ensuremath{\Phi}}
\newunicodechar{Ψ}{\ensuremath{\Psi}}
\newunicodechar{Ω}{\ensuremath{\Omega}}
\newunicodechar{↦}{\ensuremath{\mapsto}}
\newunicodechar{∈}{\ensuremath{\in}}
\newunicodechar{∉}{\ensuremath{\notin}}
\newunicodechar{∧}{\ensuremath{\wedge}}
\newunicodechar{⇔}{\ensuremath{\Leftrightarrow}}
\newunicodechar{⟺}{\ensuremath{\Longleftrightarrow}}
\newunicodechar{⇒}{\ensuremath{\Rightarrow}}
\newunicodechar{⟹}{\ensuremath{\Longrightarrow}}
\newunicodechar{∘}{\ensuremath{\circ}}
\newunicodechar{∪}{\ensuremath{\cup}}
\newunicodechar{∩}{\ensuremath{\cap}}
\newunicodechar{≡}{\ensuremath{\equiv}}
\newunicodechar{⊂}{\ensuremath{\subset}}
\newunicodechar{⊥}{\ensuremath{\perp}}
\newunicodechar{∃}{\ensuremath{\exists}}
\newunicodechar{∀}{\ensuremath{\forall}}
\newunicodechar{∛}{\ensuremath{\sqrt[3]{\,}}}
\newunicodechar{∜}{\ensuremath{\sqrt[4]{\,}}}
\newunicodechar{⃗}{\ensuremath{{}^{\to}}}
\newunicodechar{ⁿ}{\ifmmode^{n}\else\textsuperscript{\ensuremath{n}}\fi}
\newunicodechar{ˣ}{\ifmmode^{x}\else\textsuperscript{\ensuremath{x}}\fi}
\newunicodechar{ᵇ}{\ifmmode^{b}\else\textsuperscript{\ensuremath{b}}\fi}
\newunicodechar{ʳ}{\ifmmode^{r}\else\textsuperscript{\ensuremath{r}}\fi}
\newunicodechar{ᵘ}{\ifmmode^{u}\else\textsuperscript{\ensuremath{u}}\fi}
\newunicodechar{ʸ}{\ifmmode^{y}\else\textsuperscript{\ensuremath{y}}\fi}
\newunicodechar{ᵃ}{\ifmmode^{a}\else\textsuperscript{\ensuremath{a}}\fi}
\newunicodechar{ᵉ}{\ifmmode^{e}\else\textsuperscript{\ensuremath{e}}\fi}
\newunicodechar{ᵏ}{\ifmmode^{k}\else\textsuperscript{\ensuremath{k}}\fi}
\newunicodechar{ᵖ}{\ifmmode^{p}\else\textsuperscript{\ensuremath{p}}\fi}
\newunicodechar{ᴺ}{\ifmmode^{N}\else\textsuperscript{\ensuremath{N}}\fi}
\newunicodechar{ᶜ}{\ifmmode^{c}\else\textsuperscript{\ensuremath{c}}\fi}
\newunicodechar{ʰ}{\ifmmode^{h}\else\textsuperscript{\ensuremath{h}}\fi}
\newunicodechar{ᵠ}{\ifmmode^{q}\else\textsuperscript{\ensuremath{q}}\fi}
\newunicodechar{ₙ}{\ifmmode_{n}\else\textsubscript{\ensuremath{n}}\fi}
\newunicodechar{ₐ}{\ifmmode_{a}\else\textsubscript{\ensuremath{a}}\fi}
\newunicodechar{ᵢ}{\ifmmode_{i}\else\textsubscript{\ensuremath{i}}\fi}
\newunicodechar{ᵣ}{\ifmmode_{r}\else\textsubscript{\ensuremath{r}}\fi}
\newunicodechar{ₖ}{\ifmmode_{k}\else\textsubscript{\ensuremath{k}}\fi}
\newunicodechar{ᵦ}{\ifmmode_{\beta}\else\textsubscript{\ensuremath{\beta}}\fi}
\newunicodechar{ₓ}{\ifmmode_{x}\else\textsubscript{\ensuremath{x}}\fi}
\newfontfamily\popdisplay{ArchivoBlack-Regular.ttf}[Path=../../../fonts/]
\newfontfamily\poplogo{SpaceGrotesk-Bold.ttf}[Path=../../../fonts/]
\newfontfamily\popsemi{PlusJakartaSans-SemiBold.ttf}[Path=../../../fonts/]
\newfontfamily\popxbold{PlusJakartaSans-ExtraBold.ttf}[Path=../../../fonts/]
\newfontfamily\popxboldls{PlusJakartaSans-ExtraBold.ttf}[Path=../../../fonts/,LetterSpace=3.0]

\setlist[enumerate]{leftmargin=1.7em,itemsep=5pt,topsep=4pt}
\setlist[itemize]{leftmargin=1.7em,itemsep=4pt,topsep=4pt,label={\color{poporange}\textbullet}}
\setlength{\parindent}{0pt}
\setlength{\parskip}{5pt}
\renewcommand{\arraystretch}{1.25}

% pgfplots : style Pop par défaut
\pgfplotsset{
  every axis/.append style={axis line style={popink,line width=1pt},tick style={popink},
    grid style={popline,line width=0.5pt},label style={font=\small},tick label style={font=\footnotesize}},
  popcurve/.style={poporange,line width=1.8pt,no markers,smooth},
}
% Même style disponible dans un \draw TikZ simple (hors pgfplots)
\tikzset{popcurve/.style={poporange,line width=1.8pt}}
\tikzset{popgrid/.style={popline,very thin}}
\colorlet{popcurve}{poporange} % le nom du style est parfois utilisé comme couleur

% ============================================================
% Pill / squiggle
% ============================================================
\newcommand{\poppill}[3]{%
  \tikz[baseline=-0.55ex]\node[draw=popink,line width=1.2pt,rounded corners=2.6pt,%
    fill=#2,inner xsep=6.5pt,inner ysep=3pt]{%
    \popxboldls\fontsize{7}{7}\selectfont\color{#3}\MakeUppercase{#1}};%
}
\newcommand{\popsquiggle}[1]{%
  \tikz[baseline=-0.4ex]{
    \draw[#1,line width=2.6pt,line cap=round]
      plot [smooth] coordinates {(0,0.09) (0.34,0.01) (0.68,0.21) (1.02,0.01) (1.36,0.10)};
  }%
}
% Mot-clé surligné (marqueur orange sous le texte)
\newcommand{\cle}[1]{{\popxbold\color{popdark}#1}}

% ============================================================
% En-tête / pied
% ============================================================
\pagestyle{fancy}
\fancyhf{}
\renewcommand{\headrulewidth}{0pt}
\renewcommand{\footrulewidth}{0pt}
\lhead{\raisebox{-0.15ex}{\poplogo\fontsize{13}{13}\selectfont\color{popink}OnBuch\color{poporange}+}}
\rhead{\poppill{\DOCMATIERE\ \textbullet\ \DOCNIVEAU\ \textbullet\ Leçon \DOCLECON}{white}{popink}}
\lfoot{\popxbold\fontsize{7.6}{7.6}\selectfont\color{popmuted}\MakeUppercase{Cours signé OnBuch+}\ \textbullet\ {\popsemi\DOCTITRE}}
\rfoot{\tikz[baseline=-0.6ex]\node[rounded corners=8.5pt,fill=popink,minimum height=17pt,minimum width=17pt,inner xsep=5pt,inner sep=0]{\popxbold\fontsize{8}{8}\selectfont\color{popcream}\thepage\,/\,\pageref*{LastPage}};}

% ============================================================
% Titres
% ============================================================
\newcommand{\chaptitre}[2]{%
  \begin{tcolorbox}[enhanced,colback=poporange,colframe=popink,boxrule=2.6pt,arc=11pt,
      drop shadow=popink,left=15pt,right=15pt,top=12pt,bottom=11pt,before skip=3pt,after skip=18pt]
    \poppill{#2}{popink}{popcream}\par\vspace{9pt}
    {\raggedright\hyphenpenalty=10000\exhyphenpenalty=10000\popdisplay\fontsize{22}{24}\selectfont\color{popcream}\MakeUppercase{#1}\par}
  \end{tcolorbox}
}
% Section numérotée : pastille orange avec le numéro + titre Archivo
\newcounter{popsec}
\newcommand{\coursec}[1]{%
  \stepcounter{popsec}\setcounter{popsub}{0}%
  \par\vspace{16pt}\noindent
  \begin{minipage}{\linewidth}
  \tikz[baseline=-0.7ex]\node[draw=popink,line width=1.6pt,fill=poporange,rounded corners=4pt,minimum size=22pt,inner sep=0]{\popdisplay\fontsize{12}{12}\selectfont\color{popcream}\thepopsec};%
  \hspace{8pt}{\popdisplay\fontsize{15}{17}\selectfont\color{popink}\MakeUppercase{#1}}\par
  \vspace{4pt}\noindent{\color{poporange}\rule{36pt}{3.6pt}}
  \end{minipage}\par\nopagebreak\vspace{8pt}
}
\newcounter{popsub}
\newcommand{\courssub}[1]{%
  \stepcounter{popsub}%
  \par\vspace{9pt}\noindent{\popxbold\fontsize{11.5}{13}\selectfont\color{popink}{\color{poporange}\thepopsec.\thepopsub}\hspace{6pt}#1}\par\nopagebreak\vspace{4pt}
}

% ============================================================
% Boîtes pédagogiques — même recette : fond blanc, bordure épaisse colorée,
% ombre dure encre, badge plein. Titre optionnel [#1].
% ============================================================
\tcbset{popbase/.style={breakable,enhanced,colback=white,boxrule=2.3pt,arc=9pt,
  shadow={3pt}{-3pt}{0pt}{popink},left=14pt,right=14pt,top=11pt,bottom=11pt,before skip=11pt,after skip=15pt,
  fontupper=\popsemi\fontsize{10.6}{14.5}\selectfont\color{popink},
  fontlower=\popsemi\fontsize{10.6}{14.5}\selectfont\color{popink}}}
% Coche dessinée (le glyphe ✓ n'existe pas dans les polices du document)
\newcommand{\popcheck}[1]{\tikz[baseline=-0.5ex]\draw[#1,line width=1.6pt,line cap=round,line join=round](-0.13,0.0)--(-0.04,-0.09)--(0.13,0.1);}
\providecommand{\coche}{\popcheck{popgreen}}
\providecommand{\resultat}[1]{\par\smallskip\noindent\fcolorbox{popgreen}{popgreenL}{\parbox{\dimexpr\linewidth-2\fboxsep-2\fboxrule\relax}{\textbf{Résultat.} #1}}\par\smallskip}
\newcommand{\popboxhead}[3]{\poppill{#1}{#2}{white}\ifx\relax#3\relax\else\hspace{6pt}{\popxbold\fontsize{10.6}{12}\selectfont\color{popink}#3}\fi\par\vspace{7pt}}

\newtcolorbox{aretenir}[1][]{popbase,colframe=popgreen,before upper={\popboxhead{À retenir}{popgreen}{#1}}}
% Texte en chinois (document de lecture, dialogue, extrait) : cadre
% d'étude. Un titre optionnel (source, auteur) s'affiche dans l'en-tête.
\newtcolorbox{textebox}[1][]{popbase,colframe=popblue,before upper={\popboxhead{文本 · Texte}{popblue}{#1}},after upper={}}
% Marques ✓ / ✗ (forme correcte / fautive) souvent utilisées par les modèles
\providecommand{\cross}{\textcolor{poppink}{\ensuremath{\times}}}
\providecommand{\xmark}{\cross}\providecommand{\crossmark}{\cross}\providecommand{\cmark}{\ensuremath{\checkmark}}
% Ligne de réponse / justification à compléter (inventée par les modèles)
\providecommand{\justification}[1]{\par\noindent\textit{Justification~:} \dotfill\par}
% \textipa (package tipa non chargé) : la police ne couvre pas l'API -> simple passage
\providecommand{\textipa}[1]{#1}
% Soulignement (ulem non chargé) : \uline, \ul
\providecommand{\uline}[1]{\underline{#1}}\providecommand{\ul}[1]{\underline{#1}}
% \glossaire{\item ...} inventée par les modèles : liste de glossaire
\providecommand{\glossaire}[1]{\begin{itemize}#1\end{itemize}}
% \alert (beamer) inventée par les modèles : texte mis en évidence
\providecommand{\alert}[1]{\textbf{\textcolor{poppink}{#1}}}
% variantes ASCII d'ordinaux écrites en macro
\providecommand{\iere}{\textsuperscript{re}}\providecommand{\ieme}{\textsuperscript{e}}\providecommand{\ere}{\textsuperscript{re}}\providecommand{\eme}{\textsuperscript{e}}\providecommand{\er}{\textsuperscript{er}}\providecommand{\ie}{\textsuperscript{e}}
% Ordinaux français écrits en macro par les modèles : 1\ière, 3\ième
\newcommand{\ière}{\textsuperscript{re}}\newcommand{\ième}{\textsuperscript{e}}
% \exemple (inventée par les modèles) : étiquette « Exemple : »
\providecommand{\exemple}{\textbf{Exemple~:}\ }
% \square utilisable aussi hors mode math (cases à cocher)
\AtBeginDocument{\let\squareMath\square\renewcommand{\square}{\ensuremath{\squareMath}}}
% Mot ou phrase en chinois à l'intérieur d'un paragraphe français
\newcommand{\zh}[1]{\ifmmode\text{#1}\else{#1}\fi}
\let\eng\zh \let\esp\zh \let\all\zh % alias tolérés
% flèches d'intonation inventées par les modèles
\providecommand{\textsearrow}{}\renewcommand{\textsearrow}{\ensuremath{\searrow}}\providecommand{\textnearrow}{}\renewcommand{\textnearrow}{\ensuremath{\nearrow}}
% \fr{..} : mot français (inventé par les modèles)
\providecommand{\fr}[1]{\textit{#1}}
% zhbox : encadré de texte chinois inventé par les modèles
\newenvironment{zhbox}{\begin{quote}}{\end{quote}}
% \vs : « versus » inventé par les modèles
\providecommand{\vs}{\textit{vs}\ }
% \blank : case à compléter inventée par les modèles
\providecommand{\blank}[1][]{\underline{\hspace{1.6em}}}
\newtcolorbox{exemplebox}[1][]{popbase,colframe=poporange,before upper={\popboxhead{Exemple}{poporange}{#1}}}
\newtcolorbox{definition}[1][]{popbase,colframe=popblue,before upper={\popboxhead{Définition}{popblue}{#1}}}
\newtcolorbox{propriete}[1][]{popbase,colframe=poppurple,before upper={\popboxhead{Propriété}{poppurple}{#1}}}
\newtcolorbox{methode}[1][]{popbase,colframe=popgold,before upper={\popboxhead{Méthode}{popgold}{#1}}}
\newtcolorbox{attention}[1][]{popbase,colframe=poppink,before upper={\popboxhead{Attention}{poppink}{#1}}}
\newtcolorbox{experience}[1][]{popbase,colframe=popblue,colback=popblueL,before upper={\popboxhead{Expérience}{popblue}{#1}}}
% « Le savais-tu ? » : carte violette pleine (contexte, culture, Cameroun)
\newtcolorbox{savaistu}[1][]{popbase,colback=poppurple,colframe=popink,
  fontupper=\popsemi\fontsize{10.4}{14.2}\selectfont\color{white},
  before upper={\poppill{Le savais-tu\,?}{white}{poppurple}\ifx\relax#1\relax\else\hspace{6pt}{\popxbold\color{white}#1}\fi\par\vspace{7pt}}}
% Exercice résolu : énoncé en haut, \tcblower puis la solution sur fond crème
\newcounter{popexo}\newcounter{popexr}
\newtcolorbox{exoresolu}[1][]{popbase,colframe=poporange,colbacklower=popcream,
  segmentation style={popink,line width=1.2pt,dashed},
  before upper={\stepcounter{popexr}\popboxhead{Exercice résolu \thepopexr}{poporange}{#1}},
  before lower={\poppill{Solution}{popink}{popcream}\par\vspace{6pt}}}
% Exercice (non résolu en place) — la correction est dans la partie Corrigés
\newtcolorbox{exercice}[1][]{popbase,colframe=popink,
  before upper={\stepcounter{popexo}\popboxhead{Exercice \thepopexo}{popink}{#1}}}
% Corrigé
\newtcolorbox{corrige}[1][]{popbase,colframe=popgreen,colback=popgreenL,
  before upper={\popboxhead{Corrigé}{popgreen}{#1}}}
% Objectifs / prérequis / bilan
\newtcolorbox{objectifs}{popbase,colframe=popink,colback=poporangeL,
  before upper={\popboxhead{Objectifs de la leçon}{popink}{}}}
\newtcolorbox{prerequis}{popbase,colframe=popink,colback=popblueL,
  before upper={\popboxhead{Prérequis}{popblue}{}}}
\newtcolorbox{bilan}{popbase,colframe=popink,colback=white,boxrule=2.8pt,
  before upper={\popboxhead{Fiche bilan}{poporange}{L'essentiel à connaître pour l'examen}}}
% Figure encadrée avec légende
\newtcolorbox{popfigure}[1][]{enhanced,breakable=false,colback=white,colframe=popline,boxrule=1.4pt,arc=8pt,
  left=10pt,right=10pt,top=10pt,bottom=8pt,before skip=10pt,after skip=12pt,halign=center,
  lower separated=false,halign lower=center,
  fontlower=\popsemi\fontsize{9}{11.5}\selectfont\color{popmuted},
  #1}
\newcommand{\legende}[1]{\par\vspace{6pt}{\popsemi\fontsize{9}{11.5}\selectfont\color{popmuted}#1}}

% ============================================================
% Signature OnBuch+
% ============================================================
\newcommand{\poplogoinline}[1]{{\poplogo\fontsize{#1}{#1}\selectfont\color{popink}OnBuch\color{poporange}+}}
\newcommand{\signatureonbuch}{%
  \begin{tcolorbox}[enhanced,colback=popink,colframe=popink,boxrule=2.2pt,arc=10pt,
      drop shadow=poporange,left=14pt,right=14pt,top=10pt,bottom=10pt,before skip=0pt,after skip=0pt]
    \tikz[baseline=-0.7ex]\node[circle,fill=popgreen,draw=popcream,line width=1.3pt,minimum size=22pt,inner sep=0]{\popcheck{white}};%
    \hspace{9pt}\begin{minipage}[c]{0.62\linewidth}
      {\popxbold\fontsize{12}{13}\selectfont\color{popcream}Signé\ \textbullet\ {\poplogo OnBuch\color{poporange}+}}\par\vspace{2pt}
      {\raggedright\popsemi\fontsize{8}{10.5}\selectfont\color{popline}\MakeUppercase{Équipe pédagogique OnBuch+}\\\MakeUppercase{Conforme au programme officiel MINESEC}\par}
    \end{minipage}\hfill
    {\poplogo\fontsize{22}{22}\selectfont\color{popcream}OnBuch\color{poporange}+}
  \end{tcolorbox}%
}

% ============================================================
% Page de garde
% #1 titre · #2 matière · #3 niveau · #4 module · #5 n° leçon · #6 durée ·
% #7 description / objectifs (texte ou itemize)
% ============================================================
\newcommand{\pagegarde}[7]{%
  \thispagestyle{empty}
  \newgeometry{a4paper,margin=1.7cm,top=1.6cm,bottom=1.4cm}%
  \begin{tikzpicture}[remember picture,overlay]
    \fill[poporange!18] ([xshift=-3.2cm,yshift=-3.6cm]current page.north east) circle (4.6cm);
    \fill[poppurple!14] ([xshift=2.4cm,yshift=7.6cm]current page.south west) circle (3.8cm);
  \end{tikzpicture}%
  {\poplogo\fontsize{30}{30}\selectfont\color{popink}OnBuch\color{poporange}+}\hfill
  \raisebox{0.6ex}{\poppill{Cours signé \textbullet\ Premium}{popink}{popcream}}\par
  \vspace{1pt}\popsquiggle{poppurple}\par
  \vspace{8pt}
  \begin{tcolorbox}[enhanced,colback=poporange,colframe=popink,boxrule=2.8pt,arc=13pt,
      drop shadow=popink,left=18pt,right=18pt,top=12pt,bottom=13pt,after skip=10pt]
    \poppill{#2 \textbullet\ #3}{popink}{popcream}\hspace{5pt}\poppill{#4}{white}{popink}\par\vspace{8pt}
    {\popdisplay\fontsize{14}{14}\selectfont\color{popink}LEÇON #5}\par\vspace{4pt}
    {\raggedright\hyphenpenalty=10000\exhyphenpenalty=10000\popdisplay\fontsize{25}{27}\selectfont\color{popcream}\MakeUppercase{#1}\par}
    \vspace{8pt}
    {\popxbold\fontsize{9.5}{9.5}\selectfont\color{popink}\MakeUppercase{Durée conseillée : #6}}
  \end{tcolorbox}
  \begin{tcolorbox}[enhanced,colback=white,colframe=popink,boxrule=2.2pt,arc=10pt,
      drop shadow=popink,left=16pt,right=16pt,top=10pt,bottom=10pt,after skip=8pt]
    \poppill{Dans cette leçon}{poppurple}{white}\par\vspace{5pt}
    \popsemi\fontsize{9.3}{12.6}\selectfont\color{popink}#7
  \end{tcolorbox}
  \noindent\begin{minipage}[t]{0.487\linewidth}
    \begin{tcolorbox}[enhanced,colback=popgreen,colframe=popink,boxrule=2.2pt,arc=10pt,
        drop shadow=popink,left=11pt,right=11pt,top=10pt,bottom=10pt,height=3.1cm,valign=center]
      \tcbox[colback=white,colframe=popink,boxrule=1.3pt,arc=5pt,boxsep=0pt,nobeforeafter,
        left=3pt,right=3pt,top=3pt,bottom=3pt,valign=center]{\qrcode[height=1.9cm]{https://play.google.com/store/apps/details?id=com.luvvix.onbuch&hl=fr}}%
      \hspace{8pt}\begin{minipage}[c]{0.5\linewidth}
        {\popdisplay\fontsize{11}{12.5}\selectfont\color{white}\MakeUppercase{Télécharge OnBuch}}\par\vspace{4pt}
        {\raggedright\popsemi\fontsize{8.4}{11}\selectfont\color{white}Gratuit \textbullet\ Google Play\\Tuteur IA, annales, concours, résultats\par}
      \end{minipage}
    \end{tcolorbox}
  \end{minipage}\hfill
  \begin{minipage}[t]{0.487\linewidth}
    \begin{tcolorbox}[enhanced,colback=poppurple,colframe=popink,boxrule=2.2pt,arc=10pt,
        drop shadow=popink,left=11pt,right=11pt,top=10pt,bottom=10pt,height=3.1cm,valign=center]
      \tikz[baseline=-0.7ex]\node[circle,fill=white,draw=popink,line width=1.3pt,minimum size=1.6cm,inner sep=0]{\popdisplay\fontsize{24}{24}\selectfont\color{poppurple}+};%
      \hspace{8pt}\begin{minipage}[c]{0.6\linewidth}
        {\popdisplay\fontsize{11}{12.5}\selectfont\color{white}\MakeUppercase{Passe à OnBuch+}}\par\vspace{4pt}
        {\raggedright\popsemi\fontsize{8.4}{11}\selectfont\color{white}Tous les cours signés, Tuteur IA illimité, fiches PDF — onglet OnBuch+ de l'app\par}
      \end{minipage}
    \end{tcolorbox}
  \end{minipage}
  \vspace{8pt}
  \popcontenu
  \vfill
  \signatureonbuch
  \restoregeometry
  \newpage
}

% Rangée « Ce cours contient » (page de garde)
\newcommand{\poptile}[3]{% #1 couleur #2 gros texte #3 légende
  \begin{tcolorbox}[enhanced,colback=#1,colframe=popink,boxrule=2pt,arc=9pt,shadow={2.5pt}{-2.5pt}{0pt}{popink},
      left=6pt,right=6pt,top=9pt,bottom=9pt,halign=center,valign=center,height=1.75cm,width=\linewidth,nobeforeafter]
    {\popdisplay\fontsize{12.5}{13}\selectfont\color{white}\MakeUppercase{#2}}\par\vspace{3pt}
    {\centering\popsemi\fontsize{7.8}{9.5}\selectfont\color{white}#3\par}
  \end{tcolorbox}}
\newcommand{\popcontenu}{%
  \par\noindent{\popxboldls\fontsize{7.5}{7.5}\selectfont\color{popmuted}CE COURS SIGNÉ CONTIENT}\par\vspace{7pt}
  \noindent\begin{minipage}[t]{0.235\linewidth}\poptile{popblue}{Cours}{complet, illustré, pas à pas}\end{minipage}\hfill
  \begin{minipage}[t]{0.235\linewidth}\poptile{poporange}{Méthodes}{et exercices résolus}\end{minipage}\hfill
  \begin{minipage}[t]{0.235\linewidth}\poptile{poppink}{Exercices}{type examen + corrigés}\end{minipage}\hfill
  \begin{minipage}[t]{0.235\linewidth}\poptile{popgreen}{Fiche bilan}{l'essentiel à retenir}\end{minipage}\par}

% Sommaire
\newtcolorbox{sommaire}{enhanced,colback=white,colframe=popink,boxrule=2.2pt,arc=10pt,
  drop shadow=popink,left=18pt,right=18pt,top=13pt,bottom=13pt,before skip=6pt,after skip=20pt,
  before upper={\poppill{Sommaire}{poppurple}{white}\par\vspace{9pt}\popsemi\fontsize{10.6}{14.5}\selectfont\color{popink}}}

% Signature de fin de cours — poussée tout en bas de la dernière page
\newcommand{\findecours}{%
  \par\vfill\nopagebreak
  \begin{center}\popsquiggle{poporange}\end{center}\vspace{4pt}
  \signatureonbuch
}
\AtBeginDocument{\def\llbracket{\mathopen{[\mkern-3.5mu[}}\def\rrbracket{\mathclose{]\mkern-3.5mu]}}} % intervalles d'entiers (glyphe ⟦ absent de la police)
% Glyphes absents de Fira Math : redéfinis à partir de glyphes disponibles
\AtBeginDocument{%
  \def\setminus{\mathbin{\backslash}}\let\smallsetminus\setminus
  \def\vdots{\mathord{\vbox{\baselineskip3.2pt\lineskiplimit0pt\kern4pt\hbox{.}\hbox{.}\hbox{.}}}}%
  \def\ddots{\mathinner{\mkern1mu\raise6.5pt\hbox{.}\mkern2mu\raise3.6pt\hbox{.}\mkern2mu\raise0.7pt\hbox{.}\mkern1mu}}%
  \def\triangle{\mathord{\scalebox{0.9}{$\symup{\Delta}$}}}%
  \def\nabla{\mathord{\raisebox{\depth}{\scalebox{1}[-1]{$\symup{\Delta}$}}}}%
  \def\checkmark{\popcheck{popdark}}%
  \def\star{\mathbin{\ast}}\def\bigstar{\mathord{\ast}}%
  \def\diamond{\mathbin{\scalebox{0.6}{$\Diamond$}}}%
  \def\bigcup{\mathop{\mathchoice{\raisebox{-0.3em}{\scalebox{1.7}{$\cup$}}}{\scalebox{1.25}{$\cup$}}{\cup}{\cup}}\displaylimits}%
  \def\bigcap{\mathop{\mathchoice{\raisebox{-0.3em}{\scalebox{1.7}{$\cap$}}}{\scalebox{1.25}{$\cap$}}{\cap}{\cap}}\displaylimits}%
  \def\lesseqqgtr{\mathrel{\lessgtr}}\def\bigoplus{\mathop{\oplus}\displaylimits}%
}
\providecommand{\ssi}{si et seulement si} % abréviation fréquente
\AtBeginDocument{\providecommand{\wideparen}{\overparen}} % arc de cercle

%% FIGFIX-BEGIN v3 — correctifs globaux des figures (docs/onbuch-generation/tools/figfix)
\usepackage{adjustbox}\usepackage{etoolbox}
% toute figure TikZ/pgfplots plus large que la ligne est réduite à la largeur disponible
\BeforeBeginEnvironment{tikzpicture}{\begin{adjustbox}{max width=\linewidth}}
\AfterEndEnvironment{tikzpicture}{\end{adjustbox}}
% légendes pgfplots lisibles par-dessus les courbes
\pgfplotsset{every axis/.append style={legend style={fill=white,fill opacity=0.92,text opacity=1,draw=black!15,inner sep=3pt}}}
% étiquettes d'arêtes (syntaxe edge["..."]) avec fond blanc
\tikzset{every edge quotes/.append style={fill=white,inner sep=1.2pt}}
% bulles de cartes mentales : texte à la ligne dans la bulle, bulle un peu plus grande
\tikzset{every concept/.append style={text width=2.0cm,align=center,minimum size=2.3cm,inner sep=2pt}}
%% FIGFIX-END
\begin{document}

\pagegarde{Leçon 6 — Grammaire : ……才/ 就...... Et 还 ; 又 ; 再}{Chinois}{Tle A}{Module 1 : Vie familiale et intégration sociale}{ZH06}{2 h}{Cette leçon de grammaire étudie les adverbes 才 et 就 (précocité/tardiveté d'une action) ainsi que 还, 又 et 再 (continuation et répétition). L'élève apprend à construire ces structures, à éviter les pièges de la traduction mot à mot du français et à les employer dans des phrases et de courts textes.
\vspace{4pt}
\begin{multicols}{2}\small
\begin{itemize}
\item À la fin de cette leçon, je sais construire la structure Sujet + 时间 + 就 + Verbe pour exprimer qu'une action a lieu tôt ou facilement
\item À la fin de cette leçon, je sais construire la structure Sujet + 时间 + 才 + Verbe pour exprimer qu'une action a lieu tard ou avec difficulté
\item À la fin de cette leçon, je sais opposer 就 et 才 dans une même situation pour exprimer un jugement (tôt/tard, peu/beaucoup)
\item À la fin de cette leçon, je sais utiliser 还 pour exprimer la continuation d'un état ou d'une action (« encore »)
\item À la fin de cette leçon, je sais utiliser 又 pour exprimer la répétition d'une action déjà réalisée (« de nouveau »)
\item À la fin de cette leçon, je sais utiliser 再 pour exprimer la répétition future ou souhaitée d'une action (« encore une fois »)
\end{itemize}\end{multicols}}

\begin{sommaire}
\begin{multicols}{2}
\begin{enumerate}
\item Observation guidée : 就 et 才 en contexte
\item Emplois et exceptions de 就
\item Emplois et exceptions de 才
\item 还, 又, 再 : continuer et répéter
\item Comparaison avec le français et pièges des francophones
\item Entraînement : application et transformation
\item Activité d'intégration
\item Méthodes et exercices résolus
\item Exercices
\item Corrigés des exercices
\item Fiche bilan
\end{enumerate}
\end{multicols}
\end{sommaire}

\begin{objectifs}
\begin{itemize}
\item À la fin de cette leçon, je sais construire la structure Sujet + 时间 + 就 + Verbe pour exprimer qu'une action a lieu tôt ou facilement
\item À la fin de cette leçon, je sais construire la structure Sujet + 时间 + 才 + Verbe pour exprimer qu'une action a lieu tard ou avec difficulté
\item À la fin de cette leçon, je sais opposer 就 et 才 dans une même situation pour exprimer un jugement (tôt/tard, peu/beaucoup)
\item À la fin de cette leçon, je sais utiliser 还 pour exprimer la continuation d'un état ou d'une action (« encore »)
\item À la fin de cette leçon, je sais utiliser 又 pour exprimer la répétition d'une action déjà réalisée (« de nouveau »)
\item À la fin de cette leçon, je sais utiliser 再 pour exprimer la répétition future ou souhaitée d'une action (« encore une fois »)
\item À la fin de cette leçon, je sais choisir correctement entre 还, 又 et 再 selon le contexte temporel
\item À la fin de cette leçon, je sais repérer et corriger les erreurs fréquentes des francophones (place de l'adverbe, confusion 又/再, traduction de « encore »)
\item À la fin de cette leçon, je sais employer ces cinq adverbes dans un court texte narratif ou un dialogue
\end{itemize}
\end{objectifs}

\begin{prerequis}
\begin{itemize}
\item La phrase de base chinoise Sujet + Verbe + Complément et la place des adverbes avant le verbe (leçons de Première)
\item L'expression de l'heure et du temps : 点, 分, 今天, 明天, 昨天
\item Les adverbes de fréquence et de degré déjà vus : 很, 都, 也, 不, 没
\item La négation au passé avec 没(有) et la particule 了 d'accomplissement
\item Le vocabulaire de la vie familiale et quotidienne du Module 1 (manger, venir, partir, étudier)
\end{itemize}
\end{prerequis}

\newpage

\coursec{Observation guidée : \zh{就} et \zh{才} en contexte}

\courssub{Situation d'accroche : deux petits mots qui changent tout}

Imagine : ton ami te demande quand tu arrives à la répétition du club de chinois à Yaoundé. Tu peux dire « J'arrive à 15 h » — mais en chinois, il y a une nuance : \zh{我三点就到了} (\emph{Wǒ sān diǎn jiù dào le.} « Je serai là dès 15 h, c'est tôt ! ») ou \zh{我五点才到} (\emph{Wǒ wǔ diǎn cái dào.} « Je n'arriverai qu'à 17 h, c'est tard ! »). De même, quand ta mère te demande si tu veux encore du ndolé, tu réponds \zh{我还要} (\emph{Wǒ hái yào.} « J'en veux encore »), \zh{我又吃了} (\emph{Wǒ yòu chī le.} « J'en ai encore mangé ») ou \zh{我再吃一点} (\emph{Wǒ zài chī yìdiǎn.} « Je vais en reprendre »). Ces petits mots \zh{才}, \zh{就}, \zh{还}, \zh{又}, \zh{再} changent tout le sens d'une phrase. Cette leçon t'apprend à les maîtriser comme un vrai sinophone.

\textbf{Problématique :} comment le chinois exprime-t-il, avec un seul petit mot placé avant le verbe, ce que le français dit avec toute une tournure (« dès », « ne... que », « seulement », « il a fallu ») ? Observons d'abord, réglons ensuite.

\courssub{Première observation : deux phrases qui se ressemblent}

Lisons attentivement le mini-dialogue suivant. Deux élèves du lycée de Buea parlent de leur ami Wang, qui devait arriver à 8 h au cours de chinois.

\begin{textebox}[Au lycée : Wang est arrivé quand ?]
\zh{A：小王六点就来了！}\\
\emph{A : Xiǎo Wáng liù diǎn jiù lái le!}\\
« A : Wang est arrivé dès six heures ! (comme c'est tôt !) »\\
\zh{B：不对，他九点才来。}\\
\emph{B : Bú duì, tā jiǔ diǎn cái lái.}\\
« B : Non, il n'est arrivé qu'à neuf heures. (comme c'est tard !) »\\
\zh{A：我五分钟就做完了作业，他两个小时才做完。}\\
\emph{A : Wǒ wǔ fēnzhōng jiù zuòwán le zuòyè, tā liǎng ge xiǎoshí cái zuòwán.}\\
« A : Moi, j'ai fini mes devoirs en cinq minutes ; lui, il lui a fallu deux heures pour finir. »
\end{textebox}

\textbf{Questions d'observation :}
\begin{enumerate}
\item Dans les deux premières répliques, quel mot change entre \zh{六点就来了} et \zh{九点才来} ?
\item Où se placent \zh{就} et \zh{才} par rapport au verbe \zh{来} ?
\item Laquelle des deux phrases exprime l'étonnement devant une arrivée \emph{tôt} ? devant une arrivée \emph{tard} ?
\item Que remarques-tu à la fin de la phrase avec \zh{就} ? Et à la fin de la phrase avec \zh{才} ?
\end{enumerate}

\courssub{Le constat : même place, sens opposés}

Les exemples montrent que \zh{就} (\emph{jiù}) et \zh{才} (\emph{cái}) occupent exactement la même position dans la phrase, mais expriment des jugements opposés du locuteur sur le moment, la durée ou la quantité.

\begin{propriete}[Place de \zh{就} et \zh{才} dans la phrase]
\zh{就} et \zh{才} sont des \cle{adverbes} : ils se placent \textbf{après le sujet} (et après le complément de temps ou de quantité) et \textbf{immédiatement avant le verbe}.\\
Schéma : \textbf{Sujet + (temps / quantité) + 就 / 才 + Verbe (+ 了)}\\
\zh{他 六点 就 来 了。} \emph{Tā liù diǎn jiù lái le.} « Il est arrivé dès six heures. »\\
\zh{他 九点 才 来。} \emph{Tā jiǔ diǎn cái lái.} « Il n'est arrivé qu'à neuf heures. »
\end{propriete}

\begin{popfigure}
\begin{tikzpicture}[x=1.15cm]
\draw[-{Stealth}, thick, popdark] (0,0) -- (11.2,0) node[below left=2pt and -30pt, popmuted] {\footnotesize temps};
\foreach \x/\h in {2/6 h, 5/8 h, 8/10 h} {
  \draw[popdark] (\x,0.12) -- (\x,-0.12);
  \node[below, popdark, font=\footnotesize] at (\x,-0.12) {\h};
}
\draw[dashed, thick, poporange] (5,0.15) -- (5,1.5);
\node[above, poporange, font=\footnotesize\bfseries, align=center] at (5,1.5) {heure prévue\\8 h};
\draw[-{Stealth}, very thick, popgreen] (2,0.15) -- (2,1.1);
\node[above, popgreen, font=\footnotesize\bfseries, align=center] at (2,1.1) {就 : 6 h\\« dès 6 h » (tôt)};
\draw[-{Stealth}, very thick, poppink] (8,0.15) -- (8,1.1);
\node[above, poppink, font=\footnotesize\bfseries, align=center] at (8,1.1) {才 : 10 h\\« seulement à 10 h » (tard)};
\end{tikzpicture}
\legende{Frise des temps : \zh{就} signale une action perçue comme tôt (avant le repère attendu), \zh{才} une action perçue comme tard (après le repère attendu).}
\end{popfigure}

\begin{aretenir}[L'idée générale]
\begin{itemize}
\item \zh{就} = l'action est perçue comme \cle{tôt}, \cle{rapide}, \cle{facile}, ou la quantité comme \cle{petite} (« dès », « en seulement », « déjà »).
\item \zh{才} = l'action est perçue comme \cle{tardive}, \cle{lente}, \cle{difficile}, ou la quantité comme \cle{insuffisante} (« ne... que », « il a fallu... pour »).
\end{itemize}
Le choix entre \zh{就} et \zh{才} dépend donc du \textbf{point de vue du locuteur}, pas d'une règle d'heure absolue : 6 h du matin, c'est « tôt » pour un cours, mais « tard » si l'on attendait quelqu'un à 5 h !
\end{aretenir}

\begin{popfigure}
\begin{tikzpicture}[node distance=0.35cm, every node/.style={font=\small}]
\node[draw=popblue, fill=popblueL, rounded corners=2pt, minimum height=0.8cm] (s) {Sujet};
\node[draw=popgreen, fill=popgreenL, rounded corners=2pt, minimum height=0.8cm, right=of s, align=center] (t){时间 / 数量\\[-2pt]\footnotesize temps / quantité};
\node[draw=poporange, fill=poporangeL, rounded corners=2pt, minimum height=0.8cm, right=of t] (j) {就 / 才};
\node[draw=poppurple, fill=poppurpleL, rounded corners=2pt, minimum height=0.8cm, right=of j] (v) {Verbe};
\node[draw=popmuted, fill=popcream, rounded corners=2pt, minimum height=0.8cm, right=of v] (l) {(了)};
\draw[-{Stealth}, thick, popdark] (s) -- (t);
\draw[-{Stealth}, thick, popdark] (t) -- (j);
\draw[-{Stealth}, thick, popdark] (j) -- (v);
\draw[-{Stealth}, thick, popdark] (v) -- (l);
\node[below=0.25cm of j, popmuted, font=\footnotesize] {adverbe : place fixe};
\node[below=0.25cm of l, popmuted, font=\footnotesize] {了 seulement avec 就};
\end{tikzpicture}
\legende{Schéma de la structure : [Sujet] → [temps / quantité] → [\zh{就} / \zh{才}] → [Verbe] → [(了)].}
\end{popfigure}

\courssub{Observation avec les quantités et les durées}

Le même contraste fonctionne avec les durées et les quantités. Compare :

\begin{exemplebox}[Durées et quantités contrastées]
\zh{我五分钟就做完了。}\\
\emph{Wǒ wǔ fēnzhōng jiù zuòwán le.} « J'ai fini en seulement cinq minutes. » (rapide !)\\
\zh{我两个小时才做完。}\\
\emph{Wǒ liǎng ge xiǎoshí cái zuòwán.} « Il m'a fallu deux heures pour finir. » (long !)\\
\zh{我才学了两年汉语。}\\
\emph{Wǒ cái xué le liǎng nián Hànyǔ.} « Je n'apprends le chinois que depuis deux ans. » (c'est peu, à mon avis)\\
\zh{他就学了两年，说得很好！}\\
\emph{Tā jiù xué le liǎng nián, shuō de hěn hǎo!} « Seulement deux ans d'études, et il parle déjà très bien ! » (c'est remarquable)
\end{exemplebox}

Remarque bien : dans \zh{我才学了两年}, le \zh{了} suit le verbe \zh{学} (action accomplie avec durée), mais il ne suit \textbf{jamais} directement \zh{才}. On dit \zh{他昨天才来} (\emph{Tā zuótiān cái lái.} « Il n'est arrivé qu'hier »), jamais \zh{他才来了}.

\courssub{Observation avec condition et conséquence}

\zh{就} possède un emploi supplémentaire très fréquent : il introduit la \textbf{conséquence} d'une condition, comme « dès que..., alors... » en français.

\begin{exemplebox}[\zh{就} de conséquence logique]
\zh{你来了，我就放心了。}\\
\emph{Nǐ lái le, wǒ jiù fàngxīn le.} « Dès que tu es arrivé, j'ai été rassuré. »\\
\zh{下雨我就不去。}\\
\emph{Xià yǔ wǒ jiù bú qù.} « S'il pleut, je n'irai pas. »\\
\zh{你做完作业就可以玩。}\\
\emph{Nǐ zuòwán zuòyè jiù kěyǐ wán.} « Quand tu auras fini tes devoirs, tu pourras jouer. »
\end{exemplebox}

Dans cet emploi, \zh{才} exprime la condition restrictive, équivalente de « seulement si » : \zh{你做完作业才能玩。} (\emph{Nǐ zuòwán zuòyè cái néng wán.} « Tu ne pourras jouer que lorsque tes devoirs seront finis. »). Nous y reviendrons dans la section 3.

\courssub{Tableau récapitulatif de l'observation}

\begin{center}
\begin{tabularx}{\linewidth}{|X|X|X|}
\hline
\rowcolor{popblueL} \textbf{Structure} & \textbf{Exemple (caractères + pinyin)} & \textbf{Traduction} \\
\hline
Sujet + heure + 就 + V + 了 & \zh{他六点就来了。} \emph{Tā liù diǎn jiù lái le.} & Il est arrivé dès 6 h (tôt). \\
\hline
Sujet + heure + 才 + V & \zh{他九点才来。} \emph{Tā jiǔ diǎn cái lái.} & Il n'est arrivé qu'à 9 h (tard). \\
\hline
Sujet + durée + 就 + V + 了 & \zh{我五分钟就做完了。} \emph{Wǒ wǔ fēnzhōng jiù zuòwán le.} & J'ai fini en cinq minutes (rapide). \\
\hline
Sujet + durée + 才 + V & \zh{我两个小时才做完。} \emph{Wǒ liǎng ge xiǎoshí cái zuòwán.} & Il m'a fallu deux heures (long). \\
\hline
Sujet + 才 / 就 + V + 了 + quantité & \zh{我才学了两年汉语。} \emph{Wǒ cái xué le liǎng nián Hànyǔ.} & Je n'apprends que depuis deux ans. \\
\hline
Condition，sujet + 就 + V & \zh{你来了，我就放心了。} \emph{Nǐ lái le, wǒ jiù fàngxīn le.} & Dès que tu es arrivé, j'ai été rassuré. \\
\hline
\end{tabularx}
\end{center}

\begin{center}
\begin{tabularx}{\linewidth}{|X|X|X|X|}
\hline
\rowcolor{popblueL} \textbf{Caractères} & \textbf{Pinyin} & \textbf{Nature} & \textbf{Traduction} \\
\hline
就 & jiù & adverbe & dès, aussitôt, alors, seulement \\
\hline
才 & cái & adverbe & seulement, ne... que, à peine \\
\hline
放心 & fàngxīn & verbe & être rassuré \\
\hline
做完 & zuòwán & verbe + complément & finir de faire \\
\hline
作业 & zuòyè & nom & devoirs \\
\hline
小时 & xiǎoshí & nom (mesure) & heure (durée) \\
\hline
\end{tabularx}
\end{center}

\courssub{Comparaison avec le français et premiers pièges}

En français, nous utilisons des tournures très différentes : « dès 6 h », « il n'est arrivé qu'à 9 h », « il m'a fallu deux heures ». Le chinois, lui, garde le même ordre de mots et ne change qu'\textbf{un seul adverbe}. C'est économique, mais cela exige de bien choisir entre \zh{就} et \zh{才} selon le jugement porté.

\begin{attention}[Le piège de \zh{只} : « seulement » ne se traduit pas partout pareil]
Les francophones traduisent souvent « seulement » par \zh{只} (\emph{zhǐ}) partout. Or, pour dire « il n'est arrivé qu'à 9 h », le chinois dit \zh{他九点才来}, et non \zh{他只九点来} (phrase incorrecte). \zh{只} limite une \textbf{quantité ou un objet} (\zh{我只有一个弟弟。} \emph{Wǒ zhǐ yǒu yí ge dìdi.} « Je n'ai qu'un petit frère »), tandis que \zh{才} limite un \textbf{moment ou une durée} devant le verbe. Autre piège : ne mets jamais \zh{了} juste après \zh{才} pour le passé : on dit \zh{他昨天才来}, jamais \zh{他昨天才来了}.
\end{attention}

\begin{exoresolu}[À toi de choisir : \zh{就} ou \zh{才} ?]
Énoncé. Complète avec \zh{就} ou \zh{才}, puis traduis :\\
1. Le cours commence à 8 h ; Marie arrive à 7 h 30. 她七点半\_\_\_来了。\\
2. Le cours commence à 8 h ; Paul arrive à 9 h. 他九点\_\_\_来。\\
3. 你好好学习，将来\_\_\_能去中国。 (Si tu étudies bien, un jour tu pourras aller en Chine.)
\tcblower
Solution détaillée.\\
1. \zh{她七点半就来了。} \emph{Tā qī diǎn bàn jiù lái le.} « Elle est arrivée dès 7 h 30. » — 7 h 30 est \emph{avant} 8 h : arrivée perçue comme tôt, donc \zh{就}, avec \zh{了} final car l'action est accomplie.\\
2. \zh{他九点才来。} \emph{Tā jiǔ diǎn cái lái.} « Il n'est arrivé qu'à 9 h. » — 9 h est \emph{après} 8 h : arrivée perçue comme tardive, donc \zh{才}, et \textbf{pas de} \zh{了} final.\\
3. \zh{你好好学习，将来就能去中国。} \emph{Nǐ hǎohǎo xuéxí, jiānglái jiù néng qù Zhōngguó.} « Si tu étudies sérieusement, tu pourras aller en Chine plus tard. » — Ici \zh{就} introduit la conséquence logique de la condition.
\end{exoresolu}

\begin{savaistu}[Un champion de la fréquence]
\zh{就} est l'un des adverbes les plus fréquents de toute la langue chinoise : on le rencontre dans presque chaque conversation, car il exprime à la fois la précocité (« dès »), la conséquence logique (« alors ») et la restriction (« seulement »). Les Chinois l'emploient même seul pour acquiescer vivement : \zh{就这样！} (\emph{Jiù zhèyàng!} « C'est ça, exactement ! »). À l'inverse, \zh{才} porte souvent une nuance d'insatisfaction ou d'admiration retenue : dire \zh{你怎么才来！} (\emph{Nǐ zěnme cái lái!} « Mais enfin, te voilà enfin ! ») à un ami en retard au marché de Douala est tout à fait naturel.
\end{savaistu}

\begin{exercice}[Observation et jugement]
Pour chaque situation, choisis \zh{就} ou \zh{才} et justifie ton choix en français.\\
1. L'examen de chinois commence à 8 h. Aïcha arrive à 7 h. 她七点\_\_\_到了。\\
2. Le car pour Bafoussam part à 10 h. Il part en réalité à 12 h. 车十二点\_\_\_走。\\
3. Mon petit frère a fini son riz en deux minutes ! 他两分钟\_\_\_吃完了。\\
4. J'ai attendu le taxi pendant une heure. 我等了一个小时\_\_\_上车。
\end{exercice}

\begin{corrige}[Exercice « Observation et jugement »]
1. \zh{就} : \zh{她七点就到了。} « Elle est arrivée dès 7 h. » — arrivée avant l'heure prévue, donc perçue comme tôt ; \zh{了} final car l'action est accomplie.\\
2. \zh{才} : \zh{车十二点才走。} « Le car n'est parti qu'à midi. » — départ après l'heure prévue, perçu comme tard ; pas de \zh{了} après \zh{才}.\\
3. \zh{就} : \zh{他两分钟就吃完了。} « Il a fini de manger en deux minutes ! » — durée très courte, action rapide.\\
4. \zh{才} : \zh{我等了一个小时才上车。} « J'ai attendu une heure avant de pouvoir monter. » — attente longue, action retardée ; le \zh{了} suit \zh{等}, jamais \zh{才}.
\end{corrige}

\begin{aretenir}[Bilan de l'observation]
\begin{itemize}
\item \zh{就} et \zh{才} se placent au même endroit : après le sujet et le complément de temps/quantité, juste avant le verbe.
\item \zh{就} juge l'action tôt, rapide, facile, peu ; \zh{才} la juge tardive, lente, difficile, seulement.
\item Action accomplie avec \zh{就} : \zh{了} en fin de phrase (\zh{就来了}). Avec \zh{才} : jamais \zh{了} juste après (\zh{才来}, pas \zh{才来了}).
\item \zh{就} exprime aussi la conséquence logique : \zh{下雨我就不去。}
\end{itemize}
\end{aretenir}

\coursec{Emplois et exceptions de 就}

Dans la section précédente, nous avons observé \zh{就} et \zh{才} en contexte : deux petits mots placés devant le verbe, qui expriment des jugements opposés sur le temps ou la quantité. Nous savons maintenant que \zh{就} dit « plus tôt / plus vite que prévu » et \zh{才} « plus tard / plus lentement que prévu ». Cette section approfondit \textbf{tous les emplois de \zh{就}}, un par un, avec leurs structures, leurs exemples et leurs pièges. C'est l'un des adverbes les plus fréquents du chinois : le maîtriser, c'est gagner en naturel à l'oral comme à l'écrit.

\courssub{Emploi 1 : la précocité — « plus tôt que prévu »}

\textbf{Observation.} Lisez ces deux phrases :

\begin{exemplebox}[Observer avant de comprendre]
\zh{她十五岁就上大学了。}\par
\emph{Tā shíwǔ suì jiù shàng dàxué le.}\par
« Elle est entrée à l'université dès quinze ans. »\par
\medskip
\zh{我弟弟六岁就会游泳了。}\par
\emph{Wǒ dìdi liù suì jiù huì yóuyǒng le.}\par
« Mon petit frère savait déjà nager à six ans. »
\end{exemplebox}

Dans les deux cas, le locuteur juge que l'événement arrive \textbf{plus tôt que ce qu'on attend normalement} : on entre à l'université vers dix-huit ou dix-neuf ans, on apprend à nager vers huit ou dix ans. \zh{就} porte ce jugement d'étonnement positif.

\begin{propriete}[就 de précocité : la structure]
\textbf{Sujet + (complément de temps) + 就 + Verbe + 了}\par
\begin{itemize}
\item \zh{就} se place \textbf{après le sujet} (et après le complément de temps s'il y en a un) et \textbf{juste devant le verbe}.
\item Le \zh{了} final marque que l'action est \textbf{accomplie} : on parle d'un fait passé.
\item Sens : « dès… », « déjà… (à tel âge, à telle heure) ».
\end{itemize}
\end{propriete}

\begin{exemplebox}[Exemples]
\zh{他十七岁就中学毕业了。}\par
\emph{Tā shíqī suì jiù zhōngxué bìyè le.}\par
« Il a fini le lycée dès dix-sept ans. »\par
\medskip
\zh{我们早上五点就出发了。}\par
\emph{Wǒmen zǎoshang wǔ diǎn jiù chūfā le.}\par
« Nous sommes partis dès cinq heures du matin. »\par
\medskip
\zh{她去年就会说汉语了。}\par
\emph{Tā qùnián jiù huì shuō Hànyǔ le.}\par
« Elle savait déjà parler chinois l'année dernière. »
\end{exemplebox}

\textbf{Comparaison avec le français.} Le français rend cet emploi par « dès », « déjà », ou un ton d'insistance : « Elle est entrée à l'université \emph{dès} quinze ans ». Le chinois, lui, \textbf{oblige} à marquer ce jugement par \zh{就} : sans lui, la phrase devient un simple constat neutre. Comparez :

\begin{center}
\begin{tabularx}{\linewidth}{|X|X|X|}\hline
\rowcolor{popblueL} Phrase & Pinyin & Nuance \\ \hline
\zh{她十五岁上大学了。} & \emph{Tā shíwǔ suì shàng dàxué le.} & Constat neutre : elle est entrée à l'université à 15 ans. \\ \hline
\zh{她十五岁\cle{就}上大学了。} & \emph{Tā shíwǔ suì jiù shàng dàxué le.} & Jugement : c'est \textbf{tôt}, remarquable ! \\ \hline
\zh{她十五岁\cle{才}上大学。} & \emph{Tā shíwǔ suì cái shàng dàxué.} & (structure de \zh{才}, vue en section 3 : jugement inverse) \\ \hline
\end{tabularx}
\end{center}

\courssub{Emploi 2 : la rapidité et la facilité — « vite, sans difficulté »}

\textbf{Observation.}

\begin{exemplebox}[Observer avant de comprendre]
\zh{这个汉字很简单，我一下就学会了。}\par
\emph{Zhège Hànzì hěn jiǎndān, wǒ yīxià jiù xuéhuì le.}\par
« Ce caractère est très simple, je l'ai appris tout de suite. »\par
\medskip
\zh{从我家到学校，走路十分钟就到了。}\par
\emph{Cóng wǒ jiā dào xuéxiào, zǒulù shí fēnzhōng jiù dào le.}\par
« De chez moi à l'école, dix minutes à pied suffisent pour arriver. »
\end{exemplebox}

Ici, \zh{就} ne porte plus sur une date ou un âge, mais sur la \textbf{durée} ou la \textbf{difficulté} : l'action est accomplie \textbf{plus vite ou plus facilement que prévu}. On le trouve souvent avec des expressions comme \zh{一下} (en un rien de temps), \zh{五分钟} (cinq minutes), \zh{很容易} (très facile).

\begin{propriete}[就 de rapidité / facilité]
\textbf{Sujet + (durée courte / 一下) + 就 + Verbe + 了}\par
Sens : « en seulement… », « tout de suite », « sans difficulté ». Le locuteur souligne que la tâche a demandé \textbf{peu de temps ou peu d'effort}.
\end{propriete}

\begin{exemplebox}[Exemples]
\zh{这个菜很容易做，我五分钟就做好了。}\par
\emph{Zhège cài hěn róngyì zuò, wǒ wǔ fēnzhōng jiù zuòhǎo le.}\par
« Ce plat est facile à préparer, je l'ai fait en cinq minutes. »\par
\medskip
\zh{从杜阿拉到雅温得，坐车三个小时就到了。}\par
\emph{Cóng Dù'ālā dào Yǎwēndé, zuò chē sān ge xiǎoshí jiù dào le.}\par
« De Douala à Yaoundé, trois heures de voiture suffisent. »\par
\medskip
\zh{这道题不难，他一下就做对了。}\par
\emph{Zhè dào tí bù nán, tā yīxià jiù zuòduì le.}\par
« Cet exercice n'est pas difficile, il l'a résolu correctement en un rien de temps. »
\end{exemplebox}

\courssub{Emploi 3 : la conséquence immédiate — 一……就…… « dès que »}

\textbf{Observation.}

\begin{exemplebox}[Observer avant de comprendre]
\zh{我一回家就做作业。}\par
\emph{Wǒ yī huí jiā jiù zuò zuòyè.}\par
« Dès que je rentre à la maison, je fais mes devoirs. »\par
\medskip
\zh{他一毕业就找到了工作。}\par
\emph{Tā yī bìyè jiù zhǎodào le gōngzuò.}\par
« Dès qu'il a eu son diplôme, il a trouvé un travail. »
\end{exemplebox}

\begin{propriete}[一 + Verbe1 + 就 + Verbe2 = « dès que…, … »]
\textbf{Sujet + 一 + Verbe1 + 就 + Verbe2}\par
\begin{itemize}
\item Les deux actions \textbf{s'enchaînent immédiatement} : dès que la première est faite, la seconde commence.
\item \zh{一} se place devant le premier verbe, \zh{就} devant le second.
\item Si les deux verbes ont des sujets différents, le second sujet se place \textbf{avant} \zh{就} : \zh{我一到，他就走了。} \emph{Wǒ yī dào, tā jiù zǒu le.} « Dès que je suis arrivé, il est parti. »
\item Cette structure décrit aussi bien une habitude (présent) qu'un enchaînement passé (avec \zh{了}).
\end{itemize}
\end{propriete}

\begin{exemplebox}[Exemples]
\zh{妈妈一下班就做饭。}\par
\emph{Māma yī xiàbān jiù zuòfàn.}\par
« Dès que maman finit le travail, elle cuisine. »\par
\medskip
\zh{我们一考完试就放暑假。}\par
\emph{Wǒmen yī kǎowán shì jiù fàng shǔjià.}\par
« Dès que les examens sont terminés, les grandes vacances commencent. »\par
\medskip
\zh{他一看足球比赛就高兴。}\par
\emph{Tā yī kàn zúqiú bǐsài jiù gāoxìng.}\par
« Dès qu'il regarde un match de football, il est content. »
\end{exemplebox}

\textbf{Comparaison avec le français.} Le français « dès que » introduit une subordonnée ; le chinois, lui, encadre les deux verbes avec \zh{一} et \zh{就}, comme deux crochets. Ne traduisez jamais « dès que » par \zh{就} seul : il faut le couple \zh{一……就……}.

\courssub{Emploi 4 : la petitesse d'une quantité — « seulement, ne… que »}

\textbf{Observation.}

\begin{exemplebox}[Observer avant de comprendre]
\zh{我们班就有五个学生学汉语。}\par
\emph{Wǒmen bān jiù yǒu wǔ ge xuésheng xué Hànyǔ.}\par
« Notre classe n'a que cinq élèves qui apprennent le chinois. »\par
\medskip
\zh{我就带了十块钱。}\par
\emph{Wǒ jiù dài le shí kuài qián.}\par
« Je n'ai apporté que dix yuans. »
\end{exemplebox}

Ici, \zh{就} juge une \textbf{quantité} : elle est \textbf{plus petite que prévu}. C'est l'équivalent du français « seulement », « ne… que ». (On retrouve l'idée opposée avec \zh{才} dans d'autres contextes ; pour les quantités jugées insuffisantes, \zh{就} insiste sur le jugement du locuteur : « c'est peu ! ».)

\begin{propriete}[就 + quantité = « seulement »]
\textbf{Sujet + 就 + Verbe + (quantité petite)}\par
Le locuteur insiste sur la \textbf{petitesse} du nombre ou de la quantité. Souvent renforcé par le ton de la voix sur \zh{就}.
\end{propriete}

\begin{exemplebox}[Exemples]
\zh{这个学校就有两位汉语老师。}\par
\emph{Zhège xuéxiào jiù yǒu liǎng wèi Hànyǔ lǎoshī.}\par
« Cette école n'a que deux professeurs de chinois. »\par
\medskip
\zh{我就学了一年汉语。}\par
\emph{Wǒ jiù xué le yī nián Hànyǔ.}\par
« Je n'ai appris le chinois que pendant un an. »
\end{exemplebox}

\courssub{Emploi 5 : l'affirmation — « c'est précisément »}

\textbf{Observation.}

\begin{exemplebox}[Observer avant de comprendre]
\zh{这就是我的老师。}\par
\emph{Zhè jiù shì wǒ de lǎoshī.}\par
« C'est précisément mon professeur. » / « Voici mon professeur. »\par
\medskip
\zh{雅温得就是喀麦隆的首都。}\par
\emph{Yǎwēndé jiù shì Kāmàilóng de shǒudū.}\par
« Yaoundé, c'est précisément la capitale du Cameroun. »
\end{exemplebox}

Devant \zh{是} (ou un verbe), \zh{就} exprime une \textbf{affirmation insistante} : « c'est bien… », « c'est exactement… ». On l'utilise pour confirmer, identifier ou corriger.

\begin{exemplebox}[Exemples]
\zh{我就要这本书。}\par
\emph{Wǒ jiù yào zhè běn shū.}\par
« C'est précisément ce livre que je veux. »\par
\medskip
\zh{他就住在市场旁边。}\par
\emph{Tā jiù zhù zài shìchǎng pángbiān.}\par
« Il habite juste à côté du marché. »
\end{exemplebox}

\courssub{Exception : 就 devant un verbe d'action future = « tout de suite »}

\begin{propriete}[Exception : 就 + verbe d'action (futur proche)]
Devant un verbe d'action \textbf{non accomplie}, \zh{就} signifie « tout de suite, immédiatement » :\par
\medskip
\zh{我就来！}\par
\emph{Wǒ jiù lái!}\par
« J'arrive tout de suite ! »\par
\medskip
\zh{你等一下，我就去。}\par
\emph{Nǐ děng yīxià, wǒ jiù qù.}\par
« Attends un peu, j'y vais tout de suite. »
\end{propriete}

C'est l'emploi qu'on entend chaque jour en Chine : quelqu'un vous appelle, vous répondez \zh{我就来！}. Ici, pas de \zh{了} final, car l'action n'est pas encore faite.

\begin{attention}[Ne pas confondre 就 (tôt) et 马上 (immédiatement)]
\begin{itemize}
\item \zh{就} \textbf{juge} une heure ou une durée \emph{par rapport à une attente} : « plus tôt que prévu ». Il regarde vers le passé ou un repère : \zh{他五点就来了。} « Il est arrivé dès cinq heures (c'était tôt). »
\item \zh{马上} \textbf{annonce} une action \emph{imminente}, sans jugement : \zh{我马上来。} « J'arrive immédiatement. »
\item Les deux peuvent se combiner : \zh{我马上就来。} « J'arrive vraiment tout de suite. » Mais remplacer \zh{就} par \zh{马上} dans \zh{她十五岁就上大学了} est une faute de sens : on perd le jugement de précocité.
\end{itemize}
\end{attention}

\begin{methode}[Comment choisir 就 ?]
\begin{enumerate}
\item Repérez ce que le locuteur \textbf{juge} : une heure, un âge, une durée, une quantité, une identification.
\item Posez-vous la question : « est-ce \textbf{plus tôt / plus vite / plus facile / moins} que prévu, ou une affirmation \textbf{précise} ? »
\item Si oui, placez \zh{就} \textbf{après le sujet et juste devant le verbe} (ou devant \zh{是}).
\item Si l'action est accomplie, ajoutez \zh{了} ; si c'est une promesse d'action immédiate, pas de \zh{了} (\zh{我就来！}).
\item Si deux actions s'enchaînent, utilisez le couple \zh{一……就……}.
\end{enumerate}
\end{methode}

\courssub{Tableau récapitulatif des emplois de 就}

\begin{center}
\begin{tabularx}{\linewidth}{|X|X|X|}\hline
\rowcolor{popblueL} Emploi & Exemple (caractères + pinyin) & Traduction \\ \hline
1. Précocité & \zh{她十五岁就上大学了。} \emph{Tā shíwǔ suì jiù shàng dàxué le.} & Elle est entrée à l'université dès 15 ans. \\ \hline
2. Rapidité / facilité & \zh{我一下就学会了。} \emph{Wǒ yīxià jiù xuéhuì le.} & Je l'ai appris tout de suite. \\ \hline
3. Dès que & \zh{我一回家就做作业。} \emph{Wǒ yī huí jiā jiù zuò zuòyè.} & Dès que je rentre, je fais mes devoirs. \\ \hline
4. Seulement & \zh{我们班就有五个学生学汉语。} \emph{Wǒmen bān jiù yǒu wǔ ge xuésheng xué Hànyǔ.} & Notre classe n'a que cinq élèves de chinois. \\ \hline
5. Précisément & \zh{这就是我的老师。} \emph{Zhè jiù shì wǒ de lǎoshī.} & C'est précisément mon professeur. \\ \hline
Exception : futur proche & \zh{我就来！} \emph{Wǒ jiù lái!} & J'arrive tout de suite ! \\ \hline
\end{tabularx}
\end{center}

\begin{popfigure}
\begin{tikzpicture}[mindmap, grow cyclic, every node/.style={concept}, concept color=popblue!40, text=popdark,
level 1/.append style={level distance=3.4cm, sibling angle=72},
level 2/.append style={level distance=2.6cm, sibling angle=50, concept color=poporange!35}]
\node[concept color=popblue!60, font=\Large\bfseries] {就 jiù}
child { node {tôt, dès} child { node[align=center, font=\small] {她十五岁就\\上大学了。} } }
child { node {vite, facile} child { node[align=center, font=\small] {我一下就\\学会了。} } }
child { node {dès que} child { node[align=center, font=\small] {我一回家\\就做作业。} } }
child { node {seulement} child { node[align=center, font=\small] {就有五个学生\\学汉语。} } }
child { node {précisément} child { node[align=center, font=\small] {这就是我的\\老师。} } };
\end{tikzpicture}
\legende{Carte mentale des cinq emplois de 就, avec un exemple par branche.}
\end{popfigure}

\begin{exoresolu}[Application : traduire et justifier]
Énoncé. Traduisez en chinois : « Dès que mon frère finit ses devoirs, il regarde la télévision. » Puis : « Je n'ai que dix yuans dans ma poche. » Indiquez quel emploi de \zh{就} vous utilisez.\par
\tcblower
Solution détaillée.\par
1. « Dès que… » appelle la structure \zh{一……就……} (emploi 3) : les deux actions s'enchaînent. Le sujet est le même pour les deux verbes, donc :\par
\zh{我弟弟一做完作业就看电视。}\par
\emph{Wǒ dìdi yī zuòwán zuòyè jiù kàn diànshì.}\par
Notez \zh{做完} « finir de faire » : le complément de résultat \zh{完} marque que la première action est terminée avant que la seconde commence.\par
2. « Ne… que dix yuans » exprime une quantité jugée petite : emploi 4, \zh{就} devant le verbe \zh{有} :\par
\zh{我口袋里就有十块钱。}\par
\emph{Wǒ kǒudài li jiù yǒu shí kuài qián.}\par
« Je n'ai que dix yuans dans ma poche. »\par
Remarque : pour un \textbf{âge} jugé petit (« il n'a que huit ans »), le chinois préfère \zh{才} : \zh{我弟弟才八岁。} — nous verrons cette structure en section 3. Retenez dès maintenant : \zh{就} juge « c'est peu » pour une quantité possédée ou constatée, \zh{才} juge « c'est peu / c'est tard » pour un âge, une heure ou une quantité attendue plus grande.
\end{exoresolu}

\begin{attention}[Erreurs fréquentes des francophones avec 就]
\begin{itemize}
\item \textbf{Mauvaise place.} \zh{就} ne se met jamais devant le sujet. Faux : \zh{就她十五岁上大学了。} Correct : \zh{她十五岁就上大学了。}
\item \textbf{Oublier 了.} Pour un fait accompli jugé tôt, le \zh{了} final est presque obligatoire : \zh{他五点就到了。}, et non \zh{他五点就到}.
\item \textbf{Traduire « seulement » par 只 partout.} \zh{只} est le mot neutre « seulement » ; \zh{就} ajoute un jugement (« c'est peu ! »). Les deux sont possibles, mais ne les cumulez pas dans les exercices de ce niveau.
\item \textbf{Confondre 一……就…… et 就 seul.} « Dès que » exige les deux mots ; \zh{就} seul devant le second verbe sans \zh{一} devant le premier ne suffit pas.
\end{itemize}
\end{attention}

\begin{exercice}[À vous de jouer]
Traduisez en chinois en utilisant \zh{就} (précisez l'emploi) :
\begin{enumerate}
\item Elle sait lire les caractères depuis l'âge de cinq ans.
\item Ce restaurant est tout près : cinq minutes à pied suffisent.
\item Dès que le professeur entre, les élèves se lèvent.
\item Je n'ai que trois cours de chinois par semaine.
\item C'est précisément ici que se trouve le marché de Yaoundé.
\end{enumerate}
\end{exercice}

\begin{corrige}[Exercice « À vous de jouer »]
\begin{enumerate}
\item \zh{她五岁就会认汉字了。} \emph{Tā wǔ suì jiù huì rèn Hànzì le.} (emploi 1, précocité)
\item \zh{这个饭馆很近，走路五分钟就到了。} \emph{Zhège fànguǎn hěn jìn, zǒulù wǔ fēnzhōng jiù dào le.} (emploi 2, rapidité)
\item \zh{老师一进来，学生就站起来。} \emph{Lǎoshī yī jìnlái, xuésheng jiù zhàn qǐlái.} (emploi 3 ; sujets différents : le second sujet précède \zh{就})
\item \zh{我一个星期就有三节汉语课。} \emph{Wǒ yī ge xīngqī jiù yǒu sān jié Hànyǔ kè.} (emploi 4, quantité petite)
\item \zh{雅温得的市场就在这里。} \emph{Yǎwēndé de shìchǎng jiù zài zhèlǐ.} (emploi 5, affirmation)
\end{enumerate}
\end{corrige}

\begin{aretenir}[L'essentiel sur 就]
\begin{itemize}
\item \zh{就} se place \textbf{après le sujet, devant le verbe} ; il exprime un jugement : \textbf{tôt, vite, facile, peu, précisément}.
\item Fait accompli : ajoutez \zh{了}. Action immédiate à venir : \zh{我就来！} sans \zh{了}.
\item « Dès que » = \zh{一 + V1 + 就 + V2} ; si les sujets diffèrent, le second sujet passe avant \zh{就}.
\item \zh{就} juge par rapport à une attente ; \zh{马上} annonce simplement une action imminente.
\end{itemize}
\end{aretenir}

\coursec{Emplois et exceptions de 才}

Nous venons d'étudier \zh{就} (jiù), l'adverbe de la \emph{précocité} et de la facilité : « plus tôt que prévu », « plus facile que prévu », « dès que ». Nous abordons maintenant son \cle{miroir exact} : \zh{才} (cái). Retenez dès maintenant l'idée directrice de cette section : \textbf{là où 就 dit « plus tôt / plus facile / plus que prévu », 才 dit « plus tard / plus difficile / moins que prévu »}. Les deux adverbes occupent exactement la même place dans la phrase (avant le verbe, après le sujet et le complément de temps), mais ils expriment des jugements opposés. Maîtriser 才, c'est donc d'abord bien comprendre 就, puis renverser la perspective.

\courssub{Observation : un même fait, deux jugements}

\begin{textebox}[Dialogue — Le rendez-vous manqué]
A：我们约好八点见面，你几点到的？\\
\emph{Wǒmen yuē hǎo bā diǎn jiànmiàn, nǐ jǐ diǎn dào de?}\\
« Nous avions rendez-vous à huit heures, tu es arrivé à quelle heure ? »\\[4pt]
B：我七点半就到了，你九点才来！我等了你一个半小时！\\
\emph{Wǒ qī diǎn bàn jiù dào le, nǐ jiǔ diǎn cái lái! Wǒ děng le nǐ yí ge bàn xiǎoshí!}\\
« Moi, je suis arrivé dès sept heures et demie, et toi tu n'es arrivé qu'à neuf heures ! Je t'ai attendu une heure et demie ! »\\[4pt]
A：对不起，公交车晚点了。我吃了晚饭才能出门。\\
\emph{Duìbuqǐ, gōngjiāo chē wǎndiǎn le. Wǒ chī le wǎnfàn cái néng chūmén.}\\
« Pardon, le bus a eu du retard. Je n'ai pu sortir qu'après avoir dîné. »
\end{textebox}

\textbf{Observons.} Le rendez-vous est fixé à \zh{八点} (8 h). Le premier locuteur arrive à 7 h 30 : il dit \zh{七点半就到了} — \cle{就}, car c'est \emph{plus tôt} que prévu. Le second arrive à 9 h : il dit \zh{九点才来} — \cle{才}, car c'est \emph{plus tard} que prévu. Même structure, même place de l'adverbe, mais jugement inverse. C'est toute la logique de 才.

\courssub{Emploi 1 : la tardiveté — « ne … que »}

\begin{propriete}[才 + moment : « ne … que (tard) »]
\textbf{Structure :} Sujet + (complément de temps) + 才 + Verbe (+ complément).\\[4pt]
才 placé avant le verbe indique que l'action se produit \textbf{plus tard que prévu, que souhaité ou que la norme}. En français, on le rend par « ne … que », « seulement à », « pas avant ».
\end{propriete}

\begin{exemplebox}[La tardiveté en contexte]
\zh{他昨天晚上十二点才睡觉。}\\
\emph{Tā zuótiān wǎnshang shí'èr diǎn cái shuìjiào.}\\
« Hier soir, il ne s'est couché qu'à minuit. » (C'est tard : on attendait un coucher plus tôt.)\\[4pt]
\zh{飞机晚点了，我们晚上十点才到杜阿拉。}\\
\emph{Fēijī wǎndiǎn le, wǒmen wǎnshang shí diǎn cái dào Dù'ālā.}\\
« L'avion avait du retard, nous ne sommes arrivés à Douala qu'à 22 h. »\\[4pt]
\zh{老师九点才来上课，同学们都等急了。}\\
\emph{Lǎoshī jiǔ diǎn cái lái shàngkè, tóngxuémen dōu děng jí le.}\\
« Le professeur n'est venu faire cours qu'à neuf heures, les élèves étaient impatients. »
\end{exemplebox}

\begin{attention}[才 et 了 ne font pas bon ménage]
Erreur fréquente des francophones : écrire \zh{他才来了。} pour « il n'est arrivé que tard ». \textbf{C'est faux !} Avec 才 au passé, on \textbf{n'utilise pas} le 了 d'accomplissement après le verbe. On dit : \zh{他昨天才来。} \emph{Tā zuótiān cái lái.} « Il n'est arrivé qu'hier. » Pourquoi ? Parce que 了 marque une action accomplie \emph{ponctuelle et close}, alors que 才 insiste sur le \emph{retard} de l'action : les deux marques entrent en conflit. (Exception apparente : \zh{我七点半就到了} avec 就 accepte 了 — c'est une asymétrie à mémoriser : \textbf{就 + 了 possible, 才 + 了 interdit} dans cet emploi.)
\end{attention}

\courssub{Emploi 2 : « seulement » devant une quantité — c'est peu}

\begin{propriete}[才 + quantité : « ne … que, à peine »]
\textbf{Structure :} Sujet + 才 + nombre + classificateur (+ nom).\\[4pt]
Devant une quantité, 才 signifie « seulement, à peine » et exprime que la quantité est \textbf{inférieure à l'attente} : c'est peu, c'est bon marché, c'est jeune. C'est l'inverse de 就 devant une quantité, qui peut souligner « déjà tant ».
\end{propriete}

\begin{exemplebox}[Quantités jugées petites]
\zh{这本书才五块钱。}\\
\emph{Zhè běn shū cái wǔ kuài qián.}\\
« Ce livre ne coûte que cinq yuans. » (C'est bon marché !)\\[4pt]
\zh{他才十五岁，就会说三种语言。}\\
\emph{Tā cái shíwǔ suì, jiù huì shuō sān zhǒng yǔyán.}\\
« Il n'a que quinze ans, et il sait déjà parler trois langues. » (Remarquez : 才 pour l'âge jugé jeune, 就 pour la compétence jugée précoce — les deux adverbes cohabitent !)\\[4pt]
\zh{从我家到学校才两公里，很近。}\\
\emph{Cóng wǒ jiā dào xuéxiào cái liǎng gōnglǐ, hěn jìn.}\\
« De chez moi à l'école, il n'y a que deux kilomètres, c'est très près. »
\end{exemplebox}

\courssub{Emploi 3 : la condition nécessaire — 只有……才……}

\begin{propriete}[只有 + condition + 才 + résultat]
\textbf{Structure :} 只有 + condition + 才 (+ 能/会) + résultat.\\[4pt]
Sens : « ce n'est qu'en … que … », « seulement si … alors … ». La condition est présentée comme \textbf{indispensable} : sans elle, pas de résultat. Cette structure est \textbf{très fréquente à l'examen} (thème et expression écrite). Ne confondez pas avec \zh{只要……就……} (« du moment que …, alors … » : condition \emph{suffisante}).
\end{propriete}

\begin{exemplebox}[La condition nécessaire]
\zh{只有努力，才能成功。}\\
\emph{Zhǐyǒu nǔlì, cái néng chénggōng.}\\
« Ce n'est qu'en travaillant dur qu'on peut réussir. »\\[4pt]
\zh{只有多练习，才能说好汉语。}\\
\emph{Zhǐyǒu duō liànxí, cái néng shuō hǎo Hànyǔ.}\\
« Ce n'est qu'en pratiquant beaucoup qu'on peut bien parler chinois. »\\[4pt]
\zh{只有身体健康，才能好好学习。}\\
\emph{Zhǐyǒu shēntǐ jiànkāng, cái néng hǎohao xuéxí.}\\
« Ce n'est qu'en étant en bonne santé qu'on peut bien étudier. »
\end{exemplebox}

\begin{attention}[只有……才…… contre 只要……就……]
\zh{只有} + 才 : condition \textbf{nécessaire} (« il faut … pour … »). \zh{只要} + 就 : condition \textbf{suffisante} (« il suffit de … pour … »). Comparez : \zh{只有努力，才能成功。} « Sans effort, pas de réussite » ; \zh{只要努力，就能成功。} « Dès que tu fais un effort, tu réussiras ». Le français « seulement » traduit les deux : méfiez-vous, et demandez-vous si la condition est \emph{indispensable} (只有…才…) ou \emph{suffisante} (只要…就…).
\end{attention}

\courssub{Emploi 4 : « pas avant que » — une action après une autre}

\begin{propriete}[Verbe₁ + 了 + … + 才 + Verbe₂]
\textbf{Structure :} Sujet + Verbe₁ + 了 + (complément) + 才 + Verbe₂.\\[4pt]
才 indique que la deuxième action n'a lieu \textbf{qu'après} l'accomplissement de la première : « ne … qu'après avoir … ». Ici, 了 accompagne le \emph{premier} verbe (action préalable accomplie), jamais le second.
\end{propriete}

\begin{exemplebox}[L'antériorité obligée]
\zh{他吃了饭才走。}\\
\emph{Tā chī le fàn cái zǒu.}\\
« Il n'est parti qu'après avoir mangé. »\\[4pt]
\zh{我做完了作业才看电视。}\\
\emph{Wǒ zuò wán le zuòyè cái kàn diànshì.}\\
« Je ne regarde la télévision qu'après avoir fini mes devoirs. »\\[4pt]
\zh{妈妈到了市场才想起没带钱。}\\
\emph{Māma dào le shìchǎng cái xiǎngqǐ méi dài qián.}\\
« Maman n'a réalisé qu'elle n'avait pas pris d'argent qu'une fois arrivée au marché. »
\end{exemplebox}

\courssub{Le miroir 就 / 才 : une image pour ne plus se tromper}

\begin{popfigure}
\begin{tikzpicture}[scale=1, every node/.style={font=\small}]
% Frise du haut : 就
\draw[thick, popblue, -{Stealth}] (0,1.6) -- (10,1.6);
\foreach \x/\h in {1.5/7h30, 5/8h, 8.5/9h} {
  \draw[popblue] (\x,1.45) -- (\x,1.75);
  \node[below, popdark] at (\x,1.45) {\h};
}
\draw[-{Stealth}, very thick, popgreen] (1.5,2.1) -- (1.5,1.75);
\node[above, popgreen, align=center] at (1.5,2.1) {Arrivée A};
\node[popdark, align=center] at (5,2.6) {\zh{他七点半就到了} — content \;\textbf{就} = plus tôt que prévu};
\node[popgreen, fill=popgreenL, circle, inner sep=2pt] at (1.5,2.5) {};
\node[popdark, align=center] at (1.5,2.8) {Content};
% Frise du bas : 才
\draw[thick, poporange, -{Stealth}] (0,-1.6) -- (10,-1.6);
\foreach \x/\h in {1.5/7h30, 5/8h, 8.5/9h} {
  \draw[poporange] (\x,-1.75) -- (\x,-1.45);
  \node[below, popdark] at (\x,-1.75) {\h};
}
\draw[-{Stealth}, very thick, poppink] (8.5,-2.3) -- (8.5,-1.75);
\node[below, poppink, align=center] at (8.5,-2.35) {Arrivée B};
\node[popdark, align=center] at (3.2,-2.7) {\zh{他九点才来} — mécontent \;\textbf{才} = plus tard que prévu};
\node[poppink, fill=poppinkL, circle, inner sep=2pt] at (8.5,-2.5) {};
\node[popdark, align=center] at (8.5,-2.8) {Mécontent};
% Rendez-vous commun
\node[popdark, align=center, fill=popgoldL, rounded corners, inner sep=2pt] at (5,0) {Rendez-vous : 8 h};
\end{tikzpicture}
\legende{Le miroir 就 / 才 : même rendez-vous à 8 h ; arrivée à 7 h 30 → 就 (tôt, satisfaction) ; arrivée à 9 h → 才 (tard, insatisfaction).}
\end{popfigure}

\begin{center}
\begin{tabularx}{\linewidth}{|X|X|X|X|}
\hline
\rowcolor{popblueL} Emploi de 才 & Structure & Exemple (caractères + pinyin) & Traduction \\
\hline
Tardiveté & S + temps + 才 + V & \zh{他十二点才睡觉。} \emph{Tā shí'èr diǎn cái shuìjiào.} & Il ne s'est couché qu'à minuit. \\
\hline
Quantité faible & S + 才 + nombre + classif. & \zh{这本书才五块钱。} \emph{Zhè běn shū cái wǔ kuài qián.} & Ce livre ne coûte que 5 yuans. \\
\hline
Condition nécessaire & 只有 + cond. + 才 + rés. & \zh{只有努力，才能成功。} \emph{Zhǐyǒu nǔlì, cái néng chénggōng.} & Ce n'est qu'en travaillant dur qu'on réussit. \\
\hline
Antériorité & V₁ + 了 + 才 + V₂ & \zh{他吃了饭才走。} \emph{Tā chī le fàn cái zǒu.} & Il n'est parti qu'après avoir mangé. \\
\hline
\end{tabularx}
\end{center}

\begin{methode}[Comment décider d'employer 才 ?]
\begin{enumerate}
\item Posez-vous la question : l'action est-elle \textbf{plus tard}, \textbf{plus difficile} ou la quantité \textbf{plus petite} que prévu ?
\item Si oui, placez \zh{才} juste avant le verbe (après le sujet et le complément de temps).
\item Au passé, \textbf{supprimez 了} après le verbe principal : \zh{他昨天才来。}
\item Si deux actions s'enchaînent (« seulement après … »), mettez 了 au premier verbe et 才 devant le second : \zh{吃了饭才走。}
\item Si la phrase exprime une condition indispensable, utilisez le cadre complet \zh{只有……才……}.
\item Test miroir : si le fait est \emph{plus tôt / plus facile / plus grand} que prévu, ce n'est pas 才, c'est \zh{就}.
\end{enumerate}
\end{methode}

\begin{exoresolu}[Traduisez en chinois]
Énoncé : Traduisez : « Hier, le bus de Yaoundé n'est arrivé qu'à dix heures du soir. » Puis : « Ce n'est qu'en écoutant attentivement qu'on peut comprendre le professeur. »
\tcblower
Solution détaillée. Phrase 1 : il s'agit d'une \textbf{tardiveté} (le bus est en retard, plus tard que prévu) → 才 avant le verbe, et \textbf{pas de 了} final malgré le passé. Traduction : \zh{昨天，雅温得的公共汽车晚上十点才到。} \emph{Zuótiān, Yǎwēndé de gōnggòng qìchē wǎnshang shí diǎn cái dào.} Phrase 2 : condition \textbf{nécessaire} → cadre 只有……才…… avec 能 : \zh{只有认真听，才能听懂老师的话。} \emph{Zhǐyǒu rènzhēn tīng, cái néng tīng dǒng lǎoshī de huà.} Piège évité : ni \zh{才到了}, ni \zh{只要……才……} (mélange des deux cadres).
\end{exoresolu}

\begin{attention}[Récapitulatif des pièges des francophones]
\begin{itemize}
\item \textbf{了 interdit} avec 才 de tardiveté : \zh{他才来了} est faux ; dites \zh{他才来}.
\item \textbf{Place} : 才 se met avant le verbe, jamais en fin de phrase comme le français « seulement » : on ne dit pas \zh{他睡觉才十二点}.
\item \textbf{Ne pas confondre} 才 (cái, « seulement/tard ») avec le nom \zh{才} dans \zh{人才} (réncái, « talent ») : même caractère, emplois différents.
\item \textbf{Français « ne … que »} : il traduit 才, mais le chinois n'a pas de négation ici — n'ajoutez pas \zh{不} : \zh{这本书才五块钱} et non \zh{这本书不才五块钱}.
\end{itemize}
\end{attention}

\begin{exercice}[À vous]
Transformez avec 才 : 1) Il est arrivé à 11 h (rendez-vous à 9 h). 2) Ce dictionnaire coûte seulement dix yuans. 3) Traduisez : « Ce n'est qu'en révisant chaque jour qu'on peut réussir à l'examen. » 4) Il n'est rentré à Buea qu'après la fin du match de football.
\end{exercice}

\begin{corrige}[Exercice « À vous »]
1) \zh{他十一点才来。} \emph{Tā shíyī diǎn cái lái.} (tardiveté, pas de 了). 2) \zh{这本词典才十块钱。} \emph{Zhè běn cídiǎn cái shí kuài qián.} (quantité faible). 3) \zh{只有每天复习，才能考试成功。} \emph{Zhǐyǒu měitiān fùxí, cái néng kǎoshì chénggōng.} (condition nécessaire). 4) \zh{足球比赛结束了，他才回布埃亚。} \emph{Zúqiú bǐsài jiéshù le, tā cái huí Bù'āiyà.} (antériorité : 了 au premier verbe).
\end{corrige}

\begin{aretenir}[L'essentiel sur 才]
\zh{才} exprime le \textbf{retard}, la \textbf{rareté} ou la \textbf{condition nécessaire} : « ne … que », « seulement », « pas avant ». Place : avant le verbe. Au passé tardif, \textbf{jamais de 了}. Cadre d'examen : \zh{只有……才……}. Miroir parfait de \zh{就} : même place, jugement inverse — 就 = plus tôt/plus que prévu, 才 = plus tard/moins que prévu.
\end{aretenir}

\coursec{还, 又, 再 : continuer et répéter}

Dans les sections précédentes, nous avons vu comment \zh{才} et \zh{就} situent une action dans le temps (tard ou tôt, avec difficulté ou avec facilité). Mais le français possède un autre petit mot redoutable : « encore ». Observez : « Il dort \emph{encore} », « Il est \emph{encore} venu hier », « Reviens \emph{encore} demain ». Un seul mot français, trois réalités différentes. Le chinois, lui, distingue rigoureusement ces trois situations avec trois adverbes : \zh{还}, \zh{又} et \zh{再}. C'est l'objet de cette section.

\courssub{Observation guidée : une journée de visites}

\begin{textebox}[Trois « encore » en une phrase]
他昨天又来了，今天还要来，说明天再来。\\
\emph{Tā zuótiān yòu lái le, jīntiān hái yào lái, shuō míngtiān zài lái.}\\
« Il est encore venu hier, il veut encore venir aujourd'hui, et il dit qu'il reviendra demain. »
\end{textebox}

Observons cette phrase construite autour du même verbe \zh{来} (venir) :

\begin{itemize}
\item \zh{昨天又来了} : hier, la visite s'est \cle{déjà reproduite} — c'est un fait accompli $\rightarrow$ \zh{又}.
\item \zh{今天还要来} : aujourd'hui, l'envie de venir \cle{continue} — pas de rupture $\rightarrow$ \zh{还}.
\item \zh{明天再来} : demain, la visite \cle{va se reproduire} — c'est un projet $\rightarrow$ \zh{再}.
\end{itemize}

La logique est limpide : le chinois regarde d'abord \textbf{où se situe l'action sur la ligne du temps}, puis choisit l'adverbe. Le français, lui, se contente d'un « encore » fourre-tout. C'est toute la difficulté — et toute la beauté — de cette leçon.

\courssub{还 (hái) : la continuation — « encore, toujours »}

\begin{propriete}[还 : continuité d'un état ou d'une action]
\zh{还} se place \textbf{avant le verbe ou l'adjectif} (et avant \zh{在} s'il y en a un) et indique qu'un état ou une action \cle{se poursuit sans interruption} jusqu'au moment où l'on parle.\\
\textbf{Schéma :} Sujet + 还 (+ 在) + Verbe / Adjectif\\
Il correspond au français « encore, toujours » (au sens de « sans changement »).
\end{propriete}

\begin{exemplebox}[还 = « encore / toujours »]
\zh{他还在睡觉。} \emph{Tā hái zài shuìjiào.} « Il dort encore. »\\
\zh{她还在雅温得上大学。} \emph{Tā hái zài Yǎwēndé shàng dàxué.} « Elle fait encore ses études à Yaoundé. »\\
\zh{现在十点，商店还开着。} \emph{Xiànzài shí diǎn, shāngdiàn hái kāizhe.} « Il est dix heures, le magasin est encore ouvert. »\\
\zh{爷爷八十岁了，身体还很好。} \emph{Yéye bāshí suì le, shēntǐ hái hěn hǎo.} « Grand-père a quatre-vingts ans, il est encore en bonne santé. »
\end{exemplebox}

\begin{propriete}[还 : « en plus, encore un » (quantité supplémentaire)]
Devant \zh{要}, \zh{想} ou un verbe d'action, \zh{还} peut aussi exprimer un \cle{ajout} : « encore un, de plus ».\\
\textbf{Schéma :} Sujet + 还 + 要/想 + (quantité) + Nom\\
\zh{我还要一杯茶。} \emph{Wǒ hái yào yì bēi chá.} « Je veux encore un thé (un de plus). »\\
\zh{他还想买一件衬衫。} \emph{Tā hái xiǎng mǎi yí jiàn chènshān.} « Il veut encore acheter une chemise (en plus). »
\end{propriete}

\begin{propriete}[还 + 没 + Verbe = « pas encore »]
Pour dire qu'une action \cle{n'a pas encore eu lieu} (mais qu'elle aura lieu), on utilise :\\
\textbf{Schéma :} Sujet + 还没 + Verbe (+ 呢)\\
\zh{我还没做完作业。} \emph{Wǒ hái méi zuòwán zuòyè.} « Je n'ai pas encore fini mes devoirs. »\\
\zh{火车还没来。} \emph{Huǒchē hái méi lái.} « Le train n'est pas encore arrivé. »\\
\zh{他还没回来呢。} \emph{Tā hái méi huílai ne.} « Il n'est pas encore rentré. »
\end{propriete}

\begin{attention}[还没 : une négation provisoire]
\zh{还没} implique que l'action \textbf{aura lieu plus tard} : \zh{我还没吃饭} = « je n'ai pas encore mangé (mais je vais manger) ». Ne le confondez pas avec une négation définitive. Et rappelez-vous : avec \zh{没}, on ne met \textbf{jamais} \zh{了} après le verbe.
\end{attention}

\courssub{又 (yòu) : la répétition constatée — « de nouveau » (passé)}

\begin{propriete}[又 : répétition déjà réalisée]
\zh{又} indique qu'une action s'est \cle{reproduite et que c'est un fait accompli}. Il regarde vers le \textbf{passé} et s'accompagne très souvent de \zh{了}.\\
\textbf{Schéma :} Sujet + 又 + Verbe + 了
\end{propriete}

\begin{exemplebox}[又 = « encore » au passé]
\zh{他又迟到了。} \emph{Tā yòu chídào le.} « Il est encore arrivé en retard. » (c'est fait, le professeur l'a vu !)\\
\zh{他又来了。} \emph{Tā yòu lái le.} « Le voilà revenu / il est encore venu. »\\
\zh{昨天又下雨了。} \emph{Zuótiān yòu xià yǔ le.} « Hier, il a encore plu. »\\
\zh{我又忘了一个汉字。} \emph{Wǒ yòu wàng le yí ge Hànzì.} « J'ai encore oublié un caractère. »
\end{exemplebox}

Remarquez la nuance : \zh{又} exprime souvent une légère \textbf{exaspération ou un constat} — « encore lui ! », « encore la pluie ! ». C'est exactement le ton du « encore » français au passé.

\begin{propriete}[又……又…… = « à la fois… et… »]
Devant deux \textbf{adjectifs ou verbes psychologiques}, la structure \zh{又} A \zh{又} B signifie « à la fois A et B ».\\
\textbf{Schéma :} Sujet + 又 + Adj$_1$ + 又 + Adj$_2$\\
\zh{这个菜又辣又好吃。} \emph{Zhège cài yòu là yòu hǎochī.} « Ce plat est à la fois piquant et délicieux. »\\
\zh{杜阿拉的夏天又热又湿。} \emph{Dùālā de xiàtiān yòu rè yòu shī.} « L'été à Douala est à la fois chaud et humide. »\\
\zh{她又聪明又努力。} \emph{Tā yòu cōngming yòu nǔlì.} « Elle est à la fois intelligente et travailleuse. »
\end{propriete}

\begin{attention}[又……又…… : pas de verbe d'action]
La structure \zh{又……又……} ne s'emploie qu'avec des \textbf{adjectifs} ou des \textbf{verbes d'état} (\zh{喜欢}, \zh{想}, \zh{怕}…), jamais avec deux verbes d'action. On ne dit pas \zh{他又吃饭又喝水} pour « il mange et il boit » ; on dira simplement \zh{他吃饭，也喝水} ou \zh{他一边吃饭，一边喝水} (il mange tout en buvant).
\end{attention}

\courssub{再 (zài) : la répétition projetée — « encore une fois » (futur)}

\begin{propriete}[再 : répétition future ou souhaitée]
\zh{再} indique qu'une action \cle{va se reproduire} : futur, souhait, ordre, promesse. Il regarde vers l'\textbf{avant} et ne prend \textbf{jamais} \zh{了}.\\
\textbf{Schéma :} Sujet + 再 + Verbe
\end{propriete}

\begin{exemplebox}[再 = « encore » au futur]
\zh{请再说一遍。} \emph{Qǐng zài shuō yí biàn.} « Veuillez répéter une fois, s'il vous plaît. »\\
\zh{我们明天再来。} \emph{Wǒmen míngtiān zài lái.} « Nous reviendrons demain. »\\
\zh{这个问题很难，请你再说一遍。} \emph{Zhège wèntí hěn nán, qǐng nǐ zài shuō yí biàn.} « Cette question est difficile, répétez encore une fois s'il vous plaît. »\\
\zh{我想再吃一块恩多雷。} \emph{Wǒ xiǎng zài chī yí kuài Ēnduōléi.} « Je veux encore manger un morceau de ndolé. »\\
\zh{等一下，我再想想。} \emph{Děng yíxià, wǒ zài xiǎngxiang.} « Attends un peu, je réfléchis encore. »
\end{exemplebox}

\begin{propriete}[再也不 = « ne… plus jamais »]
La négation de \zh{再} se fait avec \zh{不} (et non \zh{没}) : \zh{再也不 + Verbe} signifie « ne plus jamais ».\\
\zh{我再也不迟到了。} \emph{Wǒ zài yě bù chídào le.} « Je n'arriverai plus jamais en retard. » (promesse)\\
\zh{他再也不去那家饭馆了。} \emph{Tā zài yě bù qù nà jiā fànguǎn le.} « Il n'ira plus jamais dans ce restaurant. »
\end{propriete}

\begin{savaistu}[再见 : la preuve par le dictionnaire]
Le mot « au revoir » en chinois est \zh{再见} \emph{zàijiàn}, littéralement « encore — se voir », c'est-à-dire « \textbf{se revoir} ». Puisqu'on se reverra \emph{dans le futur}, c'est bien \zh{再} (et non \zh{又}) qui est employé. Retenez ce mot comme une démonstration vivante de la règle : \zh{再} regarde toujours vers l'avant !
\end{savaistu}

\courssub{La frise des temps : tout comprendre d'un coup d'œil}

\begin{popfigure}
\begin{center}
\begin{tikzpicture}[scale=1.0]
% Axe du temps
\draw[-{Stealth}, thick, popdark] (-4.5,0) -- (4.8,0);
\node[below, popdark] at (-3.6,0) {\textbf{PASSÉ}};
\node[below, popdark] at (0,0) {\textbf{PRÉSENT}};
\node[below, popdark] at (3.6,0) {\textbf{FUTUR}};
\fill[popdark] (0,0) circle (2pt);
% 还 : flèche continue traversant le présent
\draw[-{Stealth}, very thick, popblue] (-2.6,0.7) -- (1.8,0.7);
\node[above, popblue, align=center] at (-0.4,0.75) {\zh{还} : ça \textbf{continue}};
% 又 : boucle côté passé
\draw[-{Stealth}, very thick, poporange] (-3.4,0.2) .. controls (-3.4,1.6) and (-1.6,1.6) .. (-1.6,0.2);
\node[above, poporange, align=center] at (-2.5,1.55) {\zh{又} : déjà \textbf{refait}};
% 再 : boucle côté futur
\draw[-{Stealth}, very thick, popgreen] (1.6,0.2) .. controls (1.6,1.6) and (3.4,1.6) .. (3.4,0.2);
\node[above, popgreen, align=center] at (2.5,1.55) {\zh{再} : va être \textbf{refait}};
\end{tikzpicture}
\end{center}
\legende{Frise des temps : 还 traverse le présent (continuité), 又 boucle dans le passé (répétition constatée), 再 boucle dans le futur (répétition projetée).}
\end{popfigure}

\courssub{Tableau récapitulatif et comparaison avec le français}

\begin{center}
\begin{tabularx}{\linewidth}{|X|X|X|X|}
\hline
\rowcolor{popblueL} Adverbe & Valeur & Structure & Exemple + traduction \\
\hline
\zh{还} hái & continuité : « encore, toujours » & Sujet + 还 + V/Adj & \zh{他还在睡觉。} Il dort encore. \\
\hline
\zh{还} hái & ajout : « encore un, de plus » & Sujet + 还 + 要/想 + Nom & \zh{我还要一杯茶。} Je veux encore un thé. \\
\hline
\zh{还没} hái méi & « pas encore » & Sujet + 还没 + V & \zh{我还没做完作业。} Je n'ai pas encore fini mes devoirs. \\
\hline
\zh{又} yòu & répétition au passé : « de nouveau » & Sujet + 又 + V + 了 & \zh{他又迟到了。} Il est encore arrivé en retard. \\
\hline
\zh{又…又…} & « à la fois… et… » & 又 + Adj + 又 + Adj & \zh{这个菜又辣又好吃。} Ce plat est à la fois piquant et bon. \\
\hline
\zh{再} zài & répétition au futur : « encore une fois » & Sujet + 再 + V & \zh{我们明天再来。} Nous reviendrons demain. \\
\hline
\zh{再也不} & « ne… plus jamais » & Sujet + 再也不 + V & \zh{我再也不迟到了。} Je n'arriverai plus jamais en retard. \\
\hline
\end{tabularx}
\end{center}

\begin{attention}[Le piège n° 1 des francophones : « encore » n'est pas toujours 还]
L'erreur classique consiste à traduire mécaniquement « encore » par \zh{还}. Posez-vous d'abord la question du temps :
\begin{itemize}
\item « Il est \emph{encore} venu hier » = \zh{他昨天又来了} — jamais \zh{他还来了} (l'action est faite $\rightarrow$ \zh{又}).
\item « Reviens \emph{encore} demain » = \zh{明天再来} — jamais \zh{明天又来} (projet $\rightarrow$ \zh{再}).
\item « Il dort \emph{encore} » = \zh{他还在睡觉} — ici seulement \zh{还} (continuité).
\end{itemize}
Autres pièges : \zh{又} s'associe à \zh{了} (action accomplie), \zh{再} ne s'associe \textbf{jamais} à \zh{了} — on ne dit pas \zh{再来了}. Enfin, la négation de \zh{再} est \zh{不} (\zh{再也不}), celle de la continuité est \zh{没} (\zh{还没}).
\end{attention}

\begin{methode}[Comment choisir entre 还, 又 et 再 ?]
\begin{enumerate}
\item Repérez le mot français (« encore », « de nouveau », « toujours ») et le \textbf{temps} de la phrase.
\item L'action est-elle \textbf{déjà faite} (constat, souvent avec exaspération) ? $\rightarrow$ \zh{又} + Verbe + \zh{了}.
\item L'action va-t-elle \textbf{se faire} (futur, ordre, souhait, promesse) ? $\rightarrow$ \zh{再} + Verbe, sans \zh{了}.
\item L'action ou l'état \textbf{continue-t-il sans interruption} jusqu'à maintenant ? $\rightarrow$ \zh{还} + Verbe/Adjectif ; négation : \zh{还没}.
\item S'agit-il de deux qualités simultanées ? $\rightarrow$ \zh{又……又……}.
\end{enumerate}
\end{methode}

\begin{exoresolu}[Traduisez en choisissant le bon adverbe]
Énoncé : traduisez en chinois. 1) « Mon frère est encore arrivé en retard ce matin. » 2) « Nous reviendrons au marché de Douala demain. » 3) « Le match n'est pas encore fini. » 4) « Ce restaurant est à la fois bon et pas cher. »
\tcblower
Solution détaillée :
\begin{enumerate}
\item « Encore » + action accomplie ce matin $\rightarrow$ \zh{又} + \zh{了} : \zh{我弟弟今天早上又迟到了。} \emph{Wǒ dìdi jīntiān zǎoshang yòu chídào le.}
\item « Reviendrons demain » = répétition future $\rightarrow$ \zh{再} : \zh{我们明天再去杜阿拉的市场。} \emph{Wǒmen míngtiān zài qù Dùālā de shìchǎng.}
\item « Pas encore fini » $\rightarrow$ \zh{还没} : \zh{比赛还没结束。} \emph{Bǐsài hái méi jiéshù.}
\item Deux qualités simultanées $\rightarrow$ \zh{又……又……} : \zh{这家饭馆又好吃又便宜。} \emph{Zhè jiā fànguǎn yòu hǎochī yòu piányi.}
\end{enumerate}
\end{exoresolu}

\begin{exercice}[À vous de jouer]
Traduisez en chinois avec \zh{还}, \zh{又} ou \zh{再} :
\begin{enumerate}
\item Il pleut encore à Buea. (continuité)
\item Elle est encore venue me voir hier. (passé)
\item Répétez encore une fois, s'il vous plaît. (ordre)
\item Je n'ai pas encore fini mon thé. (négation)
\item Ce caractère est à la fois beau et difficile. (deux qualités)
\item Il a promis de ne plus jamais mentir. (négation future)
\end{enumerate}
\end{exercice}

\begin{corrige}[Exercice « À vous de jouer »]
\begin{enumerate}
\item \zh{布埃亚还在下雨。} \emph{Bù'āiyà hái zài xià yǔ.} — continuité $\rightarrow$ \zh{还}.
\item \zh{她昨天又来看我了。} \emph{Tā zuótiān yòu lái kàn wǒ le.} — passé $\rightarrow$ \zh{又} + \zh{了}.
\item \zh{请再说一遍。} \emph{Qǐng zài shuō yí biàn.} — ordre, futur $\rightarrow$ \zh{再}.
\item \zh{我还没喝完茶。} \emph{Wǒ hái méi hēwán chá.} — « pas encore » $\rightarrow$ \zh{还没}.
\item \zh{这个汉字又漂亮又难。} \emph{Zhège Hànzì yòu piàoliang yòu nán.} — deux qualités $\rightarrow$ \zh{又……又……}.
\item \zh{他保证再也不说谎了。} \emph{Tā bǎozhèng zài yě bù shuōhuǎng le.} — « ne… plus jamais » $\rightarrow$ \zh{再也不}.
\end{enumerate}
\end{corrige}

\begin{aretenir}[L'essentiel de la section]
\begin{itemize}
\item \zh{还} (hái) = continuité : « encore, toujours » ; devant \zh{要}/\zh{想}, il peut aussi marquer l'ajout (« encore un ») ; \zh{还没} = « pas encore ».
\item \zh{又} (yòu) = répétition \textbf{constatée} (passé), presque toujours avec \zh{了} ; \zh{又……又……} = « à la fois… et… » (adjectifs et verbes d'état uniquement).
\item \zh{再} (zài) = répétition \textbf{projetée} (futur, ordre, souhait), jamais avec \zh{了} ; négation : \zh{再也不} = « ne… plus jamais » ; \zh{再见} = « au revoir », littéralement « se revoir ».
\item Question réflexe : déjà fait ? $\rightarrow$ \zh{又}. Va se faire ? $\rightarrow$ \zh{再}. Continue ? $\rightarrow$ \zh{还}.
\end{itemize}
\end{aretenir}

\coursec{Comparaison avec le français et pièges des francophones}

Nous venons d'étudier \zh{还}, \zh{又} et \zh{再}, trois adverbes qui expriment la continuation et la répétition. Avant de passer à l'entraînement, il est indispensable de confronter ces structures — ainsi que \zh{就} et \zh{才} — à la langue française. En effet, c'est presque toujours la traduction mot à mot depuis le français qui provoque les erreurs des élèves camerounais : un même mot français correspond à plusieurs mots chinois, et l'ordre des mots n'est pas le même. Cette section est donc une « zone de déminage » : nous allons désamorcer un à un les pièges les plus fréquents.

\courssub{Un mot français, plusieurs mots chinois : le problème général}

Commençons par une observation. Traduisez mentalement en chinois les trois phrases françaises suivantes :

\begin{enumerate}
\item « Il n'a que huit ans. »
\item « Il est encore à la maison. »
\item « Dès qu'il arrive, il mange. »
\end{enumerate}

Chaque phrase contient un petit mot français (\emph{que}, \emph{encore}, \emph{dès que}) qui ne peut pas être traduit par un seul mot chinois universel. Le français est « pauvre » là où le chinois est « précis » : il faut d'abord analyser le sens exact, puis choisir la bonne structure chinoise.

\begin{center}
\begin{tabularx}{\linewidth}{|X|X|X|}
\hline
\rowcolor{popblueL} \textbf{Français} & \textbf{Chinois correct} & \textbf{Erreur à éviter} \\
\hline
seulement / ne\ldots que (tard, peu) & \zh{才} \emph{cái} & ne pas ajouter \zh{了} au passé \\
\hline
seulement / dès / tôt (tôt, facile) & \zh{就} \emph{jiù} & ne pas le placer après le verbe \\
\hline
encore (= toujours, pas encore) & \zh{还} \emph{hái} & ne pas lire \emph{huán} \\
\hline
encore (= de nouveau, passé constaté) & \zh{又} \emph{yòu} & ne pas l'utiliser pour un ordre \\
\hline
encore (= de nouveau, futur souhaité) & \zh{再} \emph{zài} & ne pas l'utiliser pour un fait passé \\
\hline
dès que & \zh{一……就……} \emph{yī\ldots jiù\ldots} & ne pas chercher un mot unique \\
\hline
\end{tabularx}
\end{center}

\begin{popfigure}
\begin{tikzpicture}[
  fr/.style={rectangle, rounded corners, draw=popblue, fill=popblueL, align=center, minimum width=2.6cm, minimum height=0.9cm, font=\small},
  zh/.style={rectangle, rounded corners, draw=popgreen, fill=popgreenL, align=center, minimum width=2.6cm, minimum height=0.9cm, font=\small},
  err/.style={rectangle, rounded corners, draw=poppink, fill=poppinkL, align=center, minimum width=3.4cm, minimum height=0.9cm, font=\small},
  fl/.style={-{Stealth[length=2.5mm]}, thick, popdark}
]
\node[fr, align=center] (f1) at (0,4.4){seulement\\(tard / peu)};
\node[zh] (z1) at (4.2,4.4) {才 \emph{cái}};
\node[err] (e1) at (8.6,4.4) {× 他昨天才来了。};
\node[fr, align=center] (f2) at (0,3.2){seulement\\(tôt / déjà)};
\node[zh] (z2) at (4.2,3.2) {就 \emph{jiù}};
\node[err, align=center] (e2) at (8.6,3.2){× adverbe après\\le verbe};
\node[fr, align=center] (f3) at (0,2.0){encore\\(toujours)};
\node[zh] (z3) at (4.2,2.0) {还 \emph{hái}};
\node[err] (e3) at (8.6,2.0) {× lire \emph{huán}};
\node[fr, align=center] (f4) at (0,0.8){encore\\(passé répété)};
\node[zh] (z4) at (4.2,0.8) {又 \emph{yòu}};
\node[err] (e4) at (8.6,0.8) {× 又做一次};
\node[fr, align=center] (f5) at (0,-0.4){encore\\(futur)};
\node[zh] (z5) at (4.2,-0.4) {再 \emph{zài}};
\node[err] (e5) at (8.6,-0.4) {× 他昨天再来};
\node[fr] (f6) at (0,-1.6) {dès que};
\node[zh] (z6) at (4.2,-1.6) {一……就……};
\node[err, align=center] (e6) at (8.6,-1.6){× traduire par\\un seul mot};
\draw[fl] (f1) -- (z1); \draw[fl] (z1) -- (e1);
\draw[fl] (f2) -- (z2); \draw[fl] (z2) -- (e2);
\draw[fl] (f3) -- (z3); \draw[fl] (z3) -- (e3);
\draw[fl] (f4) -- (z4); \draw[fl] (z4) -- (e4);
\draw[fl] (f5) -- (z5); \draw[fl] (z5) -- (e5);
\draw[fl] (f6) -- (z6); \draw[fl] (z6) -- (e6);
\end{tikzpicture}
\legende{Correspondances français -- chinois : pour chaque mot français, la bonne structure chinoise et l'erreur typique du francophone.}
\end{popfigure}

\courssub{« Seulement / ne\ldots que » : 才 ou 就 ?}

Le français utilise le même mot « seulement » (ou la tournure « ne\ldots que ») dans deux situations que le chinois distingue nettement.

\begin{exemplebox}[Observer avant de conclure]
\zh{他才八岁。} \emph{Tā cái bā suì.} « Il n'a que huit ans. » (c'est peu, c'est surprenant : il est si jeune !)\\
\zh{他就八岁，已经会说三种语言了。} \emph{Tā jiù bā suì, yǐjīng huì shuō sān zhǒng yǔyán le.} « Il n'a que huit ans, et il sait déjà parler trois langues. » (huit ans suffisent : c'est tôt, c'est rapide.)
\end{exemplebox}

\begin{propriete}[才 = « seulement » qui insiste sur le peu ou le tard ; 就 = « seulement » qui insiste sur le tôt ou le facile]
\begin{itemize}
\item \cle{才} + nombre/adjectif : la quantité est jugée \textbf{faible}, le moment \textbf{tardif}. \zh{这部电影才一个小时。} \emph{Zhè bù diànyǐng cái yí ge xiǎoshí.} « Ce film ne dure qu'une heure. »
\item \cle{就} + nombre/verbe : la quantité est jugée \textbf{suffisante}, le moment \textbf{précoce}. \zh{他六点就起床了。} \emph{Tā liù diǎn jiù qǐchuáng le.} « Il s'est levé dès six heures. »
\end{itemize}
En français, les deux se traduisent souvent par « seulement / ne\ldots que / dès » ; c'est le \textbf{jugement du locuteur} (trop tard ? étonnamment tôt ?) qui décide du choix en chinois.
\end{propriete}

\courssub{« Encore » : trois mots chinois pour un seul mot français}

C'est le piège majeur de cette leçon. Le français « encore » recouvre trois réalités différentes :

\begin{exemplebox}[Trois « encore », trois chinois]
\zh{他还在睡觉。} \emph{Tā hái zài shuìjiào.} « Il dort encore. » (l'action \textbf{continue})\\
\zh{他又迟到了。} \emph{Tā yòu chídào le.} « Il est encore arrivé en retard. » (cela s'est \textbf{déjà reproduit}, constat)\\
\zh{请再说一遍。} \emph{Qǐng zài shuō yí biàn.} « Répétez encore une fois, s'il vous plaît. » (cela va se reproduire, \textbf{souhait})
\end{exemplebox}

\begin{methode}[Comment traduire « encore » en trois questions]
\begin{enumerate}
\item L'action \textbf{continue-t-elle} (ou n'a-t-elle pas encore eu lieu) ? $\rightarrow$ \cle{还}. Ex. : \zh{我还没吃饭。} « Je n'ai pas encore mangé. »
\item L'action s'est-elle \textbf{déjà répétée} (constat, souvent avec \zh{了}) ? $\rightarrow$ \cle{又}. Ex. : \zh{他又来了。} « Il est encore venu. »
\item L'action va-t-elle \textbf{se répéter} (futur, ordre, souhait) ? $\rightarrow$ \cle{再}. Ex. : \zh{明天再来。} « Reviens demain. »
\end{enumerate}
\end{methode}

\begin{attention}[还没 = « pas encore », pas « pas »]
\zh{我还没吃饭。} \emph{Wǒ hái méi chīfàn.} ne veut pas dire « je n'ai pas mangé » mais « je n'ai pas \textbf{encore} mangé » (je vais manger bientôt). Pour dire simplement « je n'ai pas mangé », on dit \zh{我没吃饭。} La présence de \zh{还} ajoute toujours l'idée d'attente : la chose est prévue, elle n'a simplement pas eu lieu \emph{jusqu'à maintenant}. Comparez : \zh{他还没来。} « Il n'est pas encore arrivé » / \zh{他没来。} « Il n'est pas venu » (c'est fini, il n'est pas venu du tout).
\end{attention}

\courssub{La place des adverbes : ne jamais calquer le français}

En français, « seulement » et « encore » flottent dans la phrase : « Il a \emph{seulement} mangé du riz », « Il a mangé \emph{seulement} du riz », « \emph{Encore} il a mangé du riz ». En chinois, cette liberté n'existe pas.

\begin{propriete}[Ordre des mots : l'adverbe précède toujours le verbe]
\begin{center}
\textbf{Sujet + (complément de temps) + 就 / 才 / 还 / 又 / 再 + Verbe + (complément)}
\end{center}
L'adverbe ne se déplace \textbf{jamais} après le verbe et ne se place jamais avant le sujet.\\
\zh{我明天才走。} \emph{Wǒ míngtiān cái zǒu.} « Je ne pars que demain. »\\
Faute typique du francophone : × \zh{我走才明天} ou × \zh{才我明天走} — calques de l'ordre français.
\end{propriete}

\courssub{Les quatre pièges du Bac : 了, 又/再, la lecture de 还, « dès que »}

\begin{attention}[Piège 1 : 了 est interdit avec 才 au passé, naturel avec 就]
Avec \zh{才}, l'action passée se présente comme un simple repère : on ne met \textbf{jamais} \zh{了}.\\
Erreur : × \zh{他昨天才来了。} $\rightarrow$ Correct : \zh{他昨天才来。} \emph{Tā zuótiān cái lái.} « Il n'est arrivé qu'hier. »\\
Avec \zh{就}, au contraire, \zh{了} est naturel, voire obligatoire pour une action accomplie : \zh{他八点钟就来了。} \emph{Tā bā diǎn zhōng jiù lái le.} « Il est arrivé dès huit heures. »
\end{attention}

\begin{attention}[Piège 2 : « refais-le » = 再做一次, jamais 又做一次]
Dans les consignes, les ordres et les souhaits, la répétition est \textbf{à venir} : il faut \zh{再}. \zh{请再做一次。} \emph{Qǐng zài zuò yí cì.} « Refais-le une fois, s'il te plaît. » La phrase × \zh{请又做一次} est impossible, car \zh{又} constate une répétition \textbf{déjà accomplie} : \zh{他又做了一次。} « Il l'a refait une fois. »
\end{attention}

\begin{attention}[Piège 3 : 还 se lit hái (adverbe), pas huán]
Le caractère \zh{还} a deux lectures : \emph{hái} quand il est adverbe (« encore, toujours ») et \emph{huán} quand il est verbe (« rendre, retourner »), comme dans \zh{还书} \emph{huán shū} « rendre un livre ». À la lecture à voix haute, vérifiez la fonction du mot avant de prononcer : \zh{我还要还书。} \emph{Wǒ hái yào huán shū.} « Je dois encore rendre les livres. » — les deux lectures coexistent dans la même phrase !
\end{attention}

\begin{attention}[Piège 4 : « dès que » exige la paire 一……就……]
Le français « dès que » ne se traduit pas par un mot unique : le chinois utilise une \textbf{structure en deux morceaux} : \cle{一} + verbe 1 + \cle{就} + verbe 2.\\
\zh{他一到家就吃饭。} \emph{Tā yí dào jiā jiù chīfàn.} « Dès qu'il arrive à la maison, il mange. »\\
\zh{我一看见他就高兴。} \emph{Wǒ yí kànjian tā jiù gāoxìng.} « Dès que je le vois, je suis content. »\\
Erreur fréquente : n'utiliser que \zh{就} en oubliant le \zh{一} initial, ce qui affaiblit ou change le sens.
\end{attention}

\begin{savaistu}[Au Bac, on vous tend le piège 又 / 再]
Les questions de traduction (thème et version) ciblent très souvent la confusion entre \zh{又} et \zh{再}. Réflexe de survie : avant de choisir, \textbf{relisez le temps de la phrase}. Si l'action répétée est déjà faite (souvent signalée par \zh{了}, \zh{昨天}, \zh{刚才}), choisissez \zh{又} ; si elle est à venir (avec \zh{明天}, \zh{以后}, un impératif), choisissez \zh{再}. Une seconde de vérification peut sauver un point.
\end{savaistu}

\courssub{Exercices résolus}

\begin{exoresolu}[Traduction : choisir la bonne structure]
Énoncé. Traduisez en chinois : 1) « Mon petit frère n'a que six ans. » 2) « Il dort encore. » 3) « Le professeur est encore arrivé en retard hier. » 4) « Revenez demain, s'il vous plaît. » 5) « Dès qu'il finit ses devoirs, il joue au football. »
\tcblower
Solution détaillée.\\
1) « Seulement » + jugement « c'est peu » $\rightarrow$ \zh{才} : \zh{我弟弟才六岁。} \emph{Wǒ dìdi cái liù suì.}\\
2) L'action continue $\rightarrow$ \zh{还} : \zh{他还在睡觉。} \emph{Tā hái zài shuìjiào.}\\
3) Répétition déjà constatée (\zh{昨天} + \zh{了}) $\rightarrow$ \zh{又} : \zh{老师昨天又迟到了。} \emph{Lǎoshī zuótiān yòu chídào le.}\\
4) Répétition à venir (ordre poli) $\rightarrow$ \zh{再} : \zh{请明天再来。} \emph{Qǐng míngtiān zài lái.}\\
5) « Dès que » $\rightarrow$ paire \zh{一……就……} : \zh{他一做完作业就踢足球。} \emph{Tā yí zuòwán zuòyè jiù tī zúqiú.}
\end{exoresolu}

\begin{exoresolu}[Correction de phrases fautives]
Énoncé. Chaque phrase contient une erreur typique de francophone. Corrigez-la et justifiez. 1) × \zh{他昨天才来了。} 2) × \zh{请又做一次。} 3) × \zh{我吃才了饭。} 4) × \zh{他就到家吃饭。} (pour « dès qu'il arrive à la maison, il mange »)
\tcblower
Solution détaillée.\\
1) \zh{他昨天才来。} — avec \zh{才}, pas de \zh{了} au passé.\\
2) \zh{请再做一次。} — une consigne porte sur une action à venir : \zh{再}, pas \zh{又}.\\
3) \zh{我才吃了饭。} — l'adverbe \zh{才} se place avant le verbe, jamais après.\\
4) \zh{他一到家就吃饭。} — « dès que » exige la paire complète \zh{一……就……} ; le \zh{一} initial ne peut pas être omis dans ce sens.
\end{exoresolu}

\begin{aretenir}[L'essentiel de la section]
\begin{itemize}
\item « Seulement » : \zh{才} si c'est peu/tard, \zh{就} si c'est tôt/facile.
\item « Encore » : \zh{还} (continue), \zh{又} (déjà répété), \zh{再} (va se répéter).
\item L'adverbe se place toujours \textbf{avant le verbe}, après le sujet et le complément de temps.
\item Pas de \zh{了} avec \zh{才} ; \zh{了} naturel avec \zh{就} au passé.
\item « Dès que » = \zh{一……就……} ; « pas encore » = \zh{还没}.
\end{itemize}
\end{aretenir}

\coursec{Entraînement : application et transformation}

Nous avons vu dans la section précédente les pièges que le français tend aux apprenants : traduire mécaniquement « déjà » par \zh{就}, « seulement » par \zh{才}, ou confondre « encore » avec \zh{还}, \zh{又} et \zh{再}. Il est temps de vérifier que ces structures sont bien acquises. Cette dernière section est entièrement consacrée à l'entraînement : exercices à trous, transformations, traductions (thème et version), correction d'erreurs et production écrite guidée. Chaque exercice est suivi d'un corrigé détaillé : travaille d'abord seul, puis compare avec la solution.

\courssub{La méthode générale : l'organigramme de décision}

Avant de te lancer dans les exercices, retiens une démarche simple et systématique. Face à une phrase à compléter ou à transformer, ne réponds jamais « au feeling » : interroge la phrase.

\begin{methode}[Méthode pour les exercices à trous]
\begin{enumerate}
\item \textbf{Repérer l'indice de temps.} Souligne d'abord l'indice temporel : \zh{昨天} (hier), \zh{明天} (demain), \zh{已经} (déjà), une heure précise. C'est lui qui décide presque toujours de la réponse.
\item \textbf{Poser la bonne question.} L'action continue-t-elle encore ? S'est-elle déjà répétée ? Va-t-elle se répéter ? Le locuteur juge-t-il l'heure tôt ou tard, la quantité grande ou petite ?
\item \textbf{Appliquer l'organigramme} ci-dessous pour choisir entre \zh{还}, \zh{又}, \zh{再}, \zh{就} et \zh{才}.
\item \textbf{Vérifier la place du mot.} Ces cinq adverbes se placent toujours \emph{avant le verbe} (et après le sujet et l'indice de temps). Relis la phrase complète à voix haute.
\end{enumerate}
\end{methode}

\begin{popfigure}
\begin{tikzpicture}[
  node distance=0.55cm and 0.9cm,
  quest/.style={rectangle, rounded corners, draw=popblue, fill=popblueL, align=center, font=\small, inner sep=4pt},
  rep/.style={rectangle, rounded corners, draw=poporange, fill=poporangeL, align=center, font=\small\bfseries, inner sep=4pt},
  lab/.style={font=\scriptsize\itshape, text=popdark},
  arr/.style={-{Stealth[length=2.5mm]}, thick, popdark}
]
\node[quest, align=center] (q1){L'action continue-t-elle\\encore (pas finie) ?};
\node[rep, right=of q1, align=center] (r1){还 hái\\(encore, toujours)};
\node[quest, below=of q1, align=center] (q2){S'est-elle déjà répétée\\(passé, constat) ?};
\node[rep, right=of q2, align=center] (r2){又 yòu\\(de nouveau, passé)};
\node[quest, below=of q2, align=center] (q3){Va-t-elle se répéter\\(futur, souhait) ?};
\node[rep, right=of q3, align=center] (r3){再 zài\\(de nouveau, futur)};
\node[quest, below=of q3, align=center] (q4){Le locuteur juge-t-il\\tôt / facile / rapide ?};
\node[rep, right=of q4, align=center] (r4){就 jiù\\(tôt, dès, facilité)};
\node[rep, below=of q4, align=center] (r5){才 cái\\(tard, seulement, difficulté)};
\draw[arr] (q1) -- node[lab, above]{oui} (r1);
\draw[arr] (q1) -- node[lab, right]{non} (q2);
\draw[arr] (q2) -- node[lab, above]{oui} (r2);
\draw[arr] (q2) -- node[lab, right]{non} (q3);
\draw[arr] (q3) -- node[lab, above]{oui} (r3);
\draw[arr] (q3) -- node[lab, right]{non} (q4);
\draw[arr] (q4) -- node[lab, above]{oui} (r4);
\draw[arr] (q4) -- node[lab, right]{non : tard / difficile} (r5);
\end{tikzpicture}
\legende{Organigramme de décision : choisir entre 还, 又, 再, 就 et 才.}
\end{popfigure}

\begin{center}
\begin{tabularx}{\linewidth}{|X|X|X|X|}
\hline
\rowcolor{popblueL} Mot & Valeur & Temps typique & Équivalent français \\
\hline
就 jiù & tôt, vite, facile, dès & passé ou futur & dès, déjà, il suffit de \\
\hline
才 cái & tard, difficile, seulement & passé surtout & ne... que, seulement à \\
\hline
还 hái & continuation d'un état & présent / passé & encore, toujours \\
\hline
又 yòu & répétition constatée & passé & de nouveau, encore une fois \\
\hline
再 zài & répétition projetée & futur & encore, à nouveau, re- \\
\hline
\end{tabularx}
\end{center}

\courssub{Exercice 1 : à trous avec 就 ou 才}

\begin{exoresolu}[Exercice 1 — 就 ou 才 ?]
Énoncé. Complète avec \zh{就} ou \zh{才} :
\begin{enumerate}
\item 他早上五点\_\_\_\_起床了。
\item 火车晚上十一点\_\_\_\_到。
\end{enumerate}
\tcblower
Solution détaillée.
\begin{enumerate}
\item \zh{他早上五点就起床了。} \emph{Tā zǎoshang wǔ diǎn jiù qǐchuáng le.} « Il s'est levé dès cinq heures du matin. » Indice : cinq heures du matin est une heure jugée \cle{tôt} ; le locuteur souligne la précocité, donc \zh{就}. Note la présence de \zh{了} final, compatible avec \zh{就}.
\item \zh{火车晚上十一点才到。} \emph{Huǒchē wǎnshang shíyī diǎn cái dào.} « Le train n'arrive qu'à onze heures du soir. » Indice : onze heures du soir est jugé \cle{tard} pour une arrivée ; le locuteur exprime l'attente, donc \zh{才}. Pas de \zh{了} final.
\end{enumerate}
\end{exoresolu}

\begin{exemplebox}[Exemples traduits supplémentaires]
\zh{他早上五点就起床了。} \emph{Tā zǎoshang wǔ diǎn jiù qǐchuáng le.} « Il s'est levé dès cinq heures du matin. »\\
\zh{我学了三个月就会说一点儿汉语了。} \emph{Wǒ xué le sān ge yuè jiù huì shuō yìdiǎnr Hànyǔ le.} « Après seulement trois mois d'étude, je savais déjà parler un peu chinois. » (facilité, rapidité)\\
\zh{他学了三年才会说一点儿汉语。} \emph{Tā xué le sān nián cái huì shuō yìdiǎnr Hànyǔ.} « Il a mis trois ans avant de savoir parler un peu chinois. » (difficulté, lenteur)
\end{exemplebox}

\courssub{Exercice 2 : à trous avec 还, 又 ou 再}

\begin{exoresolu}[Exercice 2 — 还, 又 ou 再 ?]
Énoncé. Complète avec \zh{还}, \zh{又} ou \zh{再} :
\begin{enumerate}
\item 他\_\_\_\_没回来。
\item 他\_\_\_\_迟到了。
\item 我们明天\_\_\_\_见。
\end{enumerate}
\tcblower
Solution détaillée.
\begin{enumerate}
\item \zh{他还没回来。} \emph{Tā hái méi huílái.} « Il n'est pas encore rentré. » L'état « pas rentré » \cle{continue} au moment de la parole : \zh{还} + négation \zh{没}.
\item \zh{他又迟到了。} \emph{Tā yòu chídào le.} « Il est encore arrivé en retard. » La répétition est \cle{constatée} (elle a déjà eu lieu, marque \zh{了}) : \zh{又}.
\item \zh{我们明天再见。} \emph{Wǒmen míngtiān zài jiàn.} « Nous nous reverrons demain. » La répétition est \cle{projetée} dans le futur (\zh{明天}) : \zh{再}.
\end{enumerate}
\end{exoresolu}

\begin{attention}[明天又见 : l'erreur classique du Bac]
Dans \zh{明天再见} \emph{míngtiān zài jiàn} (« à demain, au revoir »), \zh{再} est \textbf{obligatoire} car l'action de se revoir est future. Écrire \zh{明天又见} est une erreur classique au baccalauréat : \zh{又} ne s'emploie que pour une répétition déjà réalisée. Astuce : \zh{又} regarde en arrière, \zh{再} regarde en avant.
\end{attention}

\courssub{Exercice 3 : transformation — du fait au jugement}

\begin{propriete}[Ajouter 就 ou 才 : le fait ne change pas, le jugement oui]
Dans la transformation, on garde les mêmes faits (même sujet, même heure, même verbe) et on insère \zh{就} ou \zh{才} avant le verbe. L'événement reste identique ; seul change le \cle{jugement du locuteur} : \zh{就} présente l'action comme tôt, rapide ou facile ; \zh{才} la présente comme tardive ou difficile. Attention : avec \zh{就} + action passée, on ajoute souvent \zh{了} en fin de phrase ; avec \zh{才}, on n'ajoute jamais \zh{了} après le verbe.
\end{propriete}

\begin{exoresolu}[Exercice 3 — Transformer une phrase neutre]
Énoncé. Phrase neutre : \zh{他七点来。} Transforme-la pour exprimer (a) qu'il est arrivé tôt, (b) qu'il est arrivé tard.
\tcblower
Solution détaillée.
\begin{enumerate}
\item[(a)] \zh{他七点就来了。} \emph{Tā qī diǎn jiù lái le.} « Il est arrivé dès sept heures. » (jugement : tôt, plus tôt que prévu)
\item[(b)] \zh{他七点才来。} \emph{Tā qī diǎn cái lái.} « Il n'est arrivé qu'à sept heures. » (jugement : tard, on l'a attendu)
\end{enumerate}
Même heure, même action : seule l'attitude du locuteur change. C'est exactement ce que le français rend par « dès » contre « ne... que ».
\end{exoresolu}

\courssub{Exercice 4 : traduction (thème)}

\begin{exercice}[Exercice 4 — Traduire en chinois]
\begin{enumerate}
\item « Il n'est rentré qu'à minuit. »
\item « Dès qu'il pleut, je reste à la maison. »
\item « Répète encore une fois. »
\end{enumerate}
\end{exercice}

\begin{corrige}[Exercice 4]
\begin{enumerate}
\item \zh{他半夜十二点才回家。} \emph{Tā bànyè shí'èr diǎn cái huí jiā.} Minuit est jugé tard : \zh{才}. Pas de \zh{了} final avec \zh{才}.
\item \zh{一下雨，我就待在家里。} \emph{Yí xià yǔ, wǒ jiù dāi zài jiā lǐ.} Structure \zh{一……就……} : « dès que... aussitôt... ». \zh{就} introduit la conséquence immédiate. Attention à ne pas confondre ce \zh{再} de \zh{待在家里} (« à la maison », préposition \zh{在}) avec l'adverbe \zh{再} : homophones, emplois différents.
\item \zh{请再说一遍。} \emph{Qǐng zài shuō yí biàn.} La répétition est souhaitée, donc future : \zh{再}, jamais \zh{又} dans un ordre.
\end{enumerate}
\end{corrige}

\courssub{Exercice 5 : traduction (version)}

\begin{exercice}[Exercice 5 — Traduire en français]
\begin{enumerate}
\item \zh{我只有好好学习，才能考上大学。}
\item \zh{她又生病了，今天不能来上课。}
\end{enumerate}
\end{exercice}

\begin{corrige}[Exercice 5]
\begin{enumerate}
\item \zh{我只有好好学习，才能考上大学。} \emph{Wǒ zhǐyǒu hǎohao xuéxí, cái néng kǎoshang dàxué.} « Ce n'est qu'en étudiant sérieusement que je pourrai entrer à l'université. » Structure \zh{只有……才……} : la condition exclusive.
\item \zh{她又生病了，今天不能来上课。} \emph{Tā yòu shēngbìng le, jīntiān bù néng lái shàngkè.} « Elle est encore tombée malade, elle ne peut pas venir en cours aujourd'hui. » \zh{又} + \zh{了} : répétition constatée, avec une nuance d'agacement.
\end{enumerate}
\end{corrige}

\begin{exemplebox}[La condition exclusive 只有……才……]
\zh{我只有好好学习，才能考上大学。} \emph{Wǒ zhǐyǒu hǎohao xuéxí, cái néng kǎoshang dàxué.} « Ce n'est qu'en étudiant sérieusement que je pourrai entrer à l'université. »\\
Autre exemple : \zh{只有多练习，才能说好汉语。} \emph{Zhǐyǒu duō liànxí, cái néng shuō hǎo Hànyǔ.} « Ce n'est qu'en s'entraînant beaucoup qu'on peut bien parler chinois. »
\end{exemplebox}

\courssub{Exercice 6 : correction d'erreurs}

\begin{exoresolu}[Exercice 6 — Retrouve et corrige l'erreur]
Énoncé. Chaque phrase contient une erreur. Retrouve-la et corrige-la.
\begin{enumerate}
\item \zh{他才来了。}
\item \zh{请又说一遍。}
\item \zh{我还昨天去了。}
\end{enumerate}
\tcblower
Solution détaillée.
\begin{enumerate}
\item \zh{他才来了。} → \zh{他才来。} \emph{Tā cái lái.} « Il n'est arrivé que tard / il vient seulement d'arriver. » Erreur : \zh{才} est incompatible avec \zh{了} final (qui marque l'accomplissement, alors que \zh{才} insiste sur le retard). On garde \zh{了} avec \zh{就} : \zh{他就来了。} « il est vite arrivé », mais jamais avec \zh{才}.
\item \zh{请又说一遍。} → \zh{请再说一遍。} \emph{Qǐng zài shuō yí biàn.} « Veuillez répéter une fois. » Erreur : une demande porte sur une action \cle{future}, donc \zh{再}. \zh{又} décrirait une répétition déjà faite : \zh{他又说了一遍。} « il l'a encore dit une fois ».
\item \zh{我还昨天去了。} → \zh{我昨天还去了。} \emph{Wǒ zuótiān hái qù le.} « Hier, j'y suis quand même allé. » Erreur : l'indice de temps \zh{昨天} se place juste après le sujet, \emph{avant} l'adverbe \zh{还}. Ordre : Sujet + temps + \zh{还} + verbe.
\end{enumerate}
\end{exoresolu}

\courssub{Exercice 7 : production guidée — ma journée d'hier}

\begin{experience}[Exercice 7 — Rédaction guidée]
Rédige cinq phrases sur ta journée d'hier en utilisant \textbf{au moins une fois chacun} : \zh{就}, \zh{才}, \zh{还}, \zh{又}, \zh{再}. Aide-toi du plan suivant : lever → départ → cours → retour → projet pour demain. Souligne chaque adverbe utilisé et vérifie sa place (après le sujet et l'indice de temps, avant le verbe). Tu peux situer ta journée à Yaoundé ou à Douala pour enrichir le contexte.
\end{experience}

\begin{exemplebox}[Modèle de production (corrigé possible)]
\zh{昨天我早上五点半就起床了。} \emph{Zuótiān wǒ zǎoshang wǔ diǎn bàn jiù qǐchuáng le.} « Hier, je me suis levé dès cinq heures et demie. » (\zh{就} : tôt)\\
\zh{我七点才到学校，因为路上堵车了。} \emph{Wǒ qī diǎn cái dào xuéxiào, yīnwèi lù shang dǔchē le.} « Je ne suis arrivé à l'école qu'à sept heures, parce qu'il y avait des embouteillages. » (\zh{才} : tard)\\
\zh{上课的时候，我还没做完作业。} \emph{Shàngkè de shíhou, wǒ hái méi zuòwán zuòyè.} « Au moment du cours, je n'avais pas encore fini mes devoirs. » (\zh{还} : continuation)\\
\zh{下午我又迟到了五分钟。} \emph{Xiàwǔ wǒ yòu chídào le wǔ fēnzhōng.} « L'après-midi, je suis encore arrivé en retard de cinq minutes. » (\zh{又} : répétition constatée)\\
\zh{明天我要再早点儿出门。} \emph{Míngtiān wǒ yào zài zǎo diǎnr chūmén.} « Demain, je partirai encore un peu plus tôt. » (\zh{再} : projet futur)
\end{exemplebox}

\begin{aretenir}[Bilan de l'entraînement]
\begin{itemize}
\item \zh{就} = tôt, vite, facile, « dès » ; compatible avec \zh{了} final ; structure \zh{一……就……} « dès que... ».
\item \zh{才} = tard, difficile, « ne... que » ; jamais de \zh{了} final ; structure clé \zh{只有……才……}.
\item \zh{还} = l'action ou l'état continue : « encore, toujours » ; se place après l'indice de temps.
\item \zh{又} = répétition déjà réalisée (passé) ; \zh{再} = répétition à venir (futur, ordre, souhait).
\item Méthode : souligner l'indice de temps, poser la bonne question, appliquer l'organigramme, vérifier la place de l'adverbe.
\end{itemize}
\end{aretenir}

\newpage

\coursec{Activité d'intégration}

\coursec{Leçon 6 — Grammaire : 就 / 才 ; 还 / 又 / 再}

\courssub{Observation guidée : 就 et 才 en contexte}

Avant d'énoncer les règles, observons comment \cle{就} et \cle{才} se comportent face à une même référence temporelle. Lisez les deux mini-dialogues ci-dessous, notez la position de l'adverbe (après l'heure, avant le verbe) et la nuance de sens (précocité vs retard).

\begin{textebox}[Dialogue A — Le bus est parti tôt]
— 你几点走的？ \emph{Nǐ jǐ diǎn zǒu de?} « À quelle heure es-tu parti ? »\\
— 我七点\cle{就}走了。 \emph{Wǒ qī diǎn \cle{jiù} zǒu le.} « Je suis parti \textbf{dès} sept heures (plus tôt que prévu). »\\
— 真早！ \emph{Zhēn zǎo!} « Vraiment tôt ! »
\end{textebox}

\begin{textebox}[Dialogue B — Le bus est arrivé tard]
— 你几点到的？ \emph{Nǐ jǐ diǎn dào de?} « À quelle heure es-tu arrivé ? »\\
— 我九点\cle{才}到。 \emph{Wǒ jiǔ diǎn \cle{cái} dào.} « Je ne suis arrivé \textbf{qu'à} neuf heures (plus tard que prévu). »\\
— 真晚！ \emph{Zhēn wǎn!} « Vraiment tard ! »
\end{textebox}

\begin{aretenir}[Constat d'observation]
\begin{itemize}
\item Dans les deux cas, la structure est : \textbf{Sujet + Moment/Heure + 就/才 + Verbe (+ 了)}.
\item \cle{就} signale que l'action s'est produite \textbf{tôt}, \textbf{vite} ou \textbf{plus tôt que l'attente} du locuteur.
\item \cle{才} signale que l'action s'est produite \textbf{tard}, \textbf{seulement à ce moment-là}, après une attente ou un délai.
\item L'adverbe se place \textbf{après le complément de temps} (heure, moment) et \textbf{avant le verbe}. C'est un point crucial : en français, « dès » et « seulement » précèdent souvent l'heure (« dès 7h », « seulement à 9h ») ; en chinois, l'heure précède l'adverbe.
\end{itemize}
\end{aretenir}

\courssub{Emplois et exceptions de 就}

\begin{propriete}[L'adverbe 就 — Précocité, rapidité, succession immédiate]
\textbf{Structure de base :} \textbf{Sujet + (Moment) + 就 + Verbe (+ 了)} \\[4pt]
\textbf{Emplois principaux :}
\begin{enumerate}
\item \textbf{Action précoce / « dès » :} L'action survient avant l'heure normale ou attendue.
\zh{我六点就起床了。} \emph{Wǒ liù diǎn jiù qǐchuáng le.} « Je me suis levé dès 6h. »
\item \textbf{Action rapide / « tout de suite / immédiatement » :} Souvent sans complément de temps explicite, marque la promptitude.
\zh{我马上就去。} \emph{Wǒ mǎshàng jiù qù.} « J'y vais tout de suite. »
\item \textbf{Succession immédiate / « alors / du coup » :} Lie deux actions rapprochées (souvent avec 就...了).
\zh{下了课我们就回家。} \emph{Xià le kè wǒmen jiù huí jiā.} « Une fois le cours fini, nous rentrons (immédiatement) à la maison. »
\item \textbf{Conditionnel / « si... alors » :} Dans une phrase hypothétique (si 如果/要是... 就...).
\zh{如果下雨，我就不去了。} \emph{Rúguǒ xiàyǔ, wǒ jiù bù qù le.} « S'il pleut, je n'irai pas. »
\end{enumerate}
\textbf{Exceptions / Points de vigilance :}
\begin{itemize}
\item \textbf{Pas de \zh{了} obligatoire} si l'action est habituelle ou future : \zh{他通常七点就起床} (Il se lève d'habitude à 7h).
\item \textbf{Ordre strict :} \zh{Sujet + Temps + 就 + V}. \textbf{Interdit :} \zh{*就六点我起床} / \zh{*我就六点起床}.
\end{itemize}
\end{propriete}

\courssub{Emplois et exceptions de 才}

\begin{propriete}[L'adverbe 才 — Tardivité, exclusivité, condition nécessaire]
\textbf{Structure de base :} \textbf{Sujet + (Moment) + 才 + Verbe} \textbf{(généralement sans 了 final)} \\[4pt]
\textbf{Emplois principaux :}
\begin{enumerate}
\item \textbf{Action tardive / « seulement à... / pas avant... » :} L'action survient après l'heure attendue, après un effort ou une attente.
\zh{他十点才来。} \emph{Tā shí diǎn cái lái.} « Il n'est venu qu'à 10h (il est en retard). »
\item \textbf{Nécessité d'une condition / « seulement si / ce n'est que quand... » :} Structure \zh{只有...才...} (« Ce n'est que si... que... »).
\zh{只有努力，才能成功。} \emph{Zhǐyǒu nǔlì, cáinéng chénggōng.} « Ce n'est qu'en travaillant dur qu'on peut réussir. »
\item \textbf{Quantité minimale / « à peine / seulement » :} Avec un chiffre + classificateur.
\zh{我才买了一本书。} \emph{Wǒ cái mǎi le yì běn shū.} « Je n'ai acheté qu'un seul livre. »
\end{enumerate}
\textbf{Exceptions / Points de vigilance :}
\begin{itemize}
\item \textbf{Pas de \zh{了} final} dans l'emploi « tardivité » (action non achevée dans la perspective du locuteur, ou simple constat de retard). On dit \zh{他十点才来} (pas *\zh{来了}).
\item \textbf{Ordre strict :} \zh{Sujet + Temps + 才 + V}. \textbf{Interdit :} \zh{*才十点他来}.
\item \textbf{Ne pas confondre avec \zh{刚} (à l'instant) :} \zh{刚} = action terminée il y a peu (\zh{我刚到} « Je viens d'arriver ») ; \zh{才} = action survenue tardivement (\zh{我十点才到} « Je ne suis arrivé qu'à 10h »).
\end{itemize}
\end{propriete}

\courssub{还, 又, 再 : Continuer et répéter}

Ces trois adverbes se traduisent souvent par « encore » ou « de nouveau » en français, mais leur emploi dépend de l'\textbf{aspect temporel} (passé accompli, présent en cours, futur) et de la \textbf{nature de la répétition}.

\begin{propriete}[Comparaison : 还 vs 又 vs 再]
\begin{center}
\begin{tabularx}{\linewidth}{|X|X|X|X|}
\hline
\rowcolor{popblueL} \textbf{Adverbe} & \textbf{Temps / Aspect} & \textbf{Sens principal} & \textbf{Structure type} \\
\hline
\cle{还} & Présent / Futur immédiat & \textbf{Encore / Toujours} (continuation d'un état ou action) & Sujet + 还 (+ 在) + Verbe \\
\hline
\cle{又} & \textbf{Passé accompli} & \textbf{De nouveau / Encore une fois} (répétition \textbf{réalisée}) & Sujet + 又 + Verbe + \textbf{了} \\
\hline
\cle{再} & \textbf{Futur / Volition} & \textbf{De nouveau / Une fois de plus} (répétition \textbf{projetée, souhaitée}) & Sujet + 再 + Verbe \textbf{(pas de 了)} \\
\hline
\end{tabularx}
\end{center}

\textbf{Exemples contrastifs (même verbe \zh{看} « regarder ») :}
\begin{itemize}
\item \zh{他还在看书。} \emph{Tā hái zài kàn shū.} « Il \textbf{est encore en train de} lire. » (Continuation \textbf{maintenant})
\item \zh{他昨天又看了这本书。} \emph{Tā zuótiān yòu kàn le zhè běn shū.} « Hier, il \textbf{a relu} ce livre. » (Répétition \textbf{accomplie} dans le passé)
\item \zh{他明天想再看这本书。} \emph{Tā míngtiān xiǎng zài kàn zhè běn shū.} « Demain, il \textbf{veut relire} ce livre. » (Répétition \textbf{future / intention})
\end{itemize}

\textbf{Deuxième emploi de \cle{还} : « En plus / Aussi » (addition).}
\zh{我买了苹果，还买了香蕉。} \emph{Wǒ mǎi le píngguǒ, hái mǎi le xiāngjiāo.} « J'ai acheté des pommes, \textbf{j'ai aussi acheté} des bananes. »
\end{propriete}

\begin{attention}[Ne pas confondre \zh{还} (aussi) et \zh{也} (aussi)]
\zh{也} lie deux sujets faisant la même chose (\zh{我去，他也去} « Je vais, il va aussi »).
\zh{还} lie deux actions du même sujet (\zh{我去，还买菜} « Je vais, et en plus j'achète à manger »).
\end{attention}


\courssub{Comparaison avec le français et pièges des francophones}

\begin{attention}[Erreurs fréquentes — Synthèse des pièges]
\begin{enumerate}
\item \textbf{Ordre des mots (Temps — Adverbe — Verbe) :}
\begin{itemize}
\item \textbf{Faux :} \zh{我就六点起床} / \zh{我才十点睡觉} / \zh{我再去明天}.
\item \textbf{Vrai :} \zh{我六点就起床了} / \zh{我十点才睡觉} / \zh{我明天再去}.
\item \emph{Règle :} L'adverbe (就, 才, 又, 再, 还) se place \textbf{toujours immédiatement avant le verbe} (ou avant l'auxiliaire modal \zh{想/要/能} s'il y en a un), \textbf{après le sujet et le complément de temps}.
\end{itemize}
\item \textbf{Confusion \zh{又} (passé) / \zh{再} (futur) :}
\begin{itemize}
\item \zh{又} + \zh{了} = action répétée \textbf{terminée}. \zh{我又吃了} « J'ai mangé encore (une fois de plus) ».
\item \zh{再} = action répétée \textbf{à faire}. \zh{我再吃} « Je remangerai / Je veux remanger ».
\item \emph{Piège :} \zh{昨天我再去了} $\rightarrow$ \textbf{Incorrect}. Hier = passé $\rightarrow$ \zh{昨天我又去了}.
\end{itemize}
\item \textbf{Oubli ou excès de \zh{了} :}
\begin{itemize}
\item Avec \zh{就} (précocité passée) : \zh{了} \textbf{requis} en fin de phrase. \zh{我六点就起床了}.
\item Avec \zh{才} (tardivité) : \zh{了} \textbf{généralement absent}. \zh{他十点才来} (pas *\zh{来了}).
\item Avec \zh{又} (répétition passée) : \zh{了} \textbf{requis}. \zh{他又迟到了}.
\item Avec \zh{再} (futur) : \zh{了} \textbf{interdit}. \zh{我们明天再来}.
\end{itemize}
\item \textbf{\zh{还} vs \zh{也} vs \zh{再} (sens « aussi ») :}
\begin{itemize}
\item \zh{也} : A aussi fait V (sujets différents ou même sujet, même action).
\item \zh{还} : A fait V \textbf{en plus} d'une autre action (même sujet, actions différentes) \textbf{OU} continuation (encore).
\item \zh{再} : « En plus » dans un sens quantitatif (un de plus) : \zh{再来一碗} « Encore un bol ».
\end{itemize}
\end{enumerate}
\end{attention}

\courssub{Vocabulaire de la leçon (Lexique HSK 3-4 — Vie quotidienne)}

\begin{center}
\begin{tabularx}{\linewidth}{|X|X|X|X|}
\hline
\rowcolor{popblueL} \textbf{Caractères} & \textbf{Pinyin} & \textbf{Nature} & \textbf{Traduction} \\
\hline
起床 & qǐchuáng & verbe & se lever \\
\hline
睡觉 & shuìjiào & verbe & dormir, se coucher \\
\hline
出发 & chūfā & verbe & partir, prendre le départ \\
\hline
到达 / 到 & dàodá / dào & verbe & arriver \\
\hline
等 & děng & verbe & attendre \\
\hline
迟到 & chídào & verbe & arriver en retard \\
\hline
吃饭 & chīfàn & verbe & manger (un repas) \\
\hline
做作业 & zuò zuòyè & verbe + obj & faire ses devoirs \\
\hline
看电视 & kàn diànshì & verbe + obj & regarder la télé \\
\hline
打扫 & dǎsǎo & verbe & balayer, nettoyer \\
\hline
听写 & tīngxiě & verbe & faire une dictée \\
\hline
帮忙 & bāngmáng & verbe & aider \\
\hline
特别 & tèbié & adjectif & spécial, particulier \\
\hline
累 & lèi & adjectif & fatigué \\
\hline
已经 & yǐjīng & adverbe & déjà \\
\hline
马上 & mǎshàng & adverbe & tout de suite, immédiatement \\
\hline
只有 & zhǐyǒu & conjonction & seulement si, ce n'est que \\
\hline
早上 & zǎoshang & nom (temps) & matin (matinée) \\
\hline
上午 & shàngwǔ & nom (temps) & matin (avant-midi) \\
\hline
中午 & zhōngwǔ & nom (temps) & midi \\
\hline
下午 & xiàwǔ & nom (temps) & après-midi \\
\hline
晚上 & wǎnshang & nom (temps) & soir, soirée \\
\hline
昨天 & zuótiān & nom (temps) & hier \\
\hline
今天 & jīntiān & nom (temps) & aujourd'hui \\
\hline
明天 & míngtiān & nom (temps) & demain \\
\hline
点 & diǎn & classif/nom & heure (o'clock) \\
\hline
半 & bàn & nombre & demi (heure et demie) \\
\hline
时候 & shíhou & nom & moment, quand \\
\hline
路上 & lùshang & nom (loc) & en route, sur la route \\
\hline
市场 & shìchǎng & nom & marché \\
\hline
公共汽车 / 大巴 & gōnggòng qìchē / dàbā & nom & bus (urbain / interurbain) \\
\hline
比赛 & bǐsài & nom & match, compétition \\
\hline
本子 & běnzi & nom & cahier \\
\hline
叔叔 & shūshu & nom & oncle (paternel cadet) \\
\hline
表弟 & biǎodì & nom & cousin (mâle, plus jeune) \\
\hline
\end{tabularx}
\end{center}

\courssub{Activité d'intégration : « Notre journée » (\zh{我们的一天})}

\courssub{Objectifs de l'activité}

À l'issue de cette activité d'intégration (2 h), l'élève sera capable de :

\begin{itemize}
\item rédiger un court récit au passé (6 à 8 phrases) sur sa journée de la veille, en employant correctement et spontanément les cinq adverbes étudiés : \cle{就} (action faite tôt / rapidement), \cle{才} (action faite tard / seulement à ce moment-là), \cle{还} (action encore en cours), \cle{又} (action répétée, déjà réalisée) et \cle{再} (action à refaire, projetée) ;
\item placer correctement ces adverbes \textbf{avant le verbe} et après le sujet (et le complément de temps), sans calquer l'ordre des mots du français ;
\item distinguer à l'oral comme à l'écrit \zh{又} (répétition accomplie, passé + \zh{了}) et \zh{再} (répétition à venir, futur / volition), confusion classique des francophones ;
\item mener une interview croisée en binôme à l'aide de questions fermées et ouvertes contenant ces adverbes ;
\item repérer à l'écoute les cinq adverbes dans le récit d'un camarade (compréhension orale ciblée) ;
\item s'auto-corriger grâce à la grille de relecture fournie.
\end{itemize}

\courssub{Situation}

Dans le cadre de la semaine des langues du lycée de Nkol-Eton (Yaoundé), le club de chinois prépare un petit journal mural bilingue intitulé \zh{我们的一天} (\emph{Wǒmen de yì tiān}, « Notre journée »). Chaque élève de Terminale A doit raconter une journée « pas comme les autres » : une journée où tout ne s'est pas passé comme prévu — lever très tôt, transport en commun en retard, match de football rejoué, devoir à refaire. Le récit sera affiché, puis présenté oralement, et les meilleurs textes seront envoyés à la classe partenaire d'un lycée de Tianjin.

Voici le modèle affiché au tableau par l'enseignant :

\begin{textebox}[Modèle de récit — 我昨天的一天]
昨天是很特别的一天。\\
Zuótiān shì hěn tèbié de yì tiān.\\
« Hier a été une journée très spéciale. »\\[4pt]
我早上五点\cle{就}起床了，因为要去Douala看叔叔。\\
Wǒ zǎoshang wǔ diǎn \cle{jiù} qǐchuáng le, yīnwèi yào qù Douala kàn shūshu.\\
« Je me suis levé \textbf{dès} cinq heures du matin (tôt !), parce que je devais aller voir mon oncle à Douala. »\\[4pt]
可是大巴晚上八点\cle{才}到，我们在路上等了三个小时。\\
Kěshì dàbā wǎnshang bā diǎn \cle{cái} dào, wǒmen zài lùshang děng le sān ge xiǎoshí.\\
« Mais le car n'est arrivé \textbf{qu'à} vingt heures (tard !) : nous avons attendu trois heures sur la route. »\\[4pt]
到了叔叔家，已经很晚了，可是表弟\cle{还}在看电视。\\
Dào le shūshu jiā, yǐjīng hěn wǎn le, kěshì biǎodì \cle{hái} zài kàn diànshì.\\
« Arrivés chez mon oncle, il était déjà très tard, mais mon cousin regardait \textbf{encore} la télévision. »\\[4pt]
下午我们\cle{又}看了一次喀麦隆队的比赛，真高兴！\\
Xiàwǔ wǒmen \cle{yòu} kàn le yí cì Kāmàilóng duì de bǐsài, zhēn gāoxìng!\\
« L'après-midi, nous avons \textbf{encore} regardé (une fois de plus) un match de l'équipe du Cameroun : quel bonheur ! »\\[4pt]
明天我想\cle{再}去一次Douala，这次要早点出发。\\
Míngtiān wǒ xiǎng \cle{zài} qù yí cì Douala, zhè cì yào zǎo diǎn chūfā.\\
« Demain, je voudrais \textbf{retourner} encore une fois à Douala ; cette fois, il faudra partir plus tôt. »
\end{textebox}

\courssub{Consignes}

\textbf{Tâche 1 — Compréhension et observation (10 min).} Lis le texte modèle à voix haute avec ton binôme. Souligne en rouge les cinq adverbes \zh{就、才、还、又、再}. Pour chacun, réponds par écrit : (a) quel verbe suit l'adverbe ? (b) l'adverbe exprime-t-il « tôt », « tard », « encore (continuation) », « de nouveau (fait) » ou « de nouveau (à faire) » ? (c) où se place-t-il par rapport au sujet et au complément de temps ?

\textbf{Tâche 2 — Lexique de la journée (10 min).} Avec le tableau de vocabulaire ci-dessus, complète ta liste personnelle : au moins huit verbes d'action de la vie quotidienne (se lever, partir, attendre, arriver, manger, étudier, regarder, dormir, nettoyer, faire une dictée…) et six expressions de temps (\zh{早上、上午、中午、下午、晚上、昨天、明天} + heures avec \zh{点、半}). Vérifie les tons dans le dictionnaire ou avec l'enseignant.

\textbf{Tâche 3 — Production écrite guidée (25 min).} Rédige ton propre récit intitulé \zh{我昨天的一天} en 6 à 8 phrases. Contraintes obligatoires :
\begin{enumerate}
\item une phrase avec \zh{就} : une action faite \textbf{tôt} (ex. lever, départ, arrivée) ;
\item une phrase avec \zh{才} : une action faite \textbf{tard} (ex. bus, devoir, coucher) ;
\item une phrase avec \zh{还} (+ \zh{在} si possible) : une action \textbf{encore en cours} à un moment de la journée ;
\item une phrase avec \zh{又} : une action \textbf{répétée hier} (regarder encore, manger encore, recommencer) — \textbf{avec \zh{了}} ;
\item une phrase avec \zh{再} : un projet de \textbf{refaire demain} — \textbf{sans \zh{了}} ;
\item au moins deux heures précises avec \zh{点} (ex. \zh{六点半、八点}) et un connecteur logique (\zh{可是、因为、然后、的时候}).
\end{enumerate}
Termine par la \textbf{grille d'auto-correction} :
\begin{itemize}
\item[$\square$] Chaque adverbe est-il placé \textbf{après le sujet (et le temps) et avant le verbe} ?
\item[$\square$] \zh{又} est-il bien réservé au passé accompli (avec \zh{led}) ? \zh{再} au futur / intention (sans \zh{led}) ?
\item[$\square$] \zh{才} n'a-t-il pas de \zh{led} final ? \zh{就} (passé) a-t-il son \zh{led} ?
\item[$\square$] Les heures sont-elles placées \textbf{avant} \zh{就/才} ?
\end{itemize}

\textbf{Tâche 4 — Interview croisée (15 min).} Échange ton texte avec ton binôme, puis interviewe-le à l'oral. Utilise au moins quatre de ces questions :
\begin{itemize}
\item \zh{你昨天几点才睡觉？} \emph{Nǐ zuótiān jǐ diǎn cái shuìjiào?} « À quelle heure seulement t'es-tu couché hier ? »
\item \zh{你早上几点就起床了？} \emph{Nǐ zǎoshang jǐ diǎn jiù qǐchuáng le?} « Dès quelle heure t'es-tu levé ? »
\item \zh{你昨天又做了什么？} \emph{Nǐ zuótiān yòu zuò le shénme?} « Qu'as-tu encore fait (de répété) hier ? »
\item \zh{你明天还想做什么？} \emph{Nǐ míngtiān hái xiǎng zuò shénme?} « Que veux-tu encore faire demain ? » (continuation ou addition)
\item \zh{你想再去看一次比赛吗？} \emph{Nǐ xiǎng zài qù kàn yí cì bǐsài ma?} « Veux-tu retourner voir un match une fois de plus ? »
\end{itemize}
Note les réponses de ton camarade en deux lignes ; tu devras les restituer à la troisième personne (ex. \zh{Il dit qu'il s'est couché à 23h}).

\textbf{Tâche 5 — Présentation et écoute active (20 min).} Deux volontaires présentent leur journée à la classe. Les auditeurs tiennent une fiche de repérage : à chaque adverbe entendu (\zh{就、才、还、又、再}), ils cochent une case et notent le verbe associé. À la fin, la classe vérifie collectivement si les cinq adverbes ont été utilisés et correctement placés.

\textbf{Tâche 6 — Correction collective (10 min).} L'enseignant relève au tableau deux ou trois erreurs entendues (place de l'adverbe, confusion \zh{又}/\zh{再}, oubli de \zh{了} avec \zh{又} ou présence erronée avec \zh{才}/\zh{再}) et la classe propose la correction.

\courssub{Ressources à mobiliser}

\begin{aretenir}[Rappel synthétique des cinq structures]
\begin{itemize}
\item \textbf{Sujet + Temps + 就 + Verbe + 了} : action faite tôt, vite, plus tôt que prévu. \zh{我六点就起床了。} « Je me suis levé dès six heures. »
\item \textbf{Sujet + Temps + 才 + Verbe} : action faite tard, seulement à ce moment. \zh{他十点才来。} « Il n'est arrivé qu'à dix heures. » (\textbf{Pas de 了})
\item \textbf{Sujet + 还 (+ 在) + Verbe} : encore, toujours (continuation présente). \zh{弟弟还在睡觉。} « Mon petit frère dort encore. »
\item \textbf{Sujet + 又 + Verbe + 了} : de nouveau, action répétée et \textbf{accomplie} (passé). \zh{他又迟到了。} « Il est \textbf{encore} arrivé en retard. »
\item \textbf{Sujet + 再 + Verbe} : de nouveau, action \textbf{à venir} (futur, volonté). \zh{我们明天再来。} « Nous reviendrons demain. » (\textbf{Pas de 了})
\end{itemize}
\end{aretenir}

\begin{attention}[Pièges à éviter pendant l'activité — Check-list rapide]
\begin{itemize}
\item \textbf{Ordre :} Jamais \zh{*我就六点起床} $\rightarrow$ \zh{我六点就起床了}. L'heure \textbf{précède} l'adverbe.
\item \textbf{\zh{又} vs \zh{再} :} \zh{又} = passé accompli (souvent + \zh{了}) ; \zh{再} = futur / intention (jamais \zh{了}). \zh{昨天我再来了} $\times$ $\rightarrow$ \zh{昨天我又来了} $\checkmark$.
\item \textbf{\zh{才} + \zh{了} ?} Non. \zh{他十点才来了} $\times$ $\rightarrow$ \zh{他十点才来} $\checkmark$.
\item \textbf{\zh{就} (passé) + \zh{led} ?} Oui. \zh{我六点就起床} $\times$ (inachevé) $\rightarrow$ \zh{我六点就起床了} $\checkmark$.
\item \textbf{\zh{还} $\neq$ \zh{也} :} \zh{还} = continuation (encore) OU addition (en plus, même sujet). \zh{也} = aussi (souvent sujets différents ou même action).
\end{itemize}
\end{attention}

\courssub{Proposition de production et corrigé commenté}

\begin{exoresolu}[Production modèle — 我昨天的一天]
Énoncé : rédiger un récit de 6 à 8 phrases respectant les cinq contraintes de la tâche 3, puis formuler deux réponses d'interview.
\tcblower
\textbf{Récit modèle :}\\[4pt]
\zh{我昨天的一天}\\
\emph{Wǒ zuótiān de yì tiān}\\[4pt]
\zh{昨天是星期一，我很忙。}\\
\emph{Zuótiān shì xīngqīyī, wǒ hěn máng.}\\
« Hier, c'était lundi, j'étais très occupé. »\\[4pt]
\zh{我早上五点半就起床了，因为要先去市场帮妈妈。}\\
\emph{Wǒ zǎoshang wǔ diǎn bàn jiù qǐchuáng le, yīnwèi yào xiān qù shìchǎng bāng māma.}\\
« Je me suis levé \textbf{dès} cinq heures et demie, car je devais d'abord aider ma mère au marché. »\\[4pt]
\zh{可是公共汽车七点才来，我等了很长时间。}\\
\emph{Kěshì gōnggòng qìchē qī diǎn cái lái, wǒ děng le hěn cháng shíjiān.}\\
« Mais le bus n'est arrivé \textbf{qu'à} sept heures : j'ai attendu très longtemps. »\\[4pt]
\zh{到学校的时候，同学们还在打扫教室。}\\
\emph{Dào xuéxiào de shíhou, tóngxuémen hái zài dǎsǎo jiàoshì.}\\
« Quand je suis arrivé à l'école, mes camarades nettoyaient \textbf{encore} la salle de classe. »\\[4pt]
\zh{下午老师又听写了一次汉字，我这次写得很好。}\\
\emph{Xiàwǔ lǎoshī yòu tīngxiě le yí cì Hànzì, wǒ zhè cì xiě de hěn hǎo.}\\
« L'après-midi, le professeur nous a \textbf{encore} fait une dictée de caractères ; cette fois, j'ai bien écrit. »\\[4pt]
\zh{晚上我十点才睡觉，真的太累了！}\\
\emph{Wǎnshang wǒ shí diǎn cái shuìjiào, zhēn de tài lèi le!}\\
« Le soir, je ne me suis couché \textbf{qu'à} dix heures : j'étais vraiment trop fatigué ! »\\[4pt]
\zh{明天是星期六，我想再去市场，还要买一个新的本子。}\\
\emph{Míngtiān shì xīngqīliù, wǒ xiǎng zài qù shìchǎng, hái yào mǎi yí ge xīn de běnzi.}\\
« Demain c'est samedi : je veux \textbf{retourner} au marché, et je dois \textbf{encore} acheter un cahier neuf. »\\[6pt]
\textbf{Commentaire des choix de langue :}
\begin{itemize}
\item \textbf{Phrase 2 (\zh{就}) :} \zh{就} placé entre le complément de temps \zh{五点半} et le verbe \zh{起床} ; \zh{了} final marque l'aspect accompli. L'adverbe exprime la précocité (« dès »).
\item \textbf{Phrase 3 (\zh{才}) :} \zh{才} après l'heure \zh{七点} exprime le retard du bus ; \textbf{pas de \zh{led} final} avec \zh{才} dans cet emploi de tardivité.
\item \textbf{Phrase 4 (\zh{还}) :} \zh{还在 + verbe} décrit une action en cours à un moment repère (\zh{的时候}) — structure de simultanéité (continuation).
\item \textbf{Phrase 5 (\zh{又}) :} \zh{又} + \zh{led} pour une dictée réellement refaite hier (passé accompli) ; le complément de degré \zh{写得很好} (verbe + \zh{得} + adjectif) enrichit la phrase.
\item \textbf{Phrase 7 (\zh{再} + \zh{还}) :} \zh{再} pour le projet de demain (futur, pas de \zh{led}) ; \zh{还要} au sens de « en plus, il faut encore » — deuxième emploi de \zh{还} (addition).
\item \textbf{Connecteurs variés} (\zh{因为、可是、的时候}) et deux heures précises (\zh{五点半、七点、十点}) : les contraintes sont toutes respectées.
\end{itemize}
\textbf{Réponses d'interview (modèles) :}\\[4pt]
— \zh{你昨天几点才睡觉？}\\
— \zh{我十点才睡觉，因为我还要做作业。} \emph{Wǒ shí diǎn cái shuìjiào, yīnwèi wǒ hái yào zuò zuòyè.} « Je ne me suis couché qu'à dix heures, car je devais \textbf{encore} faire mes devoirs. »\\[4pt]
— \zh{你明天还想做什么？}\\
— \zh{我明天还想再看一次比赛。} \emph{Wǒ míngtiān hái xiǎng zài kàn yí cì bǐsài.} « Demain, je veux \textbf{encore} \textbf{revoir} un match. » (Combinaison \zh{还想} + \zh{再}).
\end{exoresolu}

\courssub{Bilan de l'activité}

Cette activité a permis de mobiliser, dans une situation de communication authentique et proche de la vie des élèves (marché, transport, école, football), l'ensemble des savoir-faire de la leçon : compréhension d'un texte modèle, repérage des adverbes, production écrite contrainte, interaction orale et écoute active. Les points critiques observés en classe — place de l'adverbe avant le verbe, opposition \zh{又} (passé accompli + \zh{led}) / \zh{再} (futur / volition - \zh{led}), présence ou absence de \zh{led} avec \zh{就} et \zh{才} — ont fait l'objet d'une correction collective immédiate. Pour prolonger, les récits corrigés seront affichés sur le journal mural \zh{我们的一天} et les deux meilleurs seront envoyés au lycée partenaire de Tianjin. En devoir maison, chaque élève rédigera le récit de la journée d'un membre de sa famille (\zh{我妈妈的一天} ou \zh{我爸爸的一天}) en réutilisant au moins quatre des cinq adverbes, afin de transférer les acquis à la troisième personne et de consolider l'ordre des mots.

\newpage

\coursec{Méthodes et exercices résolus}

\courssub{Fiches méthodes et exercices résolus}

\begin{methode}[Construire correctement une phrase avec 就 ou 才]
Pour exprimer « tôt / facilement » (就) ou « tard / seulement » (才) avec un moment ou une quantité :
\begin{enumerate}
\item Repérer le \cle{point de référence} : l'heure, l'âge, la durée ou la quantité mentionnée dans la phrase.
\item Se demander quel est le jugement du locuteur : c'est \emph{plus tôt / plus facile / plus rapide} que prévu $\Rightarrow$ \zh{就} ; c'est \emph{plus tard / plus difficile} que prévu $\Rightarrow$ \zh{才}.
\item Placer l'adverbe \cle{juste avant le verbe}, jamais avant le sujet : Sujet + moment/quantité + 就/才 + Verbe (+ 了).
\item Vérifier la négation : on ne dit pas \zh{他不八点就来} ; formes correctes : \zh{他不是八点就来} « ce n'est pas à huit heures qu'il vient » ou \zh{他不到八点就来了} « il est arrivé avant huit heures ».
\item Lire la phrase à voix haute en respectant les tons : \emph{jiù} (4\textsuperscript{e} ton), \emph{cái} (2\textsuperscript{e} ton).
\end{enumerate}
Structure à retenir : \textbf{Sujet + (complément de temps/quantité) + 就/才 + Verbe + (了)}
\end{methode}

\begin{aretenir}[就 et 才 : le jugement du locuteur]
\begin{center}
\begin{tabularx}{\linewidth}{|X|X|X|}\hline
\rowcolor{popblueL} Structure & Exemple (caractères + pinyin) & Traduction \\ \hline
Sujet + moment + 就 + V + 了 & \zh{他七点就来了。}\emph{Tā qī diǎn jiù lái le.} & Il est arrivé dès sept heures (tôt). \\ \hline
Sujet + moment + 才 + V & \zh{他七点才来。}\emph{Tā qī diǎn cái lái.} & Il n'est arrivé qu'à sept heures (tard). \\ \hline
Sujet + quantité + 就 + V + 了 & \zh{我五分钟就吃完了。}\emph{Wǒ wǔ fēnzhōng jiù chī wán le.} & J'ai fini de manger en cinq minutes (vite). \\ \hline
Sujet + quantité + 才 + V & \zh{他学了三年才会说。}\emph{Tā xué le sān nián cái huì shuō.} & Il a fallu trois ans d'étude pour qu'il sache parler. \\ \hline
\end{tabularx}
\end{center}
Remarque : avec 就, la marque \zh{了} est fréquente (action accomplie plus tôt que prévu) ; avec 才 suivi d'un fait passé, on emploie souvent la phrase sans 了 final (\zh{他昨天才回来。}).
\end{aretenir}

\begin{exoresolu}[Choisir entre 就 et 才]
\textbf{Énoncé.} Complétez avec 就 ou 才, puis traduisez :\\
1. 他六岁\_\_\_上学了。\\
2. 飞机晚上十一点\_\_\_到雅温得。\\
3. 我等了两个小时，他\_\_\_来。\\
4. 从我家到学校很近，走五分钟\_\_\_到了。
\tcblower
\textbf{Solution détaillée.}\\
1. 他六岁\textbf{就}上学了。\emph{Tā liù suì jiù shàngxué le.} « Il est entré à l'école dès six ans. » Justification : six ans, c'est \emph{tôt} par rapport à l'âge habituel (sept ans) ; le locuteur souligne la précocité $\Rightarrow$ 就.\\
2. 飞机晚上十一点\textbf{才}到雅温得。\emph{Fēijī wǎnshang shíyī diǎn cái dào Yǎwēndé.} « L'avion n'est arrivé à Yaoundé qu'à onze heures du soir. » Justification : onze heures du soir, c'est \emph{tard} ; le locuteur exprime l'attente ou la déception $\Rightarrow$ 才.\\
3. 我等了两个小时，他\textbf{才}来。\emph{Wǒ děng le liǎng ge xiǎoshí, tā cái lái.} « J'ai attendu deux heures avant qu'il n'arrive (enfin). » Justification : après une longue attente, l'arrivée est \emph{tardive} $\Rightarrow$ 才.\\
4. 走五分钟\textbf{就}到了。\emph{Zǒu wǔ fēnzhōng jiù dào le.} « En cinq minutes de marche, on y est (déjà). » Justification : cinq minutes, c'est \emph{peu}, le trajet est rapide $\Rightarrow$ 就.\\
\textbf{Conclusion.} Le choix ne dépend pas du fait objectif mais du \emph{jugement du locuteur} : 就 = « déjà, dès » ; 才 = « seulement, pas avant ».
\end{exoresolu}

\begin{methode}[Employer 就 pour enchaîner deux actions (一……就……)]
Pour dire « dès que…, aussitôt… » :
\begin{enumerate}
\item Identifier les deux actions : l'action déclencheur (A) et la conséquence immédiate (B).
\item Construire le schéma : \textbf{一 + Verbe A + 就 + Verbe B}. Exemple : \zh{我一到家就写作业。}
\item Si le sujet des deux actions est le même, on ne le répète qu'une fois, au début.
\item Si les sujets sont différents, placer le deuxième sujet \emph{après} 就 : \zh{我一回家，妈妈就做饭。}
\item Ne pas confondre avec le français « dès que » suivi d'un futur : en chinois, pas de marque de futur, 就 suffit.
\end{enumerate}
\end{methode}

\begin{exoresolu}[Transformation avec 一……就……]
\textbf{Énoncé.} Transformez les phrases suivantes avec la structure 一……就……, puis traduisez :\\
1. 他回到家。他马上吃饭。\\
2. 下课了。学生们去操场。\\
3. 我到了杜阿拉。我给你们打电话。
\tcblower
\textbf{Solution détaillée.}\\
1. Même sujet (il) $\Rightarrow$ un seul sujet au début : 他\textbf{一}回到家\textbf{就}吃饭。\emph{Tā yī huí dào jiā jiù chīfàn.} « Dès qu'il rentre à la maison, il mange. »\\
2. Même sujet (les élèves) : 学生们\textbf{一}下课\textbf{就}去操场。\emph{Xuéshengmen yī xiàkè jiù qù cāochǎng.} « Dès que le cours finit, les élèves vont au terrain de sport. »\\
3. Même sujet (je) : 我\textbf{一}到杜阿拉\textbf{就}给你们打电话。\emph{Wǒ yī dào Dù'ālā jiù gěi nǐmen dǎ diànhuà.} « Dès que j'arrive à Douala, je vous téléphone. »\\
\textbf{Conclusion.} La structure 一……就…… supprime les mots français « dès que » et « aussitôt » : la rapidité de l'enchaînement est portée par 就 seul. Attention à ne pas ajouter 了 après le premier verbe dans cette structure.
\end{exoresolu}

\begin{methode}[Distinguer 还, 又 et 再]
Trois adverbes, trois idées différentes :
\begin{enumerate}
\item \cle{还} (hái) = « encore, toujours » : un état ou une action \emph{continue} au présent. Schéma : Sujet + 还 + Verbe/adjectif. \zh{他还在学习。} « Il étudie encore. »
\item \cle{又} (yòu) = « de nouveau » : une action \emph{déjà répétée} (passé ou constat). Souvent avec 了. \zh{他又迟到了。} « Il est encore arrivé en retard. »
\item \cle{再} (zài) = « encore une fois » : une répétition \emph{projetée} (futur, souhait, ordre). \zh{请再说一遍。} « Répétez encore une fois, s'il vous plaît. »
\item Réflexe de vérification : l'action est-elle en cours ? $\Rightarrow$ 还. Déjà refaite ? $\Rightarrow$ 又. À refaire ? $\Rightarrow$ 再.
\item Place : toujours \emph{avant le verbe}, après le sujet.
\end{enumerate}
\end{methode}

\begin{aretenir}[Tableau récapitulatif : 还, 又, 再]
\begin{center}
\begin{tabularx}{\linewidth}{|X|X|X|}\hline
\rowcolor{popblueL} Adverbe et valeur & Exemple (caractères + pinyin) & Traduction \\ \hline
还 : continuité (présent) & \zh{他还在睡觉。}\emph{Tā hái zài shuìjiào.} & Il dort encore. \\ \hline
又 : répétition accomplie (passé) & \zh{他又迟到了。}\emph{Tā yòu chídào le.} & Il est encore arrivé en retard. \\ \hline
再 : répétition projetée (futur) & \zh{我们明天再来。}\emph{Wǒmen míngtiān zài lái.} & Nous reviendrons demain. \\ \hline
\end{tabularx}
\end{center}
Cas particulier : \zh{再见} « au revoir » signifie littéralement « se revoir encore » — preuve que 再 regarde vers l'avenir.
\end{aretenir}

\begin{exoresolu}[Choisir entre 还, 又 et 再]
\textbf{Énoncé.} Complétez avec 还, 又 ou 再, puis traduisez :\\
1. 明天请你\_\_\_来一次。\\
2. 都十点了，他\_\_\_没起床。\\
3. 这个电影我看过，昨天我\_\_\_看了一遍。\\
4. 喀麦隆队\_\_\_进了一个球！（比赛刚结束）\\
5. 我想\_\_\_学一年汉语。
\tcblower
\textbf{Solution détaillée.}\\
1. 明天请你\textbf{再}来一次。\emph{Míngtiān qǐng nǐ zài lái yī cì.} « Revenez encore une fois demain. » Justification : la visite est \emph{à venir}, c'est un projet $\Rightarrow$ 再.\\
2. 都十点了，他\textbf{还}没起床。\emph{Dōu shí diǎn le, tā hái méi qǐchuáng.} « Il est déjà dix heures et il n'est toujours pas levé. » Justification : l'état (dormir) \emph{continue} $\Rightarrow$ 还.\\
3. 昨天我\textbf{又}看了一遍。\emph{Zuótiān wǒ yòu kàn le yī biàn.} « Hier, je l'ai regardé une fois de plus. » Justification : l'action est \emph{terminée et répétée} (hier + 了) $\Rightarrow$ 又.\\
4. 喀麦隆队\textbf{又}进了一个球！\emph{Kāmàilóng duì yòu jìn le yī ge qiú!} « L'équipe du Cameroun a encore marqué un but ! » Justification : le but vient d'être marqué, c'est un fait accompli $\Rightarrow$ 又.\\
5. 我想\textbf{再}学一年汉语。\emph{Wǒ xiǎng zài xué yī nián Hànyǔ.} « Je voudrais étudier le chinois encore un an. » Justification : projet futur $\Rightarrow$ 再.\\
\textbf{Conclusion.} Le critère décisif est le \emph{temps} : 又 regarde en arrière, 再 regarde en avant, 还 décrit le présent continu.
\end{exoresolu}

\begin{methode}[Construire un court texte avec 就, 才, 还, 又, 再]
Pour une production écrite (type Bac : 5 à 8 phrases) :
\begin{enumerate}
\item Rédiger d'abord un plan simple en français : situation, moment, jugement (tôt/tard), continuité ou répétition.
\item Pour chaque phrase, choisir l'adverbe selon la logique : précocité $\Rightarrow$ 就 ; retard $\Rightarrow$ 才 ; continuité $\Rightarrow$ 还 ; répétition passée $\Rightarrow$ 又 ; répétition future $\Rightarrow$ 再.
\item Vérifier la place de chaque adverbe : \emph{après le sujet, avant le verbe}.
\item Contrôler les tons en relisant le pinyin : \emph{jiù, cái, hái, yòu, zài}.
\item Terminer par une phrase de conclusion naturelle, par exemple avec \zh{以后} « plus tard » ou \zh{所以} « donc ».
\end{enumerate}
\end{methode}

\begin{exoresolu}[Composition guidée : une journée chargée]
\textbf{Énoncé.} Rédigez un court texte (5 phrases minimum) sur le thème « Une journée chargée à Yaoundé », en utilisant au moins une fois 就, 才, 还 et 又 ou 再.
\tcblower
\textbf{Solution détaillée.}\\
Plan : lever tôt (就) $\rightarrow$ embouteillages, arrivée tardive au lycée (才) $\rightarrow$ le soir, encore des devoirs (还) $\rightarrow$ un match revu (又) $\rightarrow$ projet de se coucher tôt demain (再).\\
Texte rédigé :\\
今天早上我五点\textbf{就}起床了。\emph{Jīntiān zǎoshang wǒ wǔ diǎn jiù qǐchuáng le.} « Ce matin, je me suis levé dès cinq heures. »\\
可是路上车很多，我八点\textbf{才}到学校。\emph{Kěshì lùshang chē hěn duō, wǒ bā diǎn cái dào xuéxiào.} « Mais il y avait beaucoup de voitures sur la route, je ne suis arrivé au lycée qu'à huit heures. »\\
晚上回到家，我\textbf{还}要写作业。\emph{Wǎnshang huí dào jiā, wǒ hái yào xiě zuòyè.} « Le soir, rentré à la maison, je dois encore faire mes devoirs. »\\
昨天我\textbf{又}看了一遍喀麦隆队的比赛。\emph{Zuótiān wǒ yòu kàn le yī biàn Kāmàilóng duì de bǐsài.} « Hier, j'ai encore regardé le match de l'équipe du Cameroun. »\\
明天我想早点儿睡觉，以后\textbf{再}看电视。\emph{Míngtiān wǒ xiǎng zǎo diǎnr shuìjiào, yǐhòu zài kàn diànshì.} « Demain je veux me coucher plus tôt, je regarderai la télévision plus tard. »\\
\textbf{Conclusion.} Chaque adverbe est placé avant le verbe et justifié par la logique temporelle : 就 (tôt), 才 (tard), 还 (continuité), 又 (passé répété), 再 (futur projeté).
\end{exoresolu}

\begin{methode}[Traduire du français vers le chinois avec 就/才/还/又/再]
Méthode de version en quatre étapes :
\begin{enumerate}
\item Souligner dans la phrase française les mots-clés : « dès, déjà » $\Rightarrow$ 就 ; « seulement, ne… que, pas avant » $\Rightarrow$ 才 ; « encore, toujours » $\Rightarrow$ 还 ; « de nouveau » (fait passé) $\Rightarrow$ 又 ; « encore, une autre fois » (futur) $\Rightarrow$ 再.
\item Identifier sujet, verbe et compléments ; poser le squelette : Sujet + Verbe.
\item Insérer l'adverbe \emph{avant le verbe} et le complément de temps \emph{avant l'adverbe} : Sujet + moment + 就/才 + Verbe.
\item Relire : ordre des mots, tons du pinyin, ponctuation chinoise pleine chasse （。 ， ？）.
\end{enumerate}
\end{methode}

\begin{exoresolu}[Version : du français au chinois]
\textbf{Énoncé.} Traduisez en chinois (caractères + pinyin) :\\
1. « Mon frère n'est rentré qu'à minuit. »\\
2. « Elle parle déjà très bien chinois après seulement six mois. »\\
3. « Il pleut encore aujourd'hui. »\\
4. « Nous visiterons encore une fois le palais des rois Bamoun. »
\tcblower
\textbf{Solution détaillée.}\\
1. Mot-clé : « ne… qu'à » $\Rightarrow$ 才. Squelette : 我哥哥 + 回来. Insertion : 我哥哥半夜十二点\textbf{才}回来。\emph{Wǒ gēge bànyè shí'èr diǎn cái huílái.}\\
2. Mot-clé : « déjà » + « seulement six mois » (durée courte, jugement positif) $\Rightarrow$ 就. 她学了六个月汉语，\textbf{就}说得很好了。\emph{Tā xué le liù ge yuè Hànyǔ, jiù shuō de hěn hǎo le.} Remarque : \zh{说得很好} avec 得 (complément de degré).\\
3. Mot-clé : « encore » (état qui continue) $\Rightarrow$ 还. 今天\textbf{还}下雨。\emph{Jīntiān hái xiàyǔ.}\\
4. Mot-clé : « encore une fois » (projet futur) $\Rightarrow$ 再. 我们\textbf{要再}参观一次巴蒙国王的宫殿。\emph{Wǒmen yào zài cānguān yī cì Bāméng guówáng de gōngdiàn.} Remarques : \zh{要} marque ici le projet futur ; \zh{一次} « une fois » se place après le verbe.\\
\textbf{Conclusion.} En version, ne jamais traduire mot à mot : c'est l'\emph{idée temporelle} du mot français qui commande le choix de l'adverbe chinois.
\end{exoresolu}

\begin{attention}[Les 5 erreurs qui coûtent des points au Bac]
\begin{enumerate}
\item \textbf{Confondre 又 et 再.} Dire \zh{我明天又去} pour « j'irai encore demain » est faux : le futur exige \textbf{再} (\zh{我明天再去}). Réservez 又 au passé constaté.
\item \textbf{Mal placer l'adverbe.} \zh{就我八点起床} est incorrect : 就/才/还/又/再 se placent \emph{après le sujet, avant le verbe} : \zh{我八点就起床了。}
\item \textbf{Confondre les caractères.} 就 (jiù, « dès ») n'a rien à voir avec 旧 (jiù, « vieux ») ; 才 (cái) se distingue de 木 (mù) par le trait oblique ; 再 (zài) ≠ 在 (zài, « être à ») : homophones parfaits, sens totalement différents. Au Bac, un caractère faux fait perdre le point même si le pinyin est juste.
\item \textbf{Calquer le français « dès que ».} N'écrivez pas \zh{当……的时候，就……} pour « dès que » : la structure correcte est \textbf{一……就……} (\zh{我一到家就吃饭。}), sans 了 après le premier verbe.
\item \textbf{Négliger les tons dans le pinyin.} \emph{hái} (还, 2\textsuperscript{e} ton) sans ton ou avec un mauvais ton est sanctionné ; de même \emph{jiù} (4\textsuperscript{e} ton) et \emph{cái} (2\textsuperscript{e} ton). Relisez toujours votre pinyin syllabe par syllabe.
\end{enumerate}
\end{attention}

\newpage

\coursec{Exercices}

\courssub{Je vérifie mes connaissances}

\begin{exercice}[QCM : choisir le bon adverbe]
Pour chaque phrase, choisis la bonne réponse parmi A, B, C et D.

1. 电影七点开始，我们六点半\underline{\hspace{1.5cm}}到了。\\
A. 才 \hspace{0.8cm} B. 就 \hspace{0.8cm} C. 又 \hspace{0.8cm} D. 再

2. 他昨天晚上十二点\underline{\hspace{1.5cm}}回家。\\
A. 就 \hspace{0.8cm} B. 才 \hspace{0.8cm} C. 还 \hspace{0.8cm} D. 又

3. 这个菜真好吃，我想\underline{\hspace{1.5cm}}吃一碗。\\
A. 又 \hspace{0.8cm} B. 还 \hspace{0.8cm} C. 再 \hspace{0.8cm} D. 就

4. 他\underline{\hspace{1.5cm}}生病了，这个月第二次了。\\
A. 再 \hspace{0.8cm} B. 就 \hspace{0.8cm} C. 才 \hspace{0.8cm} D. 又
\end{exercice}

\begin{exercice}[Vrai ou faux ? Justifie ta réponse]
Dis si chaque phrase est correcte (V) ou fautive (F), puis justifie ta réponse en français en citant la règle de grammaire concernée.

1. 他才来了昨天。

2. 她十五岁就会说三种语言了。

3. 请又说一遍。

4. 我明天再来看你。

5. 我还昨天去了市场。
\end{exercice}

\begin{exercice}[Vocabulaire : complète le tableau]
Complète le tableau suivant en donnant le pinyin (avec les tons) et la traduction française de chaque adverbe, puis précise son emploi principal en une phrase courte en français.

\begin{center}
\begin{tabularx}{\linewidth}{|X|X|X|X|}
\hline
\rowcolor{popblueL} Caractères & Pinyin & Traduction & Emploi principal \\
\hline
就 & & & \\
\hline
才 & & & \\
\hline
还 & & & \\
\hline
又 & & & \\
\hline
再 & & & \\
\hline
\end{tabularx}
\end{center}
\end{exercice}

\begin{exercice}[Associe le pinyin aux caractères]
Associe chaque caractère de la colonne de gauche au bon pinyin de la colonne de droite, puis donne la traduction française de chaque mot.

\begin{center}
\begin{tabularx}{\linewidth}{|X|X|X|}
\hline
\rowcolor{popblueL} Caractères & Pinyin & Traduction \\
\hline
1. 就 & a. cái & \\
\hline
2. 才 & b. hái & \\
\hline
3. 还 & c. zài & \\
\hline
4. 又 & d. jiù & \\
\hline
5. 再 & e. yòu & \\
\hline
\end{tabularx}
\end{center}
\end{exercice}

\courssub{Je m'entraîne}

\begin{exercice}[Traduis en français]
Traduis les phrases suivantes en français. Attention au sens exact apporté par l'adverbe (jugement de tôt/tard, répétition, continuation).

1. 我们六点半就到了。\\
\emph{Wǒmen liù diǎn bàn jiù dào le.}

2. 他昨天晚上十二点才回家。\\
\emph{Tā zuótiān wǎnshang shí'èr diǎn cái huí jiā.}

3. 我的弟弟还小，才五岁。\\
\emph{Wǒ de dìdi hái xiǎo, cái wǔ suì.}

4. 他又生病了，这个月第二次了。\\
\emph{Tā yòu shēngbìng le, zhège yuè dì-èr cì le.}

5. 这个菜真好吃，我想再吃一碗。\\
\emph{Zhège cài zhēn hǎochī, wǒ xiǎng zài chī yì wǎn.}
\end{exercice}

\begin{exercice}[Texte à trous : 就, 才, 还, 又 ou 再 ?]
Complète chaque phrase avec l'adverbe qui convient : \zh{就}, \zh{才}, \zh{还}, \zh{又} ou \zh{再}. Justifie ton choix en français.

1. 电影七点\underline{\hspace{1.2cm}}开始了，我们六点半就到了。

2. 他昨天晚上十二点\underline{\hspace{1.2cm}}回家。

3. 我的弟弟\underline{\hspace{1.2cm}}小，才五岁。

4. 他\underline{\hspace{1.2cm}}生病了，这个月第二次了。

5. 这个菜真好吃，我想\underline{\hspace{1.2cm}}吃一碗。

6. 作业太多了，我做了三个小时\underline{\hspace{1.2cm}}做完。

7. 别走！我\underline{\hspace{1.2cm}}有一个问题要问你。

8. 今天的课听不懂，请老师\underline{\hspace{1.2cm}}讲一遍。
\end{exercice}

\begin{exercice}[Remets les mots dans l'ordre]
Remets les mots suivants dans le bon ordre pour former une phrase correcte, puis traduis-la en français.

1. 才 / 他 / 起床 / 九点 / 早上 / 今天

2. 就 / 我 / 完了 / 作业 / 做 / 一个小时

3. 又 / 迟到 / 他 / 今天 / 了

4. 还 / 在 / 睡觉 / 弟弟 / 呢

5. 再 / 请 / 说 / 一遍 / 你
\end{exercice}

\begin{exercice}[Transformation : exprime un jugement]
Réécris chaque phrase en utilisant \zh{就} ou \zh{才} pour exprimer le jugement indiqué entre parenthèses, puis traduis les deux phrases obtenues en français.

1. 他八点来。\\
\emph{Tā bā diǎn lái.} (Il arrive à huit heures.)\\
a) tôt (c'est tôt, c'est bien) : \ldots\\
b) tard (c'est tard, ce n'est pas bien) : \ldots

2. 我用了三十分钟做完作业。\\
\emph{Wǒ yòng le sānshí fēnzhōng zuòwán zuòyè.} (J'ai mis trente minutes à finir mes devoirs.)\\
a) facile (c'était facile, c'est rapide) : \ldots\\
b) difficile (c'était difficile, c'est long) : \ldots

3. 她十八岁大学毕业。\\
\emph{Tā shíbā suì dàxué bìyè.} (Elle a fini l'université à dix-huit ans.)\\
a) jeune et remarquable : \ldots\\
b) tard par rapport aux autres : \ldots
\end{exercice}

\courssub{Je me prépare au Bac}

\begin{exercice}[Type Bac — Compréhension de l'écrit et questions de langue]
Lis le texte suivant, puis réponds aux questions.

\begin{textebox}[Une journée d'examen]
今天是汉语考试的日子。我早上五点半就起床了，因为我还想复习一遍生词。\\
\emph{Jīntiān shì Hànyǔ kǎoshì de rìzi. Wǒ zǎoshang wǔ diǎn bàn jiù qǐchuáng le, yīnwèi wǒ hái xiǎng fùxí yí biàn shēngcí.}\\
« Aujourd'hui, c'est le jour de l'examen de chinois. Je me suis levé dès cinq heures et demie du matin, parce que je voulais encore réviser une fois le vocabulaire. »\\[4pt]
考试八点开始，我七点半就到了学校。可是考试太难了，我用了两个小时才做完。\\
\emph{Kǎoshì bā diǎn kāishǐ, wǒ qī diǎn bàn jiù dào le xuéxiào. Kěshì kǎoshì tài nán le, wǒ yòng le liǎng ge xiǎoshí cái zuòwán.}\\
« L'examen commence à huit heures, je suis arrivé à l'école dès sept heures et demie. Mais l'examen était trop difficile, j'ai mis deux heures à le finir. »\\[4pt]
考完以后，老师说："下个星期还有一次考试，大家再努力吧！"唉，又要考试了！\\
\emph{Kǎowán yǐhòu, lǎoshī shuō : "Xià ge xīngqī hái yǒu yí cì kǎoshì, dàjiā zài nǔlì ba !" Ài, yòu yào kǎoshì le !}\\
« Après l'examen, le professeur a dit : "La semaine prochaine, il y a encore un examen, faites encore des efforts !" Hélas, encore un examen ! »
\end{textebox}

\textbf{A. Questions de compréhension} (réponds en français)

1. Pourquoi le narrateur s'est-il levé si tôt ce matin-là ?
2. À quelle heure commence l'examen ? À quelle heure le narrateur arrive-t-il à l'école ?
3. L'examen était-il facile ? Justifie ta réponse à l'aide du texte.
4. Quelle annonce le professeur fait-il à la fin du texte ?

\textbf{B. Questions de langue}

1. Dans la phrase \zh{我早上五点半就起床了}, que exprime \zh{就} ? Recopie la règle et donne un autre exemple de ta composition.
2. Dans \zh{我用了两个小时才做完}, pourquoi utilise-t-on \zh{才} et non \zh{就} ?
3. Relevez dans le texte deux emplois différents de \zh{还} et expliquez chacun en français.
4. Pourquoi dit-on \zh{又要考试了} et non \zh{再要考试了} à la fin du texte ?
\end{exercice}

\begin{exercice}[Type Bac — Thème et version]
\textbf{A. Thème.} Traduis les phrases suivantes en chinois (caractères simplifiés), en choisissant correctement l'adverbe.

1. « Il n'est arrivé qu'à dix heures hier soir. »
2. « Dès que je rentre, je fais mes devoirs. »
3. « Il dort encore, ne le réveille pas. »
4. « Il est encore en retard aujourd'hui ! »
5. « Revenez demain, s'il vous plaît. »

\textbf{B. Version.} Traduis les phrases suivantes en français.

1. 只有多练习，才能学好汉语。\\
\emph{Zhǐyǒu duō liànxí, cái néng xuéhǎo Hànyǔ.}

2. 她十五岁就会说三种语言了。\\
\emph{Tā shíwǔ suì jiù huì shuō sān zhǒng yǔyán le.}

3. 他又忘了带书，老师很不高兴。\\
\emph{Tā yòu wàng le dài shū, lǎoshī hěn bù gāoxìng.}
\end{exercice}

\begin{exercice}[Type Bac — Correction de phrases et expression écrite guidée]
\textbf{A. Correction de phrases fautives.} Chaque phrase suivante contient une erreur. Identifie l'erreur, écris la phrase corrigée en caractères, puis justifie la correction en citant la règle de grammaire.

1. 他才来了昨天。

2. 请又说一遍。

3. 我还昨天去了市场。

4. 他明天又来看你。

\textbf{B. Expression écrite guidée.} Rédige en chinois un paragraphe de \textbf{5 à 8 phrases} (environ 80 à 120 caractères) racontant une journée d'examen au lycée (au Cameroun). Tu utiliseras \textbf{au moins trois} des cinq adverbes étudiés : \zh{就}, \zh{才}, \zh{还}, \zh{又}, \zh{再}.

Pistes : l'heure du lever, le trajet jusqu'à l'école, le déroulement de l'épreuve, tes sentiments, la fin de la journée.

\textbf{Barème indicatif :} exactitude de la place des adverbes dans la phrase (4 pts) ; choix correct entre \zh{又} (passé) et \zh{再} (futur) (2 pts) ; emploi correct de \zh{了} (2 pts) ; richesse du vocabulaire et correction des caractères (2 pts).
\end{exercice}

\newpage

\coursec{Corrigés des exercices}

\begin{corrige}[Exercice 1 — QCM : choisir le bon adverbe]

\textbf{1. B. \cle{就}}\\
\zh{电影七点开始，我们六点半就到了。}\\
\emph{Diànyǐng qī diǎn kāishǐ, wǒmen liù diǎn bàn jiù dào le.}\\
« Le film commence à sept heures, nous sommes arrivés dès six heures et demie. »\\
Justification : l'arrivée (6 h 30) a lieu \emph{avant} le moment attendu (7 h) ; le locuteur juge que c'est tôt. Règle : \zh{就} + verbe exprime une action accomplie plus tôt ou plus facilement que prévu (jugement positif).

\textbf{2. B. \cle{才}}\\
\zh{他昨天晚上十二点才回家。}\\
\emph{Tā zuótiān wǎnshang shí'èr diǎn cái huí jiā.}\\
« Hier soir, il n'est rentré qu'à minuit. »\\
Justification : minuit est jugé tard ; \zh{才} + verbe exprime une action accomplie plus tard ou plus difficilement que prévu (jugement négatif).

\textbf{3. C. \cle{再}}\\
\zh{这个菜真好吃，我想再吃一碗。}\\
\emph{Zhège cài zhēn hǎochī, wǒ xiǎng zài chī yì wǎn.}\\
« Ce plat est vraiment bon, je veux en manger encore un bol. »\\
Justification : la répétition est projetée (futur, « une fois de plus ») ; on utilise \zh{再}. \zh{又} est réservé à une répétition déjà réalisée.

\textbf{4. D. \cle{又}}\\
\zh{他又生病了，这个月第二次了。}\\
\emph{Tā yòu shēngbìng le, zhège yuè dì'èr cì le.}\\
« Il est encore tombé malade, c'est la deuxième fois ce mois-ci. »\\
Justification : la répétition est déjà réalisée (passé, avec \zh{了}) ; on utilise \zh{又}.
\end{corrige}

\begin{corrige}[Exercice 2 — Vrai ou faux ? Justifie ta réponse]

\textbf{1. F — \zh{他才来了昨天。}}\\
Correction : \zh{他昨天才来。} \emph{Tā zuótiān cái lái.} « Il n'est arrivé qu'hier. »\\
Règle : le complément de temps (\zh{昨天}) se place \emph{avant} l'adverbe \zh{才} et le verbe, jamais après. Schéma : Sujet + temps + \zh{才} + verbe. De plus, \zh{才} n'appelle pas \zh{了} après le verbe.

\textbf{2. V — \zh{她十五岁就会说三种语言了。}}\\
\emph{Tā shíwǔ suì jiù huì shuō sān zhǒng yǔyán le.} « À quinze ans, elle savait déjà parler trois langues. »\\
Règle : \zh{就} devant un verbe (ici \zh{会}) exprime que l'action est accomplie plus tôt / plus facilement que prévu ; jugement positif du locuteur. Phrase correcte.

\textbf{3. F — \zh{请又说一遍。}}\\
Correction : \zh{请再说一遍。} \emph{Qǐng zài shuō yí biàn.} « Répétez une fois, s'il vous plaît. »\\
Règle : pour demander de répéter une action \emph{à venir}, on utilise \zh{再} (futur). \zh{又} ne s'emploie que pour une répétition déjà réalisée (passé).

\textbf{4. V — \zh{我明天再来看你。}}\\
\emph{Wǒ míngtiān zài lái kàn nǐ.} « Je reviendrai te voir demain. »\\
Règle : \zh{再} + verbe exprime une action qui se fera plus tard (futur) ; avec \zh{明天}, c'est correct.

\textbf{5. F — \zh{我还昨天去了市场。}}\\
Correction : \zh{我昨天还去了市场。} \emph{Wǒ zuótiān hái qù le shìchǎng.} « Hier, je suis en plus allé au marché. »\\
Règle : le complément de temps \zh{昨天} se place juste après le sujet, \emph{avant} l'adverbe \zh{还}. Schéma : Sujet + temps + \zh{还} + verbe.
\end{corrige}

\begin{corrige}[Exercice 3 — Vocabulaire : complète le tableau]

\begin{center}
\begin{tabularx}{\linewidth}{|X|X|X|X|}
\hline
\rowcolor{popblueL} Caractères & Pinyin & Nature & Traduction \\
\hline
就 & jiù & adv. & dès, aussitôt, seulement (tôt) ; action accomplie plus tôt ou plus facilement que prévu \\
\hline
才 & cái & adv. & seulement, ne\ldots que (tard) ; action accomplie plus tard ou plus difficilement que prévu \\
\hline
还 & hái & adv. & encore, aussi, en plus ; continuation d'un état ou ajout d'un élément \\
\hline
又 & yòu & adv. & de nouveau, encore (passé) ; répétition d'une action déjà réalisée (souvent avec \zh{了}) \\
\hline
再 & zài & adv. & encore, de nouveau (futur) ; répétition projetée ou action remise à plus tard \\
\hline
\end{tabularx}
\end{center}

\cle{Point de vigilance} : \zh{又} regarde vers le passé (répétition constatée), \zh{再} vers le futur (répétition voulue ou programmée) ; \zh{还} n'exprime ni l'un ni l'autre mais la continuation ou l'ajout.
\end{corrige}

\begin{corrige}[Exercice 4 — Associe le pinyin aux caractères]

\begin{center}
\begin{tabularx}{\linewidth}{|X|X|X|}
\hline
\rowcolor{popblueL} Caractères & Pinyin & Traduction \\
\hline
1. 就 & d. jiù & dès, aussitôt ; alors \\
\hline
2. 才 & a. cái & seulement, ne\ldots que \\
\hline
3. 还 & b. hái & encore, aussi \\
\hline
4. 又 & e. yòu & de nouveau (action réalisée) \\
\hline
5. 再 & c. zài & encore, de nouveau (à venir) \\
\hline
\end{tabularx}
\end{center}

\cle{Attention} : \zh{还} a deux lectures : \emph{hái} (« encore », adverbe) et \emph{huán} (« rendre », verbe, comme dans \zh{还书} \emph{huán shū}, « rendre un livre »). Ici, c'est l'adverbe, donc \emph{hái}.
\end{corrige}

\begin{corrige}[Exercice 5 — Traduis en français]

\textbf{1.} \zh{我们六点半就到了。}\\
« Nous sommes arrivés dès six heures et demie. »\\
\zh{就} marque un jugement : c'est plus tôt que prévu (nuance de satisfaction).

\textbf{2.} \zh{他昨天晚上十二点才回家。}\\
« Hier soir, il n'est rentré qu'à minuit. »\\
\zh{才} marque un jugement : c'est plus tard que prévu (nuance de reproche ou d'étonnement).

\textbf{3.} \zh{我的弟弟还小，才五岁。}\\
« Mon petit frère est encore petit, il n'a que cinq ans. »\\
\zh{还} exprime la continuation d'un état (« encore ») ; \zh{才} insiste sur la petitesse du chiffre (« seulement »).

\textbf{4.} \zh{他又生病了，这个月第二次了。}\\
« Il est encore tombé malade ; c'est déjà la deuxième fois ce mois-ci. »\\
\zh{又} + \zh{了} : répétition déjà réalisée, avec une nuance d'agacement.

\textbf{5.} \zh{这个菜真好吃，我想再吃一碗。}\\
« Ce plat est vraiment bon, je veux en manger encore un bol. »\\
\zh{再} : répétition projetée (le second bol n'est pas encore mangé).
\end{corrige}

\begin{corrige}[Exercice 6 — Texte à trous : 就, 才, 还, 又 ou 再 ?]

\textbf{1. \cle{就}} — \zh{电影七点就开始了，我们六点半就到了。}\\
\emph{Diànyǐng qī diǎn jiù kāishǐ le, wǒmen liù diǎn bàn jiù dào le.}\\
« Le film a commencé dès sept heures, nous sommes arrivés dès six heures et demie. »\\
Justification : le début du film est jugé tôt / ponctuel ; \zh{就} exprime « dès ».

\textbf{2. \cle{才}} — \zh{他昨天晚上十二点才回家。}\\
\emph{Tā zuótiān wǎnshang shí'èr diǎn cái huí jiā.}\\
« Hier soir, il n'est rentré qu'à minuit. »\\
Justification : minuit est jugé tard ; \zh{才} exprime « ne\ldots que ».

\textbf{3. \cle{还}} — \zh{我的弟弟还小，才五岁。}\\
\emph{Wǒ de dìdi hái xiǎo, cái wǔ suì.}\\
« Mon petit frère est encore petit, il n'a que cinq ans. »\\
Justification : \zh{还} exprime la continuation d'un état (« encore »).

\textbf{4. \cle{又}} — \zh{他又生病了，这个月第二次了。}\\
\emph{Tā yòu shēngbìng le, zhège yuè dì'èr cì le.}\\
« Il est encore tombé malade, la deuxième fois ce mois-ci. »\\
Justification : répétition déjà réalisée (passé, \zh{了}) ; \zh{又}.

\textbf{5. \cle{再}} — \zh{这个菜真好吃，我想再吃一碗。}\\
\emph{Zhège cài zhēn hǎochī, wǒ xiǎng zài chī yì wǎn.}\\
« Ce plat est vraiment bon, je veux en manger encore un bol. »\\
Justification : répétition projetée (futur) ; \zh{再}.

\textbf{6. \cle{才}} — \zh{作业太多了，我做了三个小时才做完。}\\
\emph{Zuòyè tài duō le, wǒ zuò le sān ge xiǎoshí cái zuòwán.}\\
« Il y avait trop de devoirs, j'ai mis trois heures à les finir. »\\
Justification : la durée est jugée longue ; schéma : verbe + \zh{了} + durée + \zh{才} + résultat.

\textbf{7. \cle{还}} — \zh{别走！我还有一个问题要问你。}\\
\emph{Bié zǒu ! Wǒ hái yǒu yí ge wèntí yào wèn nǐ.}\\
« Ne pars pas ! J'ai encore une question à te poser. »\\
Justification : \zh{还} + \zh{有} exprime l'ajout (« en plus »).

\textbf{8. \cle{再}} — \zh{今天的课听不懂，请老师再讲一遍。}\\
\emph{Jīntiān de kè tīng bu dǒng, qǐng lǎoshī zài jiǎng yí biàn.}\\
« Je n'ai pas compris la leçon d'aujourd'hui, que le professeur veuille bien l'expliquer encore une fois. »\\
Justification : on demande une répétition à venir ; \zh{再}, jamais \zh{又}.
\end{corrige}

\begin{corrige}[Exercice 7 — Remets les mots dans l'ordre]

\textbf{1.} \zh{他今天早上九点才起床。}\\
\emph{Tā jīntiān zǎoshang jiǔ diǎn cái qǐchuáng.}\\
« Ce matin, il ne s'est levé qu'à neuf heures. »\\
Schéma : Sujet + temps (\zh{今天早上}) + heure + \zh{才} + verbe.

\textbf{2.} \zh{我一个小时就做完了作业。} (ou \zh{我一个小时就把作业做完了。})\\
\emph{Wǒ yí ge xiǎoshí jiù zuòwán le zuòyè.}\\
« J'ai fini mes devoirs en une heure seulement. »\\
Schéma : Sujet + durée + \zh{就} + verbe-résultat + \zh{了} + complément.

\textbf{3.} \zh{他今天又迟到了。}\\
\emph{Tā jīntiān yòu chídào le.}\\
« Il est encore arrivé en retard aujourd'hui. »\\
Schéma : Sujet + temps + \zh{又} + verbe + \zh{了} (répétition réalisée).

\textbf{4.} \zh{弟弟还在睡觉呢。}\\
\emph{Dìdi hái zài shuìjiào ne.}\\
« Le petit frère dort encore. »\\
Schéma : Sujet + \zh{还在} + verbe + \zh{呢} (action en cours qui continue).

\textbf{5.} \zh{请你再说一遍。}\\
\emph{Qǐng nǐ zài shuō yí biàn.}\\
« Répète une fois, s'il te plaît. »\\
Schéma : \zh{请} + sujet + \zh{再} + verbe + complément de fréquence (\zh{一遍}).
\end{corrige}

\begin{corrige}[Exercice 8 — Transformation : exprime un jugement]

\textbf{1.} 他八点来。

a) tôt : \zh{他八点就来了。} \emph{Tā bā diǎn jiù lái le.}\\
« Il est arrivé dès huit heures. » (jugement positif : c'est tôt, c'est bien.)

b) tard : \zh{他八点才来。} \emph{Tā bā diǎn cái lái.}\\
« Il n'est arrivé qu'à huit heures. » (jugement négatif : c'est tard.)

\cle{Remarque} : avec \zh{就} on ajoute souvent \zh{了} (action achevée) ; avec \zh{才}, on ne met pas \zh{了} après le verbe dans ce schéma.

\textbf{2.} 我用了三十分钟做完作业。

a) facile : \zh{我用了三十分钟就做完了作业。} \emph{Wǒ yòng le sānshí fēnzhōng jiù zuòwán le zuòyè.}\\
« Je n'ai mis que trente minutes à finir mes devoirs. » (c'était facile, rapide.)

b) difficile : \zh{我用了三十分钟才做完作业。} \emph{Wǒ yòng le sānshí fēnzhōng cái zuòwán zuòyè.}\\
« J'ai mis trente minutes entières à finir mes devoirs. » (c'était long, difficile.)

\cle{Remarque} : la même durée (30 minutes) est jugée courte avec \zh{就} et longue avec \zh{才} : l'adverbe porte le jugement du locuteur.

\textbf{3.} 她十八岁大学毕业。

a) jeune et remarquable : \zh{她十八岁就大学毕业了。} \emph{Tā shíbā suì jiù dàxué bìyè le.}\\
« Elle a fini l'université dès dix-huit ans. » (admiration.)

b) tard : \zh{她十八岁才大学毕业。} \emph{Tā shíbā suì cái dàxué bìyè.}\\
« Elle n'a fini l'université qu'à dix-huit ans. » (jugement : plus tard que la norme attendue.)
\end{corrige}

\begin{corrige}[Exercice 9 — Type Bac : Compréhension de l'écrit et questions de langue]

\textbf{A. Questions de compréhension}

\textbf{1.} Le narrateur s'est levé à cinq heures et demie parce qu'il voulait encore réviser une fois le vocabulaire avant l'examen de chinois (\zh{我还想复习一遍生词}).

\textbf{2.} L'examen commence à huit heures (\zh{考试八点开始}) ; le narrateur arrive à l'école dès sept heures et demie (\zh{我七点半就到了学校}).

\textbf{3.} Non, l'examen était trop difficile : le texte dit \zh{考试太难了} et le narrateur précise qu'il a mis deux heures à le finir (\zh{我用了两个小时才做完}), ce qui montre la difficulté grâce à l'adverbe \zh{才} (jugement de longueur).

\textbf{4.} Le professeur annonce qu'il y aura encore un examen la semaine suivante (\zh{下个星期还有一次考试}) et encourage les élèves à faire encore des efforts (\zh{大家再努力吧！}).

\textbf{B. Questions de langue}

\textbf{1.} Dans \zh{我早上五点半就起床了}, \zh{就} exprime que l'action a été accomplie plus tôt que prévu : le locuteur juge que cinq heures et demie, c'est tôt. Règle : Sujet + (temps) + \zh{就} + verbe + \zh{了} = « dès\ldots ». Exemple personnel : \zh{我六点就到了教室。} \emph{Wǒ liù diǎn jiù dào le jiàoshì.} « Je suis arrivé en classe dès six heures. »

\textbf{2.} Dans \zh{我用了两个小时才做完}, on utilise \zh{才} parce que le locuteur juge la durée longue : l'examen était difficile, deux heures, c'est beaucoup. \zh{就} aurait signifié le contraire (« en seulement deux heures », jugement positif), ce qui contredirait \zh{考试太难了}. Schéma : Sujet + verbe + \zh{了} + durée + \zh{才} + résultat.

\textbf{3.} Deux emplois de \zh{还} :
\begin{itemize}
\item \zh{我还想复习一遍生词} : \zh{还} exprime l'ajout / la continuation d'une intention — « je voulais \emph{encore} réviser » (en plus du reste).
\item \zh{下个星期还有一次考试} : \zh{还} + \zh{有} exprime l'existence d'un élément supplémentaire — « il y a \emph{encore} un examen » (en plus de celui-ci).
\end{itemize}

\textbf{4.} La phrase \zh{又要考试了} (que le narrateur pourrait penser) illustre l'emploi de \zh{又} : l'examen est déjà annoncé, la situation « avoir un examen » se reproduit et est perçue comme une répétition subie (avec agacement, d'où \zh{唉}). \zh{再} exprimerait une répétition volontaire ou projetée (« refaire »), ce qui ne convient pas ici : le narrateur subit l'examen, il ne le décide pas. De plus, dans \zh{大家再努力吧}, \zh{再} est correct car le professeur invite à un effort futur.

\textbf{Barème indicatif (sur 20) :} A : 4 $\times$ 2 pts = 8 pts ; B : 3 pts par question (règle exacte 1,5 pt + exemple ou justification 1,5 pt) = 12 pts.
\end{corrige}

\begin{corrige}[Exercice 10 — Type Bac : Thème et version]

\textbf{A. Thème}

\textbf{1.} « Il n'est arrivé qu'à dix heures hier soir. »\\
\zh{他昨天晚上十点才到。}\\
\emph{Tā zuótiān wǎnshang shí diǎn cái dào.}\\
Points de vigilance : « ne\ldots que » (tard) = \zh{才} ; le temps \zh{昨天晚上} se place avant \zh{才} ; pas de \zh{了} après le verbe avec \zh{才}.

\textbf{2.} « Dès que je rentre, je fais mes devoirs. »\\
\zh{我一回家就做作业。}\\
\emph{Wǒ yì huí jiā jiù zuò zuòyè.}\\
Points de vigilance : « dès que\ldots » = schéma \zh{一\ldots 就\ldots} ; pas de \zh{了} car c'est une habitude.

\textbf{3.} « Il dort encore, ne le réveille pas. »\\
\zh{他还在睡觉，别叫醒他。}\\
\emph{Tā hái zài shuìjiào, bié jiàoxǐng tā.}\\
Points de vigilance : continuation d'une action en cours = \zh{还在} + verbe ; « ne\ldots pas » = \zh{别}.

\textbf{4.} « Il est encore en retard aujourd'hui ! »\\
\zh{他今天又迟到了！}\\
\emph{Tā jīntiān yòu chídào le !}\\
Points de vigilance : répétition déjà réalisée (le retard est constaté) = \zh{又} + \zh{了}, avec nuance d'agacement ; ne pas utiliser \zh{再}.

\textbf{5.} « Revenez demain, s'il vous plaît. »\\
\zh{请明天再来。}\\
\emph{Qǐng míngtiān zài lái.}\\
Points de vigilance : action remise au futur = \zh{再} ; ordre : \zh{请} + temps + \zh{再} + verbe.

\textbf{B. Version}

\textbf{1.} \zh{只有多练习，才能学好汉语。}\\
\emph{Zhǐyǒu duō liànxí, cáinéng xué hǎo Hànyǔ.}\\
« Ce n'est qu'en s'entraînant beaucoup qu'on peut bien apprendre le chinois. » (ou : « Il faut beaucoup s'exercer pour pouvoir bien apprendre le chinois. »)\\
Structure : \zh{只有\ldots 才\ldots} = « seulement si\ldots alors\ldots » (condition nécessaire).

\textbf{2.} \zh{她十五岁就会说三种语言了。}\\
\emph{Tā shíwǔ suì jiù huì shuō sān zhǒng yǔyán le.}\\
« À quinze ans, elle savait déjà parler trois langues. »\\
\zh{就} exprime l'admiration : c'est tôt, c'est remarquable.

\textbf{3.} \zh{他又忘了带书，老师很不高兴。}\\
\emph{Tā yòu wàng le dài shū, lǎoshī hěn bù gāoxìng.}\\
« Il a encore oublié d'apporter son livre, le professeur n'est pas content du tout. »\\
\zh{又} + \zh{了} : répétition constatée, nuance de reproche.

\textbf{Barème indicatif (sur 20) :} Thème 5 $\times$ 2 pts = 10 pts (choix de l'adverbe 1 pt, ordre des mots et caractères 1 pt) ; Version 3 $\times$ 3 pts + 1 pt de qualité globale = 10 pts.
\end{corrige}

\begin{corrige}[Exercice 11 — Type Bac : Correction de phrases et expression écrite guidée]

\textbf{A. Correction de phrases fautives}

\textbf{1.} Phrase fautive : \zh{他才来了昨天。}\\
Correction : \zh{他昨天才来。} \emph{Tā zuótiān cái lái.} « Il n'est arrivé qu'hier. »\\
Règle : le complément de temps se place avant l'adverbe et le verbe (Sujet + temps + \zh{才} + verbe) ; de plus, avec \zh{才} on ne met pas \zh{了} après le verbe.

\textbf{2.} Phrase fautive : \zh{请又说一遍。}\\
Correction : \zh{请再说一遍。} \emph{Qǐng zài shuō yí biàn.} « Répétez une fois, s'il vous plaît. »\\
Règle : pour une répétition demandée (à venir), on utilise \zh{再} ; \zh{又} ne s'emploie que pour une répétition déjà réalisée.

\textbf{3.} Phrase fautive : \zh{我还昨天去了市场。}\\
Correction : \zh{我昨天还去了市场。} \emph{Wǒ zuótiān hái qù le shìchǎng.} « Hier, je suis en plus allé au marché. »\\
Règle : le complément de temps \zh{昨天} précède l'adverbe \zh{还} (Sujet + temps + \zh{还} + verbe).

\textbf{4.} Phrase fautive : \zh{他明天又来看你。}\\
Correction : \zh{他明天再来看你。} \emph{Tā míngtiān zài lái kàn nǐ.} « Il reviendra te voir demain. »\\
Règle : avec un complément de temps futur (\zh{明天}), la répétition est projetée ; on utilise \zh{再}, jamais \zh{又}.

\textbf{B. Expression écrite guidée — proposition de rédaction modèle}

\zh{今天是汉语考试的日子。我早上五点就起床了，因为我还想复习一遍生词。我六点半就到了学校，可是考试九点才开始。考试不太难，我用了一个小时就做完了。考完以后，老师说我考得很好。明天我还要考数学，我要再努力！}\\
\emph{Jīntiān shì Hànyǔ kǎoshì de rìzi. Wǒ zǎoshang wǔ diǎn jiù qǐchuáng le, yīnwèi wǒ hái xiǎng fùxí yí biàn shēngcí. Wǒ liù diǎn bàn jiù dào le xuéxiào, kěshì kǎoshì jiǔ diǎn cái kāishǐ. Kǎoshì bú tài nán, wǒ yòng le yí ge xiǎoshí jiù zuòwán le. Kǎowán yǐhòu, lǎoshī shuō wǒ kǎo de hěn hǎo. Míngtiān wǒ hái yào kǎo shùxué, wǒ yào zài nǔlì !}\\
« Aujourd'hui, c'est le jour de l'examen de chinois. Je me suis levé dès cinq heures du matin, parce que je voulais encore réviser une fois le vocabulaire. Je suis arrivé à l'école dès six heures et demie, mais l'examen n'a commencé qu'à neuf heures. L'examen n'était pas trop difficile, je l'ai fini en une heure seulement. Après l'examen, le professeur a dit que j'avais bien réussi. Demain, j'ai encore l'épreuve de mathématiques, je dois encore faire des efforts ! »

Adverbes employés : \zh{就} (trois fois : lever tôt, arrivée tôt, devoirs finis vite), \zh{还} (deux fois : ajout, continuation), \zh{才} (début tardif de l'épreuve), \zh{再} (effort futur).

\cle{Points de vigilance pour les élèves :}
\begin{itemize}
\item Placer toujours le complément de temps (\zh{早上}, \zh{明天}\ldots) avant l'adverbe.
\item \zh{就} + verbe + \zh{了} pour le passé jugé tôt ; \zh{才} + verbe \emph{sans} \zh{了}.
\item Réserver \zh{又} au passé répété (avec \zh{了}) et \zh{再} au futur ou à la demande.
\item Ne pas oublier la particule \zh{了} pour les actions accomplies.
\end{itemize}

\textbf{Barème indicatif :} A : 4 $\times$ 1,5 pt = 6 pts (erreur identifiée 0,5 pt + phrase corrigée 0,5 pt + règle citée 0,5 pt) ; B : 10 pts selon le barème de l'énoncé (place des adverbes 4 pts, opposition \zh{又}/\zh{再} 2 pts, emploi de \zh{了} 2 pts, vocabulaire et caractères 2 pts).
\end{corrige}

\newpage

\coursec{Fiche bilan}

\begin{popfigure}
\begin{tikzpicture}[mindmap, grow cyclic, every node/.style={concept, align=center, font=\small}, root concept/.append style={concept color=popblue, font=\bfseries}, level 1/.append style={level distance=3.4cm, sibling angle=60}, text=white]
\node[root concept, align=center]{就 / 才\\还 / 又 / 再}
  child[concept color=poporange]{ node[align=center]{就\\« tôt, dès »\\avant l'heure prévue} }
  child[concept color=poppurple]{ node[align=center]{才\\« seulement, pas avant »\\après l'heure prévue} }
  child[concept color=popgreen]{ node[align=center]{还\\« encore, aussi »\\continuation} }
  child[concept color=poppink]{ node[align=center]{又\\« de nouveau »\\action passée répétée} }
  child[concept color=popgold]{ node[align=center]{再\\« encore une fois »\\action future répétée} }
  child[concept color=popdark]{ node[align=center]{Pièges\\ordre des mots\\又 $\neq$ 再} };
\end{tikzpicture}
\legende{Carte mentale de la leçon : les adverbes 就, 才, 还, 又, 再}
\end{popfigure}

\begin{bilan}
\textbf{1. Lexique clé à connaître par cœur}
\begin{itemize}
  \item \zh{就} \emph{jiù} : tôt, dès, seulement (avant le moment prévu) ; \zh{才} \emph{cái} : seulement, pas avant (après le moment prévu).
  \item \zh{还} \emph{hái} : encore, aussi, en plus ; \zh{又} \emph{yòu} : de nouveau (répétition déjà réalisée) ; \zh{再} \emph{zài} : encore une fois (répétition à venir).
  \item Mots associés : \zh{点} \emph{diǎn} (heure), \zh{起床} \emph{qǐchuáng} (se lever), \zh{回来} \emph{huílai} (revenir), \zh{一遍} \emph{yí biàn} (une fois), \zh{再说} \emph{zàishuō} (en reparler plus tard).
\end{itemize}

\textbf{2. Règles de grammaire à connaître par cœur}
\begin{center}
\begin{tabularx}{\linewidth}{|X|X|X|}
\hline
\rowcolor{popblueL} Structure & Exemple (caractères + pinyin) & Traduction \\
\hline
Heure / durée + \cle{就} + V & \zh{他六点就起床了。} \emph{Tā liù diǎn jiù qǐchuáng le.} & Il s'est levé dès six heures. \\
\hline
Heure / durée + \cle{才} + V & \zh{他九点才起床。} \emph{Tā jiǔ diǎn cái qǐchuáng.} & Il ne s'est levé qu'à neuf heures. \\
\hline
\cle{还} + V / Adj & \zh{我还要一个。} \emph{Wǒ hái yào yí ge.} & J'en veux encore un. \\
\hline
\cle{又} + V + 了 & \zh{他又迟到了。} \emph{Tā yòu chídào le.} & Il est encore arrivé en retard (fait accompli). \\
\hline
\cle{再} + V (futur) & \zh{请再说一遍。} \emph{Qǐng zài shuō yí biàn.} & Répétez encore une fois, s'il vous plaît. \\
\hline
\end{tabularx}
\end{center}

\textbf{3. Points critiques}
\begin{itemize}
  \item 就 et 才 se placent \textbf{avant le verbe}, après le complément de temps ou de quantité ; jamais en tête de phrase.
  \item Avec 就, on peut ajouter \zh{了} en fin de phrase ; avec 才, \textbf{jamais de} \zh{了} pour un fait passé.
  \item \zh{又} = répétition déjà accomplie (souvent avec \zh{了}) ; \zh{再} = répétition non encore réalisée (futur, souhait, impératif).
  \item \zh{还} exprime la continuation (« encore ») ou l'addition (« aussi, en plus ») ; ne pas le confondre avec \zh{也} (simple « aussi » entre sujets).
\end{itemize}

\textbf{4. Méthode express}
\begin{enumerate}
  \item Identifier le sens voulu : tôt / tard / continuation / répétition passée / répétition future.
  \item Choisir l'adverbe : 就 (tôt), 才 (tard), 还 (continue), 又 (refait), 再 (refera).
  \item Placer l'adverbe juste avant le verbe, après le sujet et le complément de temps.
  \item Vérifier la présence ou l'absence de \zh{了} final.
\end{enumerate}
\end{bilan}

\begin{aretenir}[Checklist avant le Bac]
\begin{itemize}
  \item $\square$ Je sais que 就 marque une action plus tôt que prévu et 才 une action plus tard que prévu.
  \item $\square$ Je place correctement 就 / 才 avant le verbe, jamais avant le sujet.
  \item $\square$ Je n'écris jamais \zh{了} final avec 才 pour un événement passé.
  \item $\square$ Je distingue 又 (répétition déjà faite, souvent avec \zh{了}) et 再 (répétition à venir).
  \item $\square$ Je sais que \zh{请再说一遍} utilise 再 car la répétition n'a pas encore eu lieu.
  \item $\square$ J'utilise 还 pour « encore » (continuation) et « en plus » (addition).
  \item $\square$ Je sais traduire « Il n'est rentré qu'à minuit » par \zh{他十二点才回来。}
  \item $\square$ Je sais traduire « Dès cinq heures, il était au stade » par \zh{他五点就到体育场了。}
  \item $\square$ Je vérifie toujours l'ordre : Sujet + temps + 就/才/还/又/再 + verbe.
  \item $\square$ Je prononce correctement les tons : jiù (4\up{e}), cái (2\up{e}), hái (2\up{e}), yòu (4\up{e}), zài (4\up{e}).
\end{itemize}
\end{aretenir}

\findecours

\end{document}
